% Worked solutions for Listas/lista-r5-asymmetric-information.md (15 exercises).
% This file writes exercises 2, 4, 8, 11, 12, 15; the other nine are already solved in
% lista-topic3-asymmetric-information-solutions.tex and are spliced in at the
% "%% INSERT-FROM" lines (one per exercise, in list order).
% Compile: pdflatex twice; check: pdffonts (Type 1 only, uni = yes).
% Every number is asserted in lista-r5-asymmetric-information-solutions_check.py.
\documentclass[11pt,a4paper]{article}
\usepackage[utf8]{inputenc}
\usepackage[T1]{fontenc}
\usepackage{lmodern}   % Type 1 vector fonts (replaces PK bitmaps)
\usepackage{cmap}      % ToUnicode CMap -> copyable text and accents
\usepackage[brazilian,english]{babel}
\usepackage{amsmath,amssymb,amsfonts}
\usepackage[a4paper,margin=2.2cm]{geometry}
\usepackage{booktabs}
\usepackage{enumitem}
\usepackage{fancyhdr}
\usepackage{xcolor}
\usepackage{microtype}

\pagestyle{fancy}\fancyhf{}
\lhead{\footnotesize Finan\c{c}as II --- Corporativas (MPE Insper)}
\rhead{\footnotesize List R5 --- Asymmetric information and capital structure}
\cfoot{\footnotesize\thepage}
\renewcommand{\headrulewidth}{0.3pt}
\setlength{\headheight}{14pt}
\setlist{topsep=2pt,itemsep=1pt}

\newcommand{\exsec}[1]{\section*{\color{blue!55!black}#1}\addcontentsline{toc}{section}{#1}}
\newcommand{\notation}{\par\noindent\textbf{Notation.}}
\newcommand{\step}[2]{\par\medskip\noindent\textbf{Step #1 --- #2.}~}
\newcommand{\intuition}{\par\medskip\noindent\textbf{Intuition.}~}
\newcommand{\bookans}{\par\medskip\noindent\textbf{Book answer.}~}
\newcommand{\refbib}{\par\smallskip\noindent\textbf{Ref:}~}
\newcommand{\assume}{\par\smallskip\noindent\textbf{Assumption.}~}
\newcommand{\pt}{\foreignlanguage{brazilian}}
% Section C (deeper solutions) macros
\newcommand{\asked}{\par\medskip\noindent\textbf{What is asked.}~}
\newcommand{\approach}{\par\smallskip\noindent\textbf{Approach.}~}
\newcommand{\answer}[1]{\par\smallskip\noindent\fbox{\textbf{Answer #1.}}~}
\newcommand{\pitfalls}{\par\medskip\noindent\textbf{Common mistakes.}}

\usepackage[hyphens]{url}
%% STATEMENT-PREAMBLE BEGIN (tools/statements_into_solutions.py)
\usepackage{longtable,array,calc}
\providecommand{\tightlist}{\setlength{\itemsep}{0pt}\setlength{\parskip}{0pt}}
\providecommand{\href}[2]{#2}
\providecommand{\real}[1]{#1}
\makeatletter\@ifundefined{c@none}{\newcounter{none}}{}\makeatother  % pandoc's \LTcaptype{none}
\newcommand{\stmtbegin}{\par\medskip\noindent\textbf{Statement.}\par\nopagebreak\begingroup\small\setlength{\tabcolsep}{3pt}\sloppy}
\newcommand{\stmtend}{\par\endgroup\nopagebreak\noindent\rule{\linewidth}{0.4pt}\par\medskip}
\DeclareUnicodeCharacter{2192}{\ensuremath{\rightarrow}}
\DeclareUnicodeCharacter{2212}{\ensuremath{-}}
\DeclareUnicodeCharacter{2022}{\textbullet{}}
\DeclareUnicodeCharacter{2248}{\ensuremath{\approx}}
\DeclareUnicodeCharacter{2032}{\ensuremath{^{\prime}}}
\DeclareUnicodeCharacter{2081}{\ensuremath{_{1}}}
\DeclareUnicodeCharacter{03B8}{\ensuremath{\theta}}
\DeclareUnicodeCharacter{03BB}{\ensuremath{\lambda}}
\DeclareUnicodeCharacter{03C8}{\ensuremath{\psi}}
\DeclareUnicodeCharacter{03B1}{\ensuremath{\alpha}}
\DeclareUnicodeCharacter{25A0}{\rule{0.7em}{0.7em}}
%% STATEMENT-PREAMBLE END
\begin{document}
\begin{center}{\Large\bfseries List R5 --- Asymmetric information and capital structure}\\[4pt]{\large Worked solutions}\end{center}

\noindent Statements are in \url{Listas/lista-r5-asymmetric-information.md}; each exercise is cited by number and printed page and its full text is reproduced right before its solution. Exercises 1, 3, 5, 6, 7, 9, 10, 13 and 14 are the same exercises as Topic~3 list exercises 5, 2, 3, 6, 11, 9, 10, 15 and 12, solved in \url{lista-topic3-asymmetric-information-solutions.tex}; their solutions are inserted at their positions below. Editions are the local copies: Tirole, \emph{The Theory of Corporate Finance} (2005); Amaro de Matos, \emph{Theoretical Foundations of Corporate Finance} (2001); Berk = Berk, DeMarzo \& Harford, \emph{Fundamentals of Corporate Finance}, 2nd ed. (2011); Grinblatt--Titman, \emph{Financial Markets and Corporate Strategy}, 2nd ed. (2002); Ross = Ross, Westerfield, Jordan \& Lamb, \emph{Fundamentos de Administra\c{c}\~ao Financeira}, 9th ed. (2013).

\tableofcontents

%% BEGIN-FROM lista-topic3-asymmetric-information-solutions.tex EX 5 AS 1
\exsec{1. Dividend suspension and the stock price --- Ross CQ 17.6, p.~614}
%% BEGIN-STATEMENT 1
\stmtbegin
\textbf{Source:} Ross, Cap. 17, Exercise CQ 17.6 (Revisão de conceitos e questões instigantes),
p.~614 (PDF 638) · Answer in book: no

Mês passado, a Companhia de Energia Elétrica de Angra dos Reis, que vinha tendo problemas com custos
acima dos previstos em uma usina nuclear em construção, anunciou que estava ``suspendendo
temporariamente os pagamentos devido à diminuição dos fluxos de caixa associados a esse programa de
investimentos''. O preço da ação da empresa caiu de \$28,50 para \$25 quando esse anúncio foi feito.
Como você interpretaria essa mudança no preço da ação (ou seja, em sua opinião, o que teria causado
isso)?
\stmtend
%% END-STATEMENT

\asked (a) The utility suspended its dividend, citing lower cash flows from a nuclear plant over budget, and its share price fell from \$28.50 to \$25. Interpret the fall: what caused it?

\approach Two benchmarks frame the answer. MM dividend irrelevance says the payout alone cannot change value in a perfect market, so the price drop must come from something else. The information content of dividends (dividends as a costly signal, lecture block 3.7) supplies it: managers smooth dividends and cut them only when they expect lower cash flows, so a suspension is bad news about future cash flows. Size the reaction first, then explain it.

\notation
\begin{itemize}
  \item $P_b=28.50$, $P_a=25$: share price before and after the announcement (\$); $r_{\text{ann}}$: announcement return.
\end{itemize}

\step{1}{Size of the reaction} The change in price and the announcement return:
\begin{align}
  P_a-P_b &= (25)-(28.50) = -3.50, \nonumber\\
  r_{\text{ann}} &= \frac{P_a-P_b}{P_b} = \frac{(-3.50)}{(28.50)} = -12.28\%. \label{t3e5-r}
\end{align}

\step{2}{What caused it} In a perfect market (MM dividend irrelevance and the homemade-dividend argument) changing the payout alone does not change value: the cash stays in the firm. The drop therefore reflects \emph{information}. Firms smooth dividends and cut them only with great reluctance, so a suspension tells the market that management no longer expects the current dividend to be sustainable. Here the firm says why: cost overruns on the nuclear plant have reduced cash flows (\pt{``diminui\c{c}\~ao dos fluxos de caixa''}). Investors revise expected future cash flows, and so expected future dividends, downward, and the price falls with their present value. The price fell because future dividends are expected to be lower, not because the payout ratio changed.

\step{3}{Secondary effects} (i) Clientele: investors who held the stock for current income may sell, adding temporary pressure. (ii) Keeping the cash avoids a costly external issue to finish the plant (pecking order), which by itself would support the price; the net fall of 12.28\% shows the bad news dominated.

\answer{(a)} The \$3.50 (12.28\%) fall is the market's reaction to the information in the suspension: lower expected future cash flows (hence lower future dividends) from the cost overruns, not the change in payout itself.

\intuition Dividends work as a signal because they are costly to maintain when cash flows are weak: only a firm that expects strong cash flows can keep paying. Cutting the dividend reveals the opposite, so the stock reprices on the news about cash flows, not on the dividend itself (lecture: payout policy as a signalling device).

\pitfalls
\begin{itemize}
  \item Saying the price fell ``because shareholders lost the dividend'': under MM the retained cash is still theirs; the payout change alone does not destroy value.
  \item Computing the return on the after-price, $3.50/25=14\%$, instead of on the before-price, $3.50/28.50=12.28\%$.
  \item Ignoring that the cut was partly anticipated (the overruns were known): the price reacts only to the \emph{new} information in the announcement.
\end{itemize}
\refbib{Ross Cap.~17, \S17.5, ``Conte\'udo informacional dos dividendos'', pp.~593--594; lecture Topic~3, ``Other signaling devices: payout policy''.}
%% END-FROM

% =============================================================================
\exsec{2. Raising debt after favourable information --- Grinblatt \& Titman 19.2, p.~687}
%% BEGIN-STATEMENT 2
\stmtbegin
\textbf{Source:} Grinblatt, Ch. 19, Exercise 19.2, p.~687 (PDF 700) · Answer in book: no

Why might a firm choose to increase its debt level in response to favorable information about its
future prospects?
\stmtend
%% END-STATEMENT

\asked Explain why a firm whose managers receive good private news about future cash flows might respond by \emph{raising} its debt, even though nothing in the MM benchmark ties leverage to news.

\approach Three mechanisms of the course each predict ``more debt after good news'', and the answer lists them: (1) \emph{signalling with debt} (Ross 1977; GT Results 19.8--19.9): a high fixed obligation is cheap for a firm that will meet it and expensive (distress, penalties for the manager) for one that will not, so only firms with good prospects can afford it and the market reads it as good news; (2) the \emph{pecking order} (Myers--Majluf; GT Result 19.10): managers who know the shares are undervalued avoid issuing them and fund with debt, or even use debt to buy the cheap shares back; (3) the \emph{trade-off} view: better expected earnings mean more taxable income to shield and a lower expected distress cost, so the value-maximizing debt level rises. The steps formalize (1), which is the chapter's main point, with the book's Example 19.6 numbers.

\notation
\begin{itemize}
  \item Two types of firm: good ($G$) and bad ($B$); annual earnings $X_G\in[50,60]$ and $X_B\in[25,30]$ (\$ million, GT Example 19.6, p.~672).
  \item $D\ge 0$: annual interest the firm commits to; default occurs when $X<D$.
  \item $V_G>V_B$: intrinsic (full-information) values of equity plus debt; $P_0$: market value today, which depends on what investors infer from $D$.
  \item $C>0$: financial distress cost borne if the firm defaults; $\omega\in(0,1)$: weight the manager puts on today's market value, $1-\omega$ on the long-run value net of distress costs (GT Result 19.1).
\end{itemize}

\assume To isolate the signalling channel, intrinsic value does not depend on $D$ unless default occurs (MM otherwise). For the numerical illustration only, each type's earnings are a perpetuity at the midpoint of its range discounted at $10\%$; this is our assumption, not the book's.

\step{1}{Manager's objective} The manager chooses $D$ to maximize a weighted sum of the current price and the long-run value:
\begin{align}
  U(D) &= \omega\,P_0(D) + (1-\omega)\bigl[V - C\cdot\mathbf 1\{X<D\}\bigr]. \label{r5e2-U}
\end{align}

\step{2}{A separating debt level} Investors believe ``$D>30$ means good type''. Any $D_G\in(30,50]$ is then safe for the good firm (its lowest earnings, 50, cover it) and certain default for the bad firm (its highest earnings, 30, do not). The book's choice $D_G=40$ lies in this window; the coverage ratios are
\begin{align}
  \frac{\min X_G}{D_G} &= \frac{(50)}{(40)} = 1.25 > 1, \nonumber\\
  \frac{\max X_B}{D_G} &= \frac{(30)}{(40)} = 0.75 < 1. \nonumber
\end{align}

\step{3}{Bad type does not mimic (credibility)} Under separation the bad type picks $D=0$ and is valued $P_0=V_B$; if it copies $D_G$ it is valued $V_G$ but defaults for sure. Using \eqref{r5e2-U}, it does not mimic iff
\begin{align}
  \omega V_B + (1-\omega)V_B &\ge \omega V_G + (1-\omega)(V_B - C) \nonumber\\
  V_B &\ge \omega V_G + V_B - \omega V_B - (1-\omega)C \nonumber\\
  (1-\omega)\,C &\ge \omega\,(V_G - V_B). \label{r5e2-cred}
\end{align}
This is GT Result 19.9: the signal is credible only if mimicking is costly enough for the bad type.

\step{4}{Good type prefers to signal} The good type never defaults at $D_G$, so signalling costs it nothing:
\begin{align}
  U_G(D_G) - U_G(0) &= \bigl[\omega V_G + (1-\omega)V_G\bigr] - \bigl[\omega V_B + (1-\omega)V_G\bigr] \nonumber\\
  &= \omega\,(V_G - V_B) > 0. \nonumber
\end{align}

\step{5}{Numbers} With midpoint earnings $55$ and $27.5$ and a $10\%$ discount rate,
\begin{align}
  V_G &= \frac{(55)}{(0.10)} = 550, \qquad V_B = \frac{(27.5)}{(0.10)} = 275, \nonumber
\end{align}
and with $\omega=0.2$, \eqref{r5e2-cred} requires
\begin{align}
  C &\ge \frac{\omega}{1-\omega}(V_G-V_B) = \frac{(0.2)}{(0.8)}\,(550-275) = (0.25)(275) = 68.75. \nonumber
\end{align}
If distress would cost the bad firm at least \$68.75 million, the debt increase cannot be faked and the price jumps from the pooled value to $V_G$ on announcement.

\step{6}{The other two channels} (i) Pecking order: if the good news is private, today's price $P_0<V_G$ undervalues the shares; issuing equity would transfer value to new shareholders, while debt (a fixed claim) is much less mispriced. A firm with news to fund prefers debt, and a firm with no project can borrow to repurchase its cheap shares (the CUC International recapitalization, GT p.~672). (ii) Taxes and distress: higher expected earnings raise the value of interest tax shields and lower the expected distress cost of any given $D$, so even the full-information optimal leverage rises (GT Result 19.8: ``the firm will be generating taxable earnings in the future'').

\answer{\unskip} Because debt is a credible signal (a fixed obligation only a firm with good prospects can carry without facing distress, condition \eqref{r5e2-cred}); because undervalued equity is a costly source of funds, so a firm with good private news funds with debt or borrows to repurchase shares; and because better prospects raise the tax value and lower the distress cost of debt.

\intuition Good news that only managers know is worth nothing until investors believe it; talk is cheap, but a promise to pay \$40 million a year is not, if the firm could not keep it. Debt does the talking because its cost falls on firms whose cash flows would not cover it. The same news makes equity look cheap to insiders, which reinforces the move towards debt.

\pitfalls
\begin{itemize}
  \item Saying ``debt is a good signal'' without the credibility condition: if the bad type's cost of mimicking ($C$, or the manager's penalty) is small, every firm borrows and the signal carries no information.
  \item Thinking the signal destroys value for the good firm: at $D_G=40$ it never defaults, so the signal is free for it (GT Example 19.6); the cost is borne only off the equilibrium path.
  \item Confusing the two information stories: signalling explains the positive price reaction to leverage increases; the pecking order explains why the undervalued firm avoids equity. Both point the same way here.
\end{itemize}

\refbib{Grinblatt \& Titman Ch.~19, \S19.5 ``The Information Content of the Debt-Equity Choice'', pp.~671--678 (Results 19.8--19.9 and Example 19.6, p.~672; Result 19.10, p.~677); Amaro de Matos \S3.2.1 pp.~79--81 (Ross 1977); lecture Topic~3, ``Signalling''. No book answer.}

%% BEGIN-FROM lista-topic3-asymmetric-information-solutions.tex EX 2 AS 3
\exsec{3. Investment and available cash --- Grinblatt \& Titman 19.4, p.~687}
%% BEGIN-STATEMENT 3
\stmtbegin
\textbf{Source:} Grinblatt, Ch. 19, Exercise 19.4, p.~687 (PDF 700) · Answer in book: no

Classical finance theory suggests that firms take projects with positive NPVs regardless of the
amount of cash the firm has available. However, empirical evidence suggests that the amount that
firms invest is heavily dependent on their available cash flows. Why might this be?
\stmtend
%% END-STATEMENT

\asked (a) The NPV rule says a firm should take every positive-NPV project whatever its cash, yet in the data investment moves with cash flow. Explain why: name the friction that makes internal cash cheaper than outside money, and show how it can make investment depend on cash.

\approach The friction is adverse selection in the issue of new securities (the lemons/pooling model of lecture block 3.2 and the Myers--Majluf pecking order of block 3.5). Outside investors cannot tell a good firm from a bad one, so they price any new claim at the pooled value; a good firm therefore sells its claims at a discount. That discount is proportional to the amount raised outside, so the more of the investment the firm pays for with its own cash, the cheaper its funding. A short model with the lecture's numbers makes the point quantitative; other frictions (moral hazard, issue costs) give the same prediction and are listed at the end.

\notation
\begin{itemize}
  \item Lecture model (Topic~3): project costs $I$, pays $R$ on success and 0 on failure; good type succeeds with probability $p$, bad type with $q<p$; $\alpha\in(0,1)$: share of good types; $m\equiv\alpha p+(1-\alpha)q$: investors' prior success probability.
  \item $A\in[0,I)$: internal cash (retained earnings) of the firm, the same for both types (the bad type copies the good type); $R_u$: repayment to outside investors on success.
  \item $U_G(A)$: good type's net payoff (gross payoff minus the $A$ it puts in).
\end{itemize}

\assume Both types hold the same cash $A$ and pool on the same external claim (as in the lecture's pecking-order slide, the bad type ``follows suit''); outside investors break even.

\step{1}{Why ``invest in every positive-NPV project'' fails} The classical rule assumes external funds are priced fairly. When managers know more than investors, new securities sell at the price of the average firm, so a good firm sells its claims too cheaply (Myers--Majluf, GT Result~19.10). Internal cash carries no such discount, so the cost of funds depends on how much of the investment the firm can pay for itself.

\step{2}{Good type's payoff with internal cash $A$} The firm borrows $I-A$ and promises $R_u$ on success. Because both types issue the same claim, investors cannot tell them apart and expect success with the pooled probability $m$; they break even when
\begin{align}
  m R_u &= I-A \label{t3e2-be}\\
  R_u &= \frac{I-A}{m}. \nonumber
\end{align}
The good type keeps $R-R_u$ on success (probability $p$) and has put in $A$. Her net payoff, substituting \eqref{t3e2-be} and then adding and subtracting $I$:
\begin{align}
  U_G(A) &= p(R-R_u) - A \nonumber\\
         &= pR - p\left(\frac{I-A}{m}\right) - A \nonumber\\
         &= pR - I + (I-A) - \frac{p}{m}(I-A) \nonumber\\
         &= (pR-I) - (I-A)\left(\frac{p}{m}-1\right). \label{t3e2-UG}
\end{align}
The first term is the NPV; the second, the adverse-selection discount, is proportional to the amount raised outside, $I-A$, and vanishes at $A=I$. It is positive because $m<p$: the good type pays for the bad type's higher failure rate.

\step{3}{Numbers from the lecture example} Take $R=50$, $I=30$, $p=0.8$, $q=0.7$, $\alpha=0.5$. The pooled probability and the NPV are
\begin{align}
  m &= (0.5)(0.8)+(0.5)(0.7) = 0.40+0.35 = 0.75, \nonumber\\
  pR-I &= (0.8)(50)-(30) = 40-30 = 10, \nonumber\\
  \frac{p}{m}-1 &= \frac{(0.8)}{(0.75)}-1 = 1.066667-1 = 0.066667. \nonumber
\end{align}
Substituting into \eqref{t3e2-UG}:
\begin{align}
  U_G(0) &= (10)-(30)(0.066667) = 10-2 = 8, \nonumber\\
  U_G(15) &= (10)-(15)(0.066667) = 10-1 = 9. \nonumber
\end{align}
Half the investment paid internally halves the discount.

\step{4}{Investment that depends on cash} Take the lecture's breakdown case $q=0.5$, $\alpha=0.2$:
\begin{align}
  m &= (0.2)(0.8)+(0.8)(0.5) = 0.16+0.40 = 0.56, \nonumber\\
  mR &= (0.56)(50) = 28 < 30 = I. \nonumber
\end{align}
With $A=0$ even handing investors the whole income ($R_u=R$) returns only 28 for a loan of 30, so nobody lends. Limited liability caps the repayment at $R_u\le R$, so by \eqref{t3e2-be} lending requires
\begin{align}
  \frac{I-A}{m} &\le R \nonumber\\
  I-A &\le mR \nonumber\\
  A &\ge I-mR = (30)-(28) = 2. \label{t3e2-Amin}
\end{align}
Two firms with the same positive-NPV project ($pR-I=10$) behave differently: one with $A\ge 2$ invests, one with $A<2$ cannot. Investment is a function of cash.

\step{5}{Other frictions with the same prediction} (i) Moral hazard: pledgeable income is limited, so borrowing capacity is a multiple of net worth (Tirole \S3.4, $I\le kA$). (ii) Issue (flotation) costs make external money dearer than internal money. (iii) Agency in the other direction: managers with free cash flow overinvest. All three make investment rise with internal cash; GT emphasise the information channel.

\answer{(a)} Because outside finance is dearer than internal cash. Under asymmetric information new claims are priced at the pooled value, so a good firm pays a discount $(I-A)(p/m-1)$ that grows with the amount raised outside (Step~2: 2 with no cash, 1 with half the investment in cash); in bad pools it cannot raise funds at all unless $A\ge I-mR$ (Step~4: $A\ge2$). Moral hazard and issue costs reinforce the same link.

\intuition External finance is costly because outsiders cannot tell good firms from bad ones and price every new claim at the pooled value. Internal cash avoids the discount, so firms invest out of cash first (the pecking order), and investment tracks cash flow even when the NPV rule says it should not.

\pitfalls
\begin{itemize}
  \item Answering ``because cash is cheaper'' without saying \emph{why}: in a perfect market (MM) a dollar of cash and a dollar raised outside cost the same.
  \item Claiming the good firm rejects a positive-NPV project because the project is bad: the project's NPV is unchanged; what changes is the transfer to new investors.
  \item Pricing the claim at the good type's probability $p$: investors only know the pool, so the break-even uses $m$.
\end{itemize}
\refbib{Grinblatt \& Titman Ch.~19, \S19.5 ``The Information Content of the Debt-Equity Choice'', pp.~671--678 (Example~19.8 and Result~19.10, pp.~676--677); Tirole \S6.2.2 pp.~244--249 and \S3.4 pp.~127--129; lecture Topic~3, ``Pecking-order hypothesis''.}
%% END-FROM

% =============================================================================
\exsec{4. Signaling with payout policy (AMC) --- Berk 17.14--17.18, pp.~528--529}
%% BEGIN-STATEMENT 4
\stmtbegin
\textbf{Source:} Berk, Ch. 17, Exercises 17.14--17.18 (shared setup), pp.~528--529 (PDF 559) ·
Answer in book: no

Use the following information to answer Problems 14 through 18:

AMC Corporation currently has an enterprise value of \$400 million and \$100 million in excess cash.
The firm has 10 million shares outstanding and no debt. Suppose AMC uses its excess cash to
repurchase shares. After the share repurchase, news will come out that will change AMC's enterprise
value to either \$600 million or \$200 million.

\begin{enumerate}
\def\labelenumi{\arabic{enumi}.}
\setcounter{enumi}{13}
\item
  What is AMC's share price prior to the share repurchase?
\item
  What would AMC's share price be after the repurchase if its enterprise value goes up? What would
  AMC's share price be after the repurchase if its enterprise value declines?
\end{enumerate}

*16. Suppose AMC waits until after the news comes out to do the share repurchase. What would AMC's
share price be after the repurchase if its enterprise value goes up? What would AMC's share price be
after the repurchase if its enterprise value declines?

\begin{enumerate}
\def\labelenumi{\arabic{enumi}.}
\setcounter{enumi}{16}
\tightlist
\item
  Suppose AMC management expects good news to come out. Based on your answers to Problems 15 and 16,
  if management wants to maximize AMC's ultimate share price, will they undertake the repurchase
  before or after the news comes out? When would management undertake the repurchase if they expect
  bad news to come out?
\end{enumerate}

*18. Given your answer to Problem 17, what effect would you expect an announcement of a share
repurchase to have on the stock price? Why?
\stmtend
%% END-STATEMENT

\asked (14) The share price before any repurchase. (15) The price after a repurchase done \emph{before} the news, in the good-news case (enterprise value \$600 million) and the bad-news case (\$200 million). (16) The same two prices if the repurchase is done \emph{after} the news. (17) The timing a value-maximizing management picks when it expects good news, and when it expects bad news. (18) What an announced repurchase should do to the stock price, and why.

\approach Pure accounting of market values with no taxes: a repurchase at a fair price does not change the price per share (MM payout irrelevance). When the price is \emph{not} fair because managers know the news, buying back at the old price transfers value between departing and remaining shareholders. Comparing (15) with (16) shows the transfer; (17)--(18) are the signalling consequence (Berk \S17.5, ``Signaling and Share Repurchases'').

\notation
\begin{itemize}
  \item $EV_0=400$: enterprise value now; $EV_H=600$, $EV_L=200$: enterprise value after good or bad news (\$ million).
  \item $C=100$: excess cash; $N_0=10$: shares outstanding (million); no debt.
  \item $P_0$: price before the repurchase; $n$: shares repurchased; $N_1=N_0-n$: shares left.
  \item $P^{\text{bef}}_H,P^{\text{bef}}_L$: final price when buying back before the news; $P^{\text{aft}}_H,P^{\text{aft}}_L$: when buying back after it.
\end{itemize}

\step{1}{(14) Price before the repurchase} Equity value is enterprise value plus excess cash (no debt):
\begin{align}
  P_0 &= \frac{EV_0 + C}{N_0} = \frac{(400)+(100)}{(10)} = 50. \label{r5e4-P0}
\end{align}
\answer{(14)} \$50 per share.

\step{2}{(15) Repurchase before the news} The \$100 million buys shares at $P_0$:
\begin{align}
  n &= \frac{C}{P_0} = \frac{(100)}{(50)} = 2, \qquad N_1 = N_0 - n = (10)-(2) = 8 \text{ million}. \nonumber
\end{align}
The cash is gone, so equity equals the new enterprise value spread over 8 million shares:
\begin{align}
  P^{\text{bef}}_H &= \frac{EV_H}{N_1} = \frac{(600)}{(8)} = 75, \label{r5e4-bH}\\
  P^{\text{bef}}_L &= \frac{EV_L}{N_1} = \frac{(200)}{(8)} = 25. \label{r5e4-bL}
\end{align}
\answer{(15)} \$75 if enterprise value rises; \$25 if it falls.

\step{3}{(16) Repurchase after the news} The news is in the price first, so the buyback is at the new fair price. Good news:
\begin{align}
  P^{\text{news}}_H &= \frac{EV_H + C}{N_0} = \frac{(600)+(100)}{(10)} = 70, \nonumber\\
  n_H &= \frac{C}{P^{\text{news}}_H} = \frac{(100)}{(70)} = 1.4286, \qquad N_{1H} = (10)-(1.4286) = 8.5714, \nonumber\\
  P^{\text{aft}}_H &= \frac{EV_H}{N_{1H}} = \frac{(600)}{(8.5714)} = 70. \label{r5e4-aH}
\end{align}
Bad news:
\begin{align}
  P^{\text{news}}_L &= \frac{EV_L + C}{N_0} = \frac{(200)+(100)}{(10)} = 30, \nonumber\\
  n_L &= \frac{C}{P^{\text{news}}_L} = \frac{(100)}{(30)} = 3.3333, \qquad N_{1L} = (10)-(3.3333) = 6.6667, \nonumber\\
  P^{\text{aft}}_L &= \frac{EV_L}{N_{1L}} = \frac{(200)}{(6.6667)} = 30. \label{r5e4-aL}
\end{align}
A repurchase at a fair price leaves the price per share unchanged, as MM predicts.
\answer{(16)} \$70 if enterprise value rises; \$30 if it falls.

\step{4}{(17) Timing} Compare \eqref{r5e4-bH} with \eqref{r5e4-aH} and \eqref{r5e4-bL} with \eqref{r5e4-aL}:
\begin{align}
  \text{good news: } P^{\text{bef}}_H - P^{\text{aft}}_H &= (75)-(70) = 5 > 0, \nonumber\\
  \text{bad news: } P^{\text{aft}}_L - P^{\text{bef}}_L &= (30)-(25) = 5 > 0. \nonumber
\end{align}
Buying at \$50 when the shares are worth \$70 is a positive-NPV deal for remaining shareholders (they gain what the sellers lose); buying at \$50 when they are worth \$30 is a negative-NPV one.
\answer{(17)} Expecting good news: repurchase before it comes out (\$75 > \$70). Expecting bad news: wait and repurchase after (\$30 > \$25).

\step{5}{(18) Announcement effect} Investors know the rule in (17): management repurchases early only when it expects good news. An announced repurchase therefore reveals that insiders think the shares are undervalued, and the price rises on the announcement (part of the \$5 transfer is competed away as the price moves towards the good-news value).
\answer{(18)} A positive price reaction: the repurchase signals that managers believe the stock is undervalued.

\intuition With symmetric information the buyback is a zero-NPV swap of cash for shares. With private information it is a trade between the firm and its own shareholders at a price the managers know is wrong, and they only trade when it favours those who stay. Because of that, the decision itself carries the information (Berk point 3: ``repurchasing shares when managers perceive the stock to be under-valued is a positive-NPV investment'').

\pitfalls
\begin{itemize}
  \item Forgetting to remove the \$100 million of cash from equity after the buyback (dividing $EV+C$ by 8 million shares gives \$87.50 and \$37.50, which are wrong).
  \item Using \$50 as the repurchase price in (16): after the news the shares trade at \$70 or \$30.
  \item Reading (18) as ``buybacks create value'': the total value of the firm is the same; the gain to remaining holders is a transfer from those who sold.
\end{itemize}

\refbib{Berk, DeMarzo \& Harford Ch.~17, \S17.5 ``Signaling with Payout Policy'', subsection ``Signaling and Share Repurchases'', pp.~519--520; lecture Topic~3, ``Payout as a signal''. No book answer.}

%% BEGIN-FROM lista-topic3-asymmetric-information-solutions.tex EX 3 AS 5
\exsec{5. Leverage, agency costs and information (IST) --- Berk 16.29, p.~496}
%% BEGIN-STATEMENT 5
\stmtbegin
\textbf{Source:} Berk, Ch. 16, Exercise 16.29 (starred), p.~496 (PDF 527) · Answer in book: no

Info Systems Technology (IST) manufactures microprocessor chips for use in appliances and other
applications. IST has no debt and 100 million shares outstanding. The correct price for these shares
is either \$14.50 or \$12.50 per share. Investors view both possibilities as equally likely, so the
shares currently trade for \$13.50.

IST must raise \$500 million to build a new production facility. Because the firm would suffer a
large loss of both customers and engineering talent in the event of financial distress, managers
believe that if IST borrows the \$500 million, the present value of financial distress costs will
exceed any tax benefits by \$20 million. At the same time, because investors believe that managers
know the correct share price, IST faces a lemons problem if it attempts to raise the \$500 million
by issuing equity.

\begin{enumerate}
\def\labelenumi{\alph{enumi}.}
\item
  Suppose that if IST issues equity, the share price will remain \$13.50. To maximize the long-term
  share price of the firm once its true value is known, would managers choose to issue equity or
  borrow the \$500 million if:

  \begin{enumerate}
  \def\labelenumii{\roman{enumii}.}
  \tightlist
  \item
    They know the correct value of the shares is \$12.50?
  \item
    They know the correct value of the shares is \$14.50?
  \end{enumerate}
\item
  Given your answer to part (a), what should investors conclude if IST issues equity? What will
  happen to the share price?
\item
  Given your answer to part (a), what should investors conclude if IST issues debt? What will happen
  to the share price in that case?
\item
  How would your answers change if there were no distress costs, but only the tax benefits of
  leverage?
\end{enumerate}
\stmtend
%% END-STATEMENT

\asked (a) Taking the equity issue price as fixed at \$13.50, compare the long-run share price (once the truth is out) under an equity issue and under a \$500 million loan, for a manager who knows the shares are worth \$12.50 (i) and \$14.50 (ii), and say which she picks. (b) What investors should infer from an equity issue and where the price goes. (c) The same for a debt issue. (d) Redo the reasoning when debt brings only tax benefits and no distress costs.

\approach This is the lemons problem in an equity issue (Berk \S16.6, lecture blocks 3.2 and 3.5). Equity is the claim whose value depends most on what the manager knows, so selling it at a pooled price is a gain for an overvalued firm and a loss for an undervalued one. Debt avoids that transfer but here costs \$20 million of net distress costs. Part (a) compares the two payoffs type by type; parts (b)--(c) use the fact that rational investors read the choice of security as a signal (a separating outcome); part (d) removes the cost that made the separation possible.

\notation
\begin{itemize}
  \item $N_0=100$ million shares; $v\in\{12.50,\,14.50\}$: true price per share before the financing (equally likely); $P_m=13.50$: market price.
  \item $K=500$ million: amount to raise; $C=20$ million: PV of distress costs net of tax benefits if IST borrows $K$.
  \item $n$: new shares issued at $P_m$; $P^E(v)$, $P^D(v)$: long-run share price (after the truth is known) under equity and debt financing; $X>0$: net tax benefit of debt in part (d).
\end{itemize}

\assume The true prices $v$ already include the facility's NPV, so the \$500 million raised is invested in assets worth what is paid for them. Under debt, the \$500 million debt claim offsets the \$500 million of cash raised, and equity loses $C$.

\step{1}{(a) Price of each policy once the truth is known} Equity: IST sells $n$ shares at $P_m$; once the truth is known the firm's equity is worth the old value $N_0v$ plus the cash raised $K$ (invested at zero NPV beyond what $v$ already includes), spread over $N_0+n$ shares:
\begin{align}
  n &= \frac{K}{P_m} = \frac{(500)}{(13.50)} = 37.037 \text{ million}, \label{t3e3-n}\\
  P^E(v) &= \frac{N_0 v + K}{N_0+n}. \label{t3e3-PE}
\end{align}
Debt: the \$500 million of assets bought is offset by the \$500 million debt claim, so equity keeps all the old shares and bears only the net distress cost $C$:
\begin{align}
  P^D(v) &= \frac{N_0 v + K - K - C}{N_0} = \frac{N_0 v - C}{N_0}. \label{t3e3-PD}
\end{align}

\textit{(i) $v=12.50$.} Substituting into \eqref{t3e3-PE} and \eqref{t3e3-PD}:
\begin{align}
  P^E(12.50) &= \frac{(100)(12.50)+(500)}{(100)+(37.037)} = \frac{1{,}750}{137.037} = 12.77, \nonumber\\
  P^D(12.50) &= \frac{(100)(12.50)-(20)}{(100)} = 12.30. \nonumber
\end{align}
Since $12.77>12.30$, managers \textbf{issue equity}: selling overvalued shares transfers $N_0[P^E(12.50)-12.50]=(100)(12.7703-12.50)=27.03$ million from new to old shareholders, while debt would cost them $C=20$ million.

\textit{(ii) $v=14.50$.}
\begin{align}
  P^E(14.50) &= \frac{(100)(14.50)+(500)}{(100)+(37.037)} = \frac{1{,}950}{137.037} = 14.23, \nonumber\\
  P^D(14.50) &= \frac{(100)(14.50)-(20)}{(100)} = \frac{1{,}430}{100} = 14.30. \nonumber
\end{align}
Since $14.30>14.23$, managers \textbf{borrow}: the net distress cost, $C/N_0=20/100=0.20$ per share, is smaller than the loss from selling undervalued shares, $14.50-14.23=0.27$ per share.

\answer{(a)} (i) Value \$12.50: issue equity (\$12.77 against \$12.30 with debt). (ii) Value \$14.50: borrow (\$14.30 against \$14.23 with equity).

\step{2}{(b) An equity issue} By Step~1 only the low-value firm issues equity, so investors who know the managers' incentives infer $v=12.50$ and the price falls from 13.50 to \textbf{12.50}. Check that this price is consistent: at 12.50 IST would sell
\begin{align}
  \frac{K}{12.50} = \frac{(500)}{(12.50)} = 40 \text{ million shares}, \qquad
  \frac{N_0(12.50)+K}{N_0+40} = \frac{(1{,}250)+(500)}{(100)+(40)} = \frac{1{,}750}{140} = 12.50, \nonumber
\end{align}
so the new shareholders pay exactly what their shares are worth.

\answer{(b)} An equity issue reveals that the shares are worth \$12.50; the price falls from \$13.50 to \$12.50.

\step{3}{(c) A debt issue} Only the high-value firm borrows, so investors infer $v=14.50$ and price in the distress cost: the price rises to $P^D(14.50)=$ \textbf{14.30}.

\answer{(c)} A debt issue reveals that the shares are worth \$14.50; net of the \$0.20 of distress costs per share the price rises to \$14.30.

\step{4}{(d) Tax benefits only} Now debt adds a net tax benefit $X>0$ instead of costing $C$, so \eqref{t3e3-PD} becomes
\begin{align}
  P^D(v) = \frac{N_0v+X}{N_0} = v+\frac{X}{N_0} = v+\frac{X}{100}. \nonumber
\end{align}
\textit{Part (a) redone.} The high-value firm borrows, since $14.50+X/100>14.50>14.23=P^E(14.50)$. With the issue price fixed at 13.50, the low-value firm issues equity only if
\begin{align}
  P^E(12.50) &> 12.50 + \frac{X}{100} \nonumber\\
  12.7703 - 12.50 &> \frac{X}{100} \nonumber\\
  X &< 100\,(0.2703) = 27.03 \text{ million}. \label{t3e3-X}
\end{align}
\textit{Parts (b)--(c) redone.} By the logic of (b), an equity issue would reveal the low type and be priced at 12.50, which leaves the low type with $12.50<12.50+X/100$ (the value under debt). So in equilibrium \textbf{both types borrow}: no equity is issued, a debt issue carries no information, and the price moves only by the tax benefit, to
\begin{align}
  \tfrac12(12.50)+\tfrac12(14.50)+\frac{X}{100} = 13.50+\frac{X}{100}. \nonumber
\end{align}
The answers to (b)--(c) change: equity would be read as bad news but is never chosen, and debt is no longer a signal of a high value.

\answer{(d)} Debt now dominates for both types. (Against a fixed \$13.50 equity price the low type would still issue equity if $X<\$27.03$ million, but investors price any equity issue at \$12.50, so it borrows too.) Both types issue debt, the issue reveals nothing, and the price rises only by the tax benefit, to $13.50+X/100$ per share: the pecking order.

\intuition Equity is the security whose value is most sensitive to what managers know, so issuing it says ``our shares are overpriced''. With distress costs, only the overvalued firm accepts that stigma, and the choice of security separates the types. Without distress costs, debt dominates for both, and the market learns nothing from the choice: this is the pecking order.

\pitfalls
\begin{itemize}
  \item Forgetting the cash raised in the equity case: the post-issue value is $N_0v+K$, not $N_0v$, divided by $N_0+n$ shares.
  \item Subtracting the \$500 million debt from equity but not adding the \$500 million of new assets (or vice versa) in the debt case; the two cancel and only $C$ is left.
  \item In (b)--(c), keeping the price at \$13.50: once the security choice reveals the type, the price must move to the revealed value.
\end{itemize}
\refbib{Berk Ch.~16, \S16.6 ``Additional Consequences of Leverage: Agency Costs and Information'', subsection ``Debt and Information'', pp.~485--487; Ch.~14 lemons principle, p.~431.}
%% END-FROM

%% BEGIN-FROM lista-topic3-asymmetric-information-solutions.tex EX 6 AS 6
\exsec{6. Privately known private benefit and market breakdown --- Tirole 6.1, p.~273}
%% BEGIN-STATEMENT 6
\stmtbegin
\textbf{Source:} Tirole, Ch. 6, Exercise 6.1, p.~273 (PDF 286) · Answer in book: yes, p.~595

Section 6.2 illustrated the possibility of market breakdown without the possibility of signaling.
This exercise supplies another illustration. Let us consider the fixed-investment model of Section
3.2 and assume that only the borrower knows the private benefit associated with misbehavior. When
the borrower has private information about this parameter, lenders are concerned that this private
benefit might be high and induce the borrower to misbehave. In the parlance of information
economics, the ``bad types'' are the types of borrower with high private benefit. We study the case
of two possible levels of private benefit (see Exercise 6.2 for an extension to a continuum of
possible types). The borrower wants to finance a fixed-size project costing \(I\), and, for
simplicity, has no equity (\(A = 0\)). The project yields \(R\) (success) or 0 (failure). The
probability of success is \(p_H\) or \(p_L\), depending on whether the borrower works or shirks,
with \(\Delta p \equiv p_H - p_L > 0\). There is no private benefit when working. The private
benefit \(B\) enjoyed by the borrower when shirking is either \(B_L > 0\) or \(B_H > B_L\). The
borrower will be labeled a ``good borrower'' when \(B = B_L\) and a ``bad borrower'' when
\(B = B_H\). At the date of contracting, the borrower knows the level of her private benefit, while
the capital market puts (common knowledge) probabilities \(\alpha\) that the borrower is a good
borrower and \(1 - \alpha\) that she is a bad borrower. All other parameters are common knowledge
between the borrower and the lenders.

To make things interesting, let us assume that under asymmetric information, the lenders are
uncertain about whether the project should be funded:

\[p_H\left(R - \frac{B_H}{\Delta p}\right) < I < p_H\left(R - \frac{B_L}{\Delta p}\right). \qquad (1)\]

Assume that investors cannot break even if the borrower shirks:

\[p_L R < I. \qquad (2)\]

\begin{enumerate}
\def\labelenumi{(\roman{enumi})}
\item
  Note that the investor cannot finance only good borrowers. Assume that the entrepreneur receives
  no reward in the case of failure (this is indeed optimal); consider the effect of rewards \(R_b\)
  in the case of success that are (a) smaller than \(B_L/\Delta p\), (b) larger than
  \(B_H/\Delta p\), (c) between these two values.
\item
  Show that there exists \(\alpha^*\), \(0 < \alpha^* < 1\), such that
\end{enumerate}

\begin{itemize}
\tightlist
\item
  no financing occurs if \(\alpha < \alpha^*\),
\item
  financing is an equilibrium if \(\alpha \geqslant \alpha^*\).
\end{itemize}

\begin{enumerate}
\def\labelenumi{(\roman{enumi})}
\setcounter{enumi}{2}
\tightlist
\item
  Describe the ``cross-subsidies'' between types that occur when borrowing is feasible.
\end{enumerate}

\paragraph{Context needed to solve it}

\emph{From Tirole §3.2.1--3.2.2 (the fixed-investment model), pp.~115--117 (PDF 128--130):}

\textbf{3.2.1 The Fixed-Investment Model.} Variants of the following entrepreneurial model will be
used in the following: an entrepreneur (also called the ``insider'' or the ``borrower,'' ``she'')
has a project. This project requires a fixed investment \(I\). The entrepreneur initially has
``assets'' (``cash on hand'' or ``net worth'') \(A < I\). For the moment we interpret these assets
as being cash or liquid securities that can be used toward covering the cost of investment.
{[}\ldots{]} The entrepreneur's cash can either be invested in the project or used for consumption.
To implement the project the entrepreneur must borrow \(I - A\) from lenders. {[}\ldots{]}

\emph{Project.} If undertaken, the project either succeeds, that is, yields verifiable income
\(R > 0\), or fails, that is, yields no income. The probability of success is denoted by \(p\). The
project is subject to moral hazard. The entrepreneur can ``behave'' (``work,'' ``exert effort,''
``take no private benefit'') or ``misbehave'' (``shirk,'' ``take a private benefit''); or,
equivalently, the entrepreneur chooses between a project with a high probability of success and
another project which ceteris paribus she prefers (is easier to implement, is more fun, has greater
spinoffs in the future for the entrepreneur, benefits a friend, delivers perks, is more
``glamorous,'' etc.) but has a lower probability of success.

Behaving yields probability \(p = p_H\) of success and no private benefit to the entrepreneur, and
misbehaving results in probability \(p = p_L < p_H\) of success and private benefit \(B > 0\)
(measured in units of account) to the entrepreneur. In the ``effort interpretation,'' \(B\) can also
be interpreted as a disutility of effort saved by the entrepreneur when shirking. Let
\(\Delta p \equiv p_H - p_L\).

\emph{Preferences and the loan agreement.} Both the borrower and the potential lenders (or
``investors'') are risk neutral. For notational simplicity, there is no time preference; the rate of
return expected by investors (which is also the riskless rate, due to risk neutrality) is taken to
be 0. The borrower is protected by limited liability, and so her income cannot take negative values.

Lenders behave competitively in the sense that the loan, if any, makes zero profit. That is, we have
in mind that several prospective lenders compete for issuing a loan to the borrower, and that, if
the most attractive loan offer made a positive profit, the borrower could turn to an alternative
lender and offer to switch for a slightly lower interest rate. {[}\ldots{]}

Let us turn to the loan contract. A contract first stipulates whether the project is financed. If
so, it further specifies how the profit is shared between the lenders and the borrower. The
borrower's limited liability will imply that both sides receive 0 in the case of failure (the gross
payoffs are the \emph{ex post} monetary payoffs and take no account of past investments and private
benefit). Intuitively, there is no point in specifying a positive transfer from the lenders to the
borrower, as such a transfer can only weaken incentives, while it has no insurance benefit under
risk neutrality. This property will be proved more rigorously and is here taken for granted. In the
case of success, the two parties share the profit \(R\); \(R_b\) goes to the borrower and \(R_l\) to
the lenders. To sum up, we posit an incentive scheme for the entrepreneur of the following form:
\(R_b\) in the case of success, 0 in the case of failure.

The zero-profit constraint for the lenders can be written as

\[p_H R_l = I - A,\]

assuming that the loan agreement induces the borrower to behave (which under our assumptions will be
the case). {[}\ldots{]} We summarize the timing in Figure 3.1.

\emph{Figure 3.1 (timeline, left to right):} Loan agreement (sharing rule) → Investment → Moral
hazard → Outcome.

We assume that the project is viable only in the absence of moral hazard. That is, the project has
positive NPV if the entrepreneur behaves,

\[p_H R - I > 0, \qquad (3.1)\]

but negative NPV, even if one includes the borrower's private benefit, if she does not,

\[p_L R - I + B < 0. \qquad (3.2)\]

{[}\ldots{]}

\textbf{3.2.2 The Lenders' Credit Analysis.} Because the project has negative NPV in the case of
misbehavior, the loan agreement must be careful to preserve enough of a stake for the borrower in
the enterprise. The borrower faces the following tradeoff once the financing has been secured: by
misbehaving, she obtains private benefit \(B\), but she reduces the probability of success from
\(p_H\) to \(p_L\). Because she has stake \(R_b\) in the firm's income (she receives \(R_b\) in the
case of success and 0 in the case of failure), the borrower will therefore behave if the following
``incentive compatibility constraint'' is satisfied:

\[p_H R_b \geqslant p_L R_b + B \quad \text{or} \quad (\Delta p) R_b \geqslant B. \qquad (\text{IC}_b)\]

From this incentive compatibility constraint we infer that the highest income in the case of success
that can be pledged to the lenders without jeopardizing the borrower's incentives is
\(R - \dfrac{B}{\Delta p}\). The \emph{(expected) pledgeable income} is then
\(\mathcal{P} = p_H\left(R - \dfrac{B}{\Delta p}\right)\).

Because the lenders must break even in order to be willing to finance the project, a necessary
condition for the borrower to receive a loan is that the expected pledgeable income exceed the
lenders' initial outlay:

\[\mathcal{P} \equiv p_H\left(R - \frac{B}{\Delta p}\right) \geqslant I - A, \qquad (\text{IR}_l)\]

where ``IR\(_l\)'' stands for the lenders' individual rationality constraint (which we will also
often call the ``breakeven constraint'' or the ``participation constraint''). Thus a necessary
condition for financing to be arranged is

\[A \geqslant \bar{A} = p_H \frac{B}{\Delta p} - (p_H R - I). \qquad (3.3)\]

{[}\ldots{]} Conversely, if \(A \geqslant \bar{A}\), the entrepreneur is able to secure financing,
and so condition (3.3) is both a necessary and a sufficient condition for financing.

\emph{From Tirole §6.2 (privately-known-prospects model) and §6.2.1.2 (Asymmetric Information),
pp.~241--243 (PDF 254--256) --- the market-breakdown illustration the statement refers to:}

\emph{Privately-known-prospects model.} A borrower/entrepreneur has no funds (\(A = 0\)) to finance
a project costing \(I\). The project yields \(R\) in the case of success and 0 in the case of
failure. The borrower and the lenders are risk neutral, and the borrower is protected by limited
liability. The interest rate in the economy is normalized at 0.

The borrower can be one of two types. A good borrower has a probability of success equal to \(p\). A
bad borrower has a probability of success \(q\). Assume that \(p > q\) and that \(pR > I\) (at least
the good type is creditworthy). There are two subcases, which we will treat separately: either
\(pR > I > qR\) (only the good type is creditworthy), or \(pR > qR > I\) (both types are
creditworthy).

The borrower has private information about her type. The capital market, which is competitive and
demands an expected rate of return equal to 0, puts probabilities \(\alpha\) and \(1 - \alpha\) on
the borrower being a good or a bad type, respectively. Under asymmetric information, the capital
market does not know whether it faces a ``\(p\)-borrower'' (a good borrower) or a ``\(q\)-borrower''
(a bad borrower). Let

\[m \equiv \alpha p + (1 - \alpha) q\]

denote the investors' prior probability of success. {[}\ldots{]}

\textbf{6.2.1.2 Asymmetric Information.} {[}\ldots{]} Let us assume that the only feasible financial
contracts are contracts that give the borrower a compensation \(R_b \geqslant 0\) in the case of
success and 0 in the case of failure. (The validity of this assumption will be discussed in the
remark below on the optimality of contracts.) Such contracts necessarily pool the two types of
borrower as each prefers receiving financing to not being funded, and conditional on being funded,
prefers contracts with a higher compensation. The investors' profit for such a contract is therefore
on average:

\[[\alpha p + (1 - \alpha) q](R - R_b) - I = m(R - R_b) - I.\]

\emph{No lending: \(mR < I\).} This case can arise only if the bad borrower is not creditworthy. It
then arises whenever the probability that the borrower is a bad borrower is large enough, or
\(\alpha < \alpha^*\), where

\[\alpha^*(pR - I) + (1 - \alpha^*)(qR - I) = 0.\]

Because the borrower cannot receive a negative compensation (\(R_b \geqslant 0\)), investors lose
money if they choose to finance the project. Accordingly they do not and the market breaks down.

The good borrower is therefore hurt by the suspicion that she might be a bad one. There is
underinvestment.

\emph{Lending: \(mR \geqslant I\).} This case corresponds either to the situation in which both
types are creditworthy or to that in which the bad borrower is not creditworthy but
\(\alpha \geqslant \alpha^*\). The borrower's compensation \(R_b\) is then set so that investors
break even on average:

\[m(R - R_b) = I.\]

This implies that, ex post, investors make money on the good type (\(p(R - R_b) > I\)) and lose
money on the bad type (\(q(R - R_b) < I\)): there is cross-subsidization.
\stmtend
%% END-STATEMENT

\asked (i) Explain why lenders cannot fund only the good (low-$B$) borrower, and work out what happens to lenders' profit for a success reward $R_b$ below $B_L/\Delta p$, above $B_H/\Delta p$, and in between. (ii) Show that there is a cutoff $\alpha^*\in(0,1)$ for the share of good borrowers: no funding below it, funding is an equilibrium at or above it. (iii) Describe who subsidises whom when funding occurs.

\approach This is the lemons/pooling model of lecture block 3.2, with the type defined by the private benefit (moral hazard parameter) instead of the success probability. The borrower's incentive constraint (IC$_b$), $(\Delta p)R_b\ge B$, decides which types work at each $R_b$; investors' break-even then depends on the pool's success probability. Because the bad type can always copy the good one, any funded contract pools both types, so the analysis reduces to one function, lenders' pooled profit, and its sign as $\alpha$ varies (a market-breakdown threshold, as in Tirole \S6.2.1).

\notation
\begin{itemize}
  \item $I$: investment; $A=0$; $R>0$: income on success (0 on failure); $p_H>p_L$: success probability when working/shirking; $\Delta p\equiv p_H-p_L$.
  \item $B\in\{B_L,B_H\}$, $0<B_L<B_H$: private benefit from shirking (good/bad borrower); $\alpha\in(0,1)$: prior that $B=B_L$.
  \item $R_b\in[0,R]$: borrower's reward on success (0 on failure); $R-R_b$: lenders' claim.
  \item $m\equiv\alpha p_H+(1-\alpha)p_L$: success probability of the pool when the good type works and the bad type shirks.
  \item $\alpha^*$: the funding threshold; $R_b^*$: reward under competitive pooled lending.
\end{itemize}

\step{1}{Benchmark and why only good types cannot be financed} With reward $R_b$ in success and 0 in failure, type $B$ gets $p_HR_b$ by working and $p_LR_b+B$ by shirking, so she works iff
\begin{align}
  p_HR_b \ge p_LR_b+B \iff (p_H-p_L)R_b\ge B \iff (\Delta p)R_b\ge B. \tag{IC$_b$}
\end{align}
Under symmetric information lenders can be repaid at most $R-B/\Delta p$ in success without breaking (IC$_b$), so type $B$ is funded iff its pledgeable income covers $I$, $p_H(R-B/\Delta p)\ge I$; by (1) the good type is and the bad type is not. Under asymmetric information a bad borrower can offer any contract the good one offers. If she offered a different contract she would reveal her type and, by (1), not be funded; any funded contract gives her at least the good type's payoff (same if she works, more if she shirks, since $\max\{p_HR_b,p_LR_b+B_H\}\ge\max\{p_HR_b,p_LR_b+B_L\}$). So any contract that funds the good type also attracts the bad type: lenders cannot fund good types alone.

\step{2}{(i) The three ranges of $R_b$} For each range, (IC$_b$) tells us which types work; lenders' profit follows.

\textit{(a) $R_b<B_L/\Delta p$.} Then $(\Delta p)R_b<B_L<B_H$, so both types shirk and succeed with $p_L$. Since $R_b\ge0$, and by (2),
\begin{align}
  p_L(R-R_b)-I \le p_L R - I < 0. \label{t3e6-a}
\end{align}
\textit{(b) $R_b\ge B_H/\Delta p$.} Then $(\Delta p)R_b\ge B_H>B_L$, so both types work and succeed with $p_H$: the claim is information-insensitive. Since $R-R_b\le R-B_H/\Delta p$, and by the left inequality of (1),
\begin{align}
  p_H(R-R_b)-I \le p_H\left(R-\frac{B_H}{\Delta p}\right)-I < 0. \label{t3e6-b}
\end{align}
\textit{(c) $B_L/\Delta p\le R_b<B_H/\Delta p$.} Then $(\Delta p)R_b\ge B_L$ but $<B_H$: the good type works and the bad type shirks. The pool succeeds with probability $\alpha p_H+(1-\alpha)p_L=m$, so lenders' expected profit is
\begin{align}
  \Pi(R_b,\alpha) = \alpha p_H(R-R_b)+(1-\alpha)p_L(R-R_b) - I = m(R-R_b) - I. \label{t3e6-c}
\end{align}
By \eqref{t3e6-a}--\eqref{t3e6-b} only range (c) can attract funds, and by Step~1 any equilibrium is \emph{pooling}: both types or neither are funded, and when funded the good type works and the bad type shirks.

\answer{(i)} Lenders cannot separate the types (the bad type copies any funded contract). Rewards below $B_L/\Delta p$ make both types shirk and rewards from $B_H/\Delta p$ up make both work, and in both ranges lenders lose money, by (2) and (1) respectively. Only $R_b\in[B_L/\Delta p,\,B_H/\Delta p)$, where the good type works and the bad type shirks, can break even, with pooled profit $m(R-R_b)-I$.

\step{3}{(ii) The threshold $\alpha^*$} In range (c), \eqref{t3e6-c} falls with $R_b$, so it is highest at the lowest reward in the range, $R_b=B_L/\Delta p$. Define
\begin{align}
  \Pi(\alpha) \equiv [\alpha p_H+(1-\alpha)p_L]\left(R-\frac{B_L}{\Delta p}\right) - I
   = [p_L+\alpha\,\Delta p]\left(R-\frac{B_L}{\Delta p}\right) - I. \label{t3e6-Pi}
\end{align}
Differentiating the second form, $\Pi'(\alpha)=\Delta p\,(R-B_L/\Delta p)$, which is positive because $p_H(R-B_L/\Delta p)>I>0$ by (1). At the end points,
\begin{align}
  \Pi(0) &= p_L\left(R-\frac{B_L}{\Delta p}\right)-I < p_LR-I < 0 \quad\text{by (2)}, \nonumber\\
  \Pi(1) &= p_H\left(R-\frac{B_L}{\Delta p}\right)-I > 0 \quad\text{by (1)}. \nonumber
\end{align}
A continuous, strictly increasing function that is negative at 0 and positive at 1 has a unique root $\alpha^*\in(0,1)$. Setting \eqref{t3e6-Pi} to zero and solving for $\alpha$:
\begin{align}
  [p_L+\alpha^*\Delta p]\left(R-\frac{B_L}{\Delta p}\right) &= I \nonumber\\
  p_L+\alpha^*\Delta p &= \frac{I}{R-B_L/\Delta p} \nonumber\\
  \alpha^* &= \frac{1}{\Delta p}\left[\frac{I}{R-B_L/\Delta p}-p_L\right]. \label{t3e6-astar}
\end{align}
\textit{No financing if $\alpha<\alpha^*$:} then $\Pi(\alpha)<0$, and every $R_b$ loses money (ranges (a), (b) by \eqref{t3e6-a}--\eqref{t3e6-b}, range (c) because $\Pi(R_b,\alpha)\le\Pi(\alpha)$). The market breaks down.

\textit{Financing if $\alpha\ge\alpha^*$:} competitive lenders break even at
\begin{align}
  m(R-R_b^*) &= I \nonumber\\
  R_b^* &= R-\frac{I}{m}. \label{t3e6-Rbstar}
\end{align}
It lies in range (c). Lower end: $R_b^*\ge B_L/\Delta p$ iff $R-I/m\ge B_L/\Delta p$ iff $m(R-B_L/\Delta p)\ge I$, i.e.\ iff $\Pi(\alpha)\ge0$, i.e.\ iff $\alpha\ge\alpha^*$. Upper end: since $m<p_H$, $I/m>I/p_H$, so $R_b^*=R-I/m<R-I/p_H$; and the left inequality of (1), $p_H(R-B_H/\Delta p)<I$, divided by $p_H$ and rearranged, reads $R-I/p_H<B_H/\Delta p$. Together, $R_b^*<B_H/\Delta p$. So the good type works, lenders break even on average, and both types offering $\{R_b^*,0\}$ is an equilibrium.

\answer{(ii)} $\alpha^*=\dfrac{1}{\Delta p}\Bigl[\dfrac{I}{R-B_L/\Delta p}-p_L\Bigr]\in(0,1)$. For $\alpha<\alpha^*$ every contract loses money and the market breaks down; for $\alpha\ge\alpha^*$ the pooling contract $R_b^*=R-I/[\alpha p_H+(1-\alpha)p_L]$ (good type works, bad type shirks) breaks even and is an equilibrium.

\step{4}{(iii) Cross-subsidies} Substituting $R-R_b^*=I/m$ from \eqref{t3e6-Rbstar}, the lenders' profit on each type is
\begin{align}
  \text{good: } p_H\left(\frac{I}{m}\right)-I &= \frac{p_HI-mI}{m} = \frac{I(p_H-m)}{m} > 0, \label{t3e6-good}\\
  \text{bad: } p_L\left(\frac{I}{m}\right)-I &= \frac{p_LI-mI}{m} = -\frac{I(m-p_L)}{m} < 0, \label{t3e6-bad}
\end{align}
and $\alpha\cdot\eqref{t3e6-good}+(1-\alpha)\cdot\eqref{t3e6-bad}=I[\alpha p_H+(1-\alpha)p_L-m]/m=0$. Lenders make money on the good type and lose on the bad type. The good type's symmetric-information reward is $R-I/p_H$ (lenders break even on her alone), so she now pays
\begin{align}
  p_H\Bigl[\Bigl(R-\frac{I}{p_H}\Bigr)-R_b^*\Bigr] = p_H\Bigl[\frac{I}{m}-\frac{I}{p_H}\Bigr] = \frac{I(p_H-m)}{m} \nonumber
\end{align}
more in expectation, exactly the lenders' profit on her. The bad type, who would not be funded at all under symmetric information, gets $p_LR_b^*+B_H>0$. There is \emph{more} lending than under symmetric information (the bad project is funded) and the quality of lending falls: NPV conditional on funding is $mR-I$ instead of $p_HR-I$.

\answer{(iii)} Good borrowers cross-subsidise bad ones: lenders earn $I(p_H-m)/m>0$ on each good borrower and lose $I(m-p_L)/m$ on each bad one, netting to zero. The bad (negative-NPV) project gets funded and the good borrower's reward falls below its symmetric-information level.

\step{5}{Numerical illustration} Take $p_H=0.8$, $p_L=0.2$ ($\Delta p=0.6$), $R=100$, $I=50$, $B_L=9$, $B_H=27$. Then
\begin{align}
  \frac{B_L}{\Delta p} = \frac{(9)}{(0.6)} = 15, \qquad \frac{B_H}{\Delta p} = \frac{(27)}{(0.6)} = 45; \nonumber
\end{align}
(1) reads $0.8(100-45)=44<50<0.8(100-15)=68$ and (2) reads $p_LR=(0.2)(100)=20<50$. From \eqref{t3e6-astar},
\begin{align}
  \alpha^* = \frac{1}{(0.6)}\left[\frac{(50)}{(100)-(15)}-(0.2)\right] = \frac{1}{0.6}\Bigl[\frac{10}{17}-\frac{1}{5}\Bigr] = \frac{1}{0.6}\cdot\frac{33}{85} = \frac{11}{17} = 0.6471. \nonumber
\end{align}
At $\alpha=0.8$:
\begin{align}
  m &= (0.8)(0.8)+(0.2)(0.2) = 0.64+0.04 = 0.68, \nonumber\\
  R_b^* &= (100)-\frac{(50)}{(0.68)} = 100-73.53 = 26.47\in[15,45), \nonumber\\
  \text{good: } & (0.8)(73.53)-50 = 58.82-50 = 8.82, \nonumber\\
  \text{bad: } & (0.2)(73.53)-50 = 14.71-50 = -35.29, \nonumber
\end{align}
which average to zero ($0.8\times8.82=0.2\times35.29=7.06$). The good type gets $0.8\times26.47=21.18$ instead of $p_HR-I=30$ under symmetric information; the bad type gets $0.2\times26.47+27=32.29$ instead of $0$.

\intuition The lender's claim is information-sensitive only in the range where it can be profitable, and there the bad type always pools. Contracts cannot screen, so the pool's quality $m$ decides everything: with few good borrowers the market breaks down (underinvestment); with many, everyone is funded at a rate that taxes good borrowers to subsidise bad ones (overinvestment). This is the lecture's $m$, $\alpha^*$, cross-subsidy logic, with types defined by the private benefit instead of the success probability.

\pitfalls
\begin{itemize}
  \item Pricing the pooled loan with $p_H$ for both types: at $R_b\in[B_L/\Delta p,B_H/\Delta p)$ the bad type \emph{shirks}, so the pool succeeds with $\alpha p_H+(1-\alpha)p_L$.
  \item Trying to screen with a menu: whatever reward funds the good type is weakly preferred by the bad type, so there is no separating contract.
  \item Deriving $\alpha^*$ at an arbitrary $R_b$: the threshold uses the most favourable contract for lenders in range (c), $R_b=B_L/\Delta p$.
\end{itemize}
\bookans Matches the book's answer (pp.~595--596): the same three cases, the same $\alpha^*$ defined by $[\alpha^*p_H+(1-\alpha^*)p_L](R-B_L/\Delta p)=I$, the pooling contract $R_b^*=R-I/[\alpha p_H+(1-\alpha)p_L]$, and the cross-subsidy and fall in lending quality.
\refbib{Tirole \S3.2.1--3.2.2 pp.~115--117 (fixed-investment model); \S6.2.1 ``Market Breakdown and Cross-Subsidization'', pp.~241--244; answer pp.~595--596.}
%% END-FROM

%% BEGIN-FROM lista-topic3-asymmetric-information-solutions.tex EX 11 AS 7
\exsec{7. Myers--Majluf: whose interest does management serve? --- Amaro de Matos 3.2.6 \#4, p.~95}
%% BEGIN-STATEMENT 7
\stmtbegin
\textbf{Source:} Amaro de Matos, Ch. 3, Exercise 3.2.6 \#4, p.~95 (PDF 113) · Answer in book: yes,
p.~260

Instead of Assumption 3.2.6, assume that the management acts in the interest of all shareholders.
How would that affect the investment decision of the management? If management acts in the interest
of the new shareholders, what is the optimal equity issuance rule?

\paragraph{Context needed to solve it}

\emph{From Amaro de Matos §3.2, p.~78 (PDF 96):}

\textbf{Assumption 3.2.1.} \emph{Managers have better information about the value of the firm than
other agents in the economy.}

\emph{From Amaro de Matos §3.2.3 (Signals Conditioned by Investment Opportunities), pp.~86--88 (PDF
104--106):}

To see how all these conclusions are derived, the model must be formalized. Let \(I\) be the
required investment and \(S\) denote the \emph{slack} of the firm (i.e., large holdings of cash or
marketable securities) or the ability to issue default-risk-free debt.

If \(S \geq I\), it will not be necessary to issue any equity at all in order to finance the
investment. The only interesting case happens when \(S < I\), where the required amount of equity to
be issued is \(E = I - S\).

As opposed to the other models previously discussed, this is a three-date model with three distinct
times---say, time 0, time 1, and time 2. Let us introduce the following notation. At time 0, the
possible value of the assets of the firm at time 1 is described by a random variable
\(\widetilde{A}\). Similarly, let \(\widetilde{B}\) denote the possible net present value of the
investment opportunity at time 1.

At time 0, the information is symmetric. All the agents, managers and market, have the same
information---namely, the distributions of \(\widetilde{A}\) and \(\widetilde{B}\) and the value of
\(S\).

At time 1, the asymmetry of information between managers and outsider investors is specified.
Assumption 3.2.1 takes the form

\textbf{Assumption 3.2.5.} \emph{At time 1, the management knows the realizations of
\(\widetilde{A}\) and \(\widetilde{B}\), denoted respectively by \(a\) and \(b\), which are unknown
to the outsider investors. Negative values of \(a\) and \(b\) are ruled out.}

At this stage, the only thing that the outsiders know is that the amount of equity to be issued is
\(E = \max[0, I - S]\). Also, at time 1, let \(P\) be the market value of the shares of old
shareholders if the stock is not issued and \(P'\) be the price of the same shares if it is issued.
At time 2, finally, all the information is disclosed.

Let \(V^\circ\) be the value of the firm for the old shareholders. If the management, knowing the
values of \(a\) and \(b\), decides not to issue equity, then the company is still owned by the same
old shareholders with assets of value \(a\) and slack \(S\). Thus

\[V^\circ = a + S.\]

If, on the other hand, knowing the value of \(a\) and \(b\), the management decides to issue equity
\(E\), then the ownership of the company is shared between the old and new shareholders. The total
value of the company will include the issued equity \(E\) and the net present value of the
investment opportunity apart from the assets of value \(a\) and slack \(S\). But the old
shareholders will have only a proportion \(P'/(P' + E)\) of all of this. Thus, in this case

\[V^\circ = \frac{P'}{P' + E}(E + S + a + b).\]

To explain the decision of the managers, the following assumption is required:

\textbf{Assumption 3.2.6.} \emph{The management acts in the best interest of the old shareholders.
Moreover, old shareholders are assumed to be passive---that is, they do not adjust their portfolios
in response to the issue--invest decision of the firm.}

If that is so, the decision to issue new equity is taken only if the old shareholders are better off
in such case---that is, only if

\[\frac{P'}{P' + E}(E + S + a + b) \geq a + S.\]

This inequality may be rewritten as

\[b \geq -E + \frac{E}{P'}(S + a) \qquad (3.20)\]

and should naturally be restricted by Assumption 3.2.5 to the region \(a \geq 0\), \(b \geq 0\).

Hence, inequality (3.20) divides the domain of the joint probability distribution of
\(\widetilde{A}\) and \(\widetilde{B}\) into two disjoint regions identified in Figure 3.1: a region
\(M'\) in which all pairs \((a, b)\) satisfy the inequality and the firm issues new equity and
invests in the new project, and a region \(M\) in which no pair \((a, b)\) satisfies the inequality
and there is no investment. In the first region, the amount of equity issued is \(E = I - S\),
whereas in the second, it is zero.

\emph{Figure 3.1 (The Regions of Investment and No-Investment Decisions According to Myers and
Majluf's Model):} horizontal axis \(a\), vertical axis \(b\), both starting at 0. An upward-sloping
straight line labelled \(b = -E + \frac{E}{P'}(S + a)\) starts on the horizontal axis at a positive
value of \(a\) (where \(b = 0\)). The area above/left of the line is labelled ``region \(M'\)''; the
area below/right of the line is labelled ``region \(M\)''.

It is simple to see that the price \(P\) for the shares when \((a, b) \in M\) is greater than the
price \(P'\) when \((a, b) \in M'\). The equilibrium prices should be given by

\[P = S + E(\widetilde{A} \mid M) \qquad (3.21)\]

and

\[P' = S + E(\widetilde{A} + \widetilde{B} \mid M').\]

For any \((a, b) \in M\), the inequality (3.20) is not satisfied and hence

\[b < -E + \frac{E}{P'}(S + a) \quad \text{or} \quad a > P'(1 + b/E) - S.\]

Since \(b/E \geq 0\) from Assumption 3.2.5, it follows that for any \((a, b) \in M\),
\(a > P' - S\). Thus \(E(\widetilde{A} \mid M) > P' - S\) and (3.21) gives directly \(P > P'\).
\stmtend
%% END-STATEMENT

\asked (a) In the Myers--Majluf model, replace ``management serves old shareholders'' (Assumption 3.2.6) with ``management serves all shareholders'' (old and new together): how does the issue-and-invest decision change? (b) If management serves the new shareholders, what is the optimal issue rule?

\approach This is the Myers--Majluf model of Amaro de Matos \S3.2.3 (lecture block 3.5). The issue price $P'$ moves wealth between old and new shareholders, while the project's NPV $b$ adds to the total. Which side of that transfer management cares about decides the rule: write each group's wealth with and without the issue, and compare. The book's rule (3.20) is the old-shareholder comparison; the all-shareholder and new-shareholder rules come from the same algebra. A last step closes the model with the book's pricing equation to see what a rational market does with the new-shareholder rule.

\notation
\begin{itemize}
  \item $I$: required investment; $S<I$: slack (cash, riskless debt capacity); $E=I-S>0$: equity to issue.
  \item $a\ge0$, $b\ge0$: realisations (known to managers at time 1) of the value of assets in place $\widetilde A$ and of the project's NPV $\widetilde B$.
  \item $P'$: market value of the old shares if equity is issued; $P$: their value if not; the new shareholders hold the fraction $E/(P'+E)$.
  \item $M'$ / $M$: regions of $(a,b)$ with / without issue; $Z\equiv\widetilde A+\widetilde B$.
\end{itemize}

\step{1}{What each group gets} If the firm does not issue, the old shareholders own $a+S$ and the potential
new shareholders keep their cash $E$. If it issues $E$ and invests, the firm is worth $E+S+a+b$, split
in proportions $P'/(P'+E)$ and $E/(P'+E)$. So
\begin{align}
  \text{old:}&\quad a+S \;\;\text{vs.}\;\; \frac{P'}{P'+E}(E+S+a+b), \label{t3e11-old}\\
  \text{new:}&\quad E \;\;\text{vs.}\;\; \frac{E}{P'+E}(E+S+a+b), \label{t3e11-new}\\
  \text{all:}&\quad a+S+E \;\;\text{vs.}\;\; E+S+a+b. \label{t3e11-all}
\end{align}
Assumption 3.2.6 uses \eqref{t3e11-old}; multiplying it by $(P'+E)/P'$ and rearranging gives rule (3.20):
\begin{align}
  \frac{P'}{P'+E}(E+S+a+b) &\ge a+S \nonumber\\
  E+S+a+b &\ge \frac{P'+E}{P'}(a+S) = (a+S)+\frac{E}{P'}(S+a) \nonumber\\
  b &\ge -E+\frac{E}{P'}(S+a). \tag{3.20}
\end{align}

\step{2}{Management serves all shareholders} By \eqref{t3e11-all}, issue and invest iff
\begin{align}
  E+S+a+b &\ge E+S+a \nonumber\\
  b &\ge 0. \label{t3e11-allrule}
\end{align}
The issue price only moves wealth between old and new shareholders, and that transfer nets out in their
joint wealth. So management takes every positive-NPV project, the first best. By Assumption 3.2.5
($b\ge0$) every firm with $S<I$ issues: $M'$ is the whole domain, $M$ is empty, the issue carries no
information, and the book's pricing equation gives the pooled price
\begin{align}
  P' = S+E(\widetilde A+\widetilde B). \label{t3e11-Pall}
\end{align}
There is no underinvestment and no price reaction to the issue.

\answer{(a)} Management then issues and invests iff $b\ge0$: every positive-NPV project is taken (first best), every firm with $S<I$ issues, and the issue carries no information.

\step{3}{Management serves the new shareholders: the issue rule} By \eqref{t3e11-new}, issue iff the new shareholders' stake is worth at least the $E$ they pay; dividing by $E/(P'+E)>0$ and subtracting $E$ from both sides,
\begin{align}
  \frac{E}{P'+E}(E+S+a+b) &\ge E \nonumber\\
  E+S+a+b &\ge P'+E \nonumber\\
  b &\ge P'-(S+a). \label{t3e11-newrule}
\end{align}
The firm issues when its post-issue value, net of the new money, is at least the price $P'$, i.e.\ when
new shares are \emph{not overpriced}.

\answer{(b)} Issue (and invest) iff $b\ge P'-(S+a)$, equivalently $a+b\ge P'-S$: only when the new shares are not overpriced.

\step{4}{Comparison with the Myers--Majluf rule} Write $a_0\equiv P'-S$. At $a=a_0$ both thresholds are
zero: \eqref{t3e11-newrule} gives $b\ge0$, and (3.20) gives $b\ge-E+\frac{E}{P'}P'=0$. The new-shareholder
threshold $P'-S-a$ falls with $a$ (slope $-1$); the old-shareholder threshold
$-E+\frac{E}{P'}(S+a)$ rises with $a$ (slope $E/P'$). Hence, for a given $P'$:
\begin{itemize}
  \item $a>a_0$ (assets in place undervalued): the new-shareholder threshold is negative, so every such firm
        issues; the Myers--Majluf rule requires $b>0$ and passes up projects. Here the new rule is \emph{less} restrictive.
  \item $a<a_0$ (assets in place overvalued): the Myers--Majluf threshold is negative, so every such firm
        issues; the new-shareholder rule requires $b\ge a_0-a>0$. Here it is \emph{more} restrictive.
\end{itemize}
For example, with $P'=10$, $S=2$, $E=5$ ($a_0=10-2=8$): at $a=4$ the thresholds are $10-2-4=4$ (new) and $-5+\frac{5}{10}(2+4)=-2$ (old); at
$a=12$ they are $10-2-12=-4$ (new) and $-5+\frac{5}{10}(2+12)=2$ (old).

\step{5}{Equilibrium price under the new-shareholder rule} Close the model with the book's pricing
equation for the issue state, $P'=S+E(\widetilde A+\widetilde B\mid M')$, and
$M'=\{a+b\ge P'-S\}$ from \eqref{t3e11-newrule}. With $c\equiv P'-S$,
\begin{align}
  P'-S &= E(\widetilde A+\widetilde B\mid \widetilde A+\widetilde B\ge P'-S) \nonumber\\
  c &= E(Z\mid Z\ge c) \nonumber\\
  0 &= E(Z-c\mid Z\ge c). \label{t3e11-fix}
\end{align}
On the event $Z\ge c$ the integrand $Z-c$ is nonnegative, so \eqref{t3e11-fix} forces $Z=c$ almost surely
on that event: only firms at the top of the support of $Z$ can issue. With a continuous distribution the
issue region has probability zero: new investors would only be offered shares that are cheap for them,
so the price rises until nothing is cheap, and the market for new equity breaks down. Every firm with
$S<I$ then passes up its project. This is worse than under Myers--Majluf, where only the projects in $M$
are lost.

\intuition Who management works for decides which way the information problem cuts. Working for old
shareholders, it sells overpriced equity and refuses to sell underpriced equity (Myers--Majluf
underinvestment). Working for all shareholders, the issue price is a transfer inside the group and only
the NPV counts. Working for the new shareholders, it sells only underpriced equity, so the rule
favours undervalued firms. But a rational market prices that in and leaves no underpriced issues.

\pitfalls
\begin{itemize}
  \item Keeping the old-shareholder rule (3.20) in part (a): for all shareholders together the issue price is an internal transfer and drops out, leaving $b\ge0$.
  \item Sign errors when isolating $b$: the old shareholders' threshold rises with $a$, the new shareholders' falls with $a$.
  \item Calling the new-shareholder rule simply ``less restrictive'': it is less restrictive only for $a>P'-S$ and more restrictive for $a<P'-S$ (Step~4).
\end{itemize}

\bookans Both rules match the book's answer (p.~260): $b\ge0$ for all shareholders and
$b\ge P'-(S+a)$ for new shareholders. The book calls the second rule ``less restrictive'' than Myers and
Majluf's. Step~4 shows this holds only for $a>P'-S$; for $a<P'-S$ it is more restrictive. The
equilibrium in Step~5 goes beyond the book.

\refbib{Amaro de Matos Ch.~3, \S3.2.3 ``Signals Conditioned by Investment Opportunities'', pp.~86--88
(rule (3.20), prices (3.21)); Appendix C.3.2, p.~260.}
%% END-FROM

% =============================================================================
\exsec{8. Ross (1977) signalling with debt --- Amaro de Matos 3.2.6 \#1, p.~95}
%% BEGIN-STATEMENT 8
\stmtbegin
\textbf{Source:} Amaro de Matos, Ch. 3, Exercise 3.2.6 \#1, p.~95 (PDF 113) · Answer in book: yes,
p.~260

Consider the signaling model of Ross (1977). Explain what happens to the optimal decision of debt
(3.14) if \(\gamma_0 = 0\). And what happens if \(\gamma_1 = 0\)?

\paragraph{Context needed to solve it}

\emph{From Amaro de Matos §3.2, p.~78 (PDF 96):}

\textbf{Assumption 3.2.1.} \emph{Managers have better information about the value of the firm than
other agents in the economy.}

\emph{From Amaro de Matos §3.2.1 (Managers' Signaling because of Fee Schedules), pp.~79--81 (PDF
97--99):}

In the model of Ross (1977), the signal for the value of the firm is the total amount of debt issued
with face value \(B\). As discussed before, the signal must imply some cost for the management,
supposedly offset by the benefits of signaling in equilibrium. Ross proposes a fee schedule
including a punishment for the case of bankruptcy.

Two additional assumptions are required: the first to allow for explicit calculations of the cost of
bankruptcy, and the second to clarify the kind of informational asymmetry incorporated.

\textbf{Assumption 3.2.2.} \emph{The returns \(X\) of the firm are represented by a random variable
uniformly distributed on \([0, k]\), where \(k\) characterizes the type of the firm and runs over
the interval \([c, d]\).}

Assumption 3.2.1 here takes the specific form:

\textbf{Assumption 3.2.3.} \emph{Managers--insiders know their own firm's \(k\) type.}

The model is defined within two points in time. Taking \(\gamma_\circ\) and \(\gamma_1\) as two real
numbers, and \(V_\circ\) as the time-0 value of the firm, the incentive schedule has two components:
one for time 0 and another for time 1. The first one is proportional to the value of the firm at
time 0, taken to present value at time 1. For a constant interest rate \(r\), the component reads

\[M_\circ = (1 + r)\gamma_\circ V_\circ. \qquad (3.7)\]

The second component is taken proportional to the \emph{expected} return minus a penalty \(L\) in
case of bankruptcy---that is, in case \(X < B\),

\[M_1 = \gamma_1 E(X - L\mathbf{1}_{\{X < B\}}),\]

where \(\mathbf{1}_{\{X < B\}}\) denotes a random variable assuming the value 1 if \(X < B\) and a
value of 0 otherwise. The higher the debt, the higher the risk of bankruptcy, and therefore the more
efficient the management should be to avoid the penalty.

If the type of the firm is known with certainty (say, \(k = t\)), the density of probability of the
returns is constant and equal to \(1/t\) in the interval \([0, t]\), and the value of the firm at
time 1 is written

\[(1 + r)V_\circ = \int_0^t x\frac{dx}{t} = \frac{t}{2},\]

which implies that

\[V_\circ = \frac{t}{2(1 + r)}. \qquad (3.8)\]

If the essence of this approach is that the debt level \(B\) should be chosen so as to signal the
type \(t\) through some function, say, \(t = a(B)\), then the incentive schedule for the manager

\[M = M_\circ + M_1\]

must be seen by an outsider as

\[M = \frac{1}{2}\gamma_\circ a(B) + \gamma_1\left[\frac{1}{2}t - L\frac{B}{t}\right]. \qquad (3.9)\]

To write \(M\) in the form (3.9), the expression (3.8) was substituted in equation (3.7) with
\(a(B)\) in the place of \(t\) and \(M_1\) was rewritten using the fact that

\[E(X) = \int_0^t x\frac{dx}{t} = \frac{t}{2} \quad \text{and} \quad E(\mathbf{1}_{\{X < B\}}) = \int_0^B \frac{dx}{t} = \frac{B}{t},\]

with \(B\) taken to be less\textsuperscript{12} than \(t\).

The signal of the true \(t\) is therefore a function \(a(B)\) at time 0. The manager chooses \(B\)
so as to maximize its schedule \(M\). The debt level \(B\) must therefore satisfy the first-order
condition\textsuperscript{13}

\[\frac{1}{2}\gamma_\circ a'(B) = \gamma_1\frac{L}{t}. \qquad (3.10)\]

However, the management will also want the signal to be efficient, which is to say that the market
perceives in equilibrium the correct type \(t\) through \(a(B)\). Hence, the condition

\[a(B) = t \qquad (3.11)\]

must be \emph{imposed}. Substituting (3.11) into (3.10) gives in differential form

\[a(B)\,da(B) = 2\frac{\gamma_1}{\gamma_\circ}L\,dB.\]

Integrating both sides, it follows that

\[\frac{1}{2}a(B)^2 = 2\frac{\gamma_1}{\gamma_\circ}LB + b, \qquad (3.12)\]

which can be solved for the signaling function \(a(B)\) up to the constant of integration \(b\):

\[a(B) = 2\sqrt{\frac{\gamma_1 L}{\gamma_\circ}}\left[B + b\frac{\gamma_\circ}{2\gamma_1 L}\right]^{1/2}. \qquad (3.13)\]

To find the constant \(b\), a \emph{boundary condition} is necessary. In the context of this model,
such a condition follows from the fact that it does not pay to the lower type \(k = c\) to signal,
that is, \(c = a(0)\), from what follows from either equation (3.12) or equation (3.13)

\[b = \frac{c^2}{2}.\]

Substitution of this value of \(b\) into (3.13) together with the equilibrium condition (3.11) gives

\[t = \sqrt{\frac{4\gamma_1 L}{\gamma_\circ}B + c^2}\]

or, equivalently,

\[B = \frac{\gamma_\circ}{4\gamma_1 L}\left[t^2 - c^2\right]. \qquad (3.14)\]

For this signaling to be an equilibrium, it should satisfy two conditions:

\begin{enumerate}
\def\labelenumi{\arabic{enumi}.}
\tightlist
\item
  It gives managers no incentive to signal a \emph{false type} as long as the probability of
  bankruptcy is different from 1, and
\item
  It permits full discrimination of the types.
\end{enumerate}

In equilibrium all types of agents other than \(c\) will pay a cost, since they will implicitly
accept a positive probability of incurring a penalty.

The optimum level of debt is used to identify the class to which the firm belongs. Note that given
the type of the firm, its value is directly given by equation (3.8), independently of its capital
structure which is compatible with the MM framework.

\emph{From Amaro de Matos, Notes to Chapter 3, p.~280 (PDF 298):}

\begin{quote}
\begin{enumerate}
\def\labelenumi{\arabic{enumi}.}
\setcounter{enumi}{11}
\item
  Otherwise, bankruptcy is certain.
\item
  This is true as far as it is an internal optimum---that is, as far as the optimal level of \(B\)
  satisfies \(0 < B < t\).
\end{enumerate}
\end{quote}

\emph{(The statement's \(\gamma_0\) is the book's \(\gamma_\circ\).)}
\stmtend
%% END-STATEMENT

\asked (a) What the optimal debt level (3.14) becomes when the manager puts no weight on the date-0 value, $\gamma_\circ=0$. (b) What it becomes when the manager puts no weight on date-1 performance (and the bankruptcy penalty), $\gamma_1=0$.

\approach Ross's model is a signalling game: the debt level $B$ is the signal, the manager's fee $M$ in (3.9) is the payoff, and separation needs a benefit of signalling ($\gamma_\circ$: a higher perceived type raises the date-0 value) and a cost that is lower for better types ($\gamma_1$: the expected bankruptcy penalty $LB/t$ falls with $t$). Setting either weight to zero removes one side of this trade-off; we redo the manager's choice from (3.9) instead of plugging zeros into (3.14), because (3.14) was derived from an interior first-order condition that may no longer hold (the book's note 13: (3.10) is valid only for $0<B<t$).

\notation
\begin{itemize}
  \item $t\in[c,d]$: true type (returns $X\sim U[0,t]$); $B\in[0,t)$: face value of debt; $L>0$: penalty to the manager if $X<B$.
  \item $\gamma_\circ,\gamma_1\ge 0$: weights on the date-0 value and on date-1 expected returns net of the penalty; $a(B)$: type the market infers from $B$.
  \item $M(B;t)=\tfrac12\gamma_\circ a(B) + \gamma_1\bigl[\tfrac12 t - LB/t\bigr]$: fee of a type-$t$ manager choosing $B$, eq.~(3.9).
\end{itemize}

\step{1}{The general case, and why it separates} With $\gamma_\circ,\gamma_1>0$, (3.14) and the market's schedule are
\begin{align}
  B(t) &= \frac{\gamma_\circ}{4\gamma_1 L}\,(t^2-c^2), \qquad a(B) = \sqrt{\frac{4\gamma_1 L}{\gamma_\circ}B + c^2}. \label{r5e8-B}
\end{align}
Differentiating $a(B)$:
\begin{align}
  a'(B) &= \frac12\Bigl(\frac{4\gamma_1 L}{\gamma_\circ}B + c^2\Bigr)^{-1/2}\frac{4\gamma_1 L}{\gamma_\circ} = \frac{2\gamma_1 L}{\gamma_\circ\,a(B)}. \nonumber
\end{align}
So the marginal fee of a type-$t$ manager who faces this schedule is
\begin{align}
  \frac{\partial M}{\partial B} &= \frac12\gamma_\circ\,a'(B) - \gamma_1\frac{L}{t}
   = \frac12\gamma_\circ\Bigl(\frac{2\gamma_1 L}{\gamma_\circ\,a(B)}\Bigr) - \gamma_1\frac{L}{t}
   = \gamma_1 L\Bigl(\frac{1}{a(B)} - \frac1t\Bigr). \label{r5e8-dM}
\end{align}
It is positive while $a(B)<t$ and negative once $a(B)>t$, so $M$ is maximized exactly where $a(B)=t$: no type gains by signalling a false type. Example: $\gamma_\circ=\gamma_1=1$, $L=10$, $c=1$, $t=2$ give
\begin{align}
  B(2) &= \frac{(1)}{4(1)(10)}\,\bigl((2)^2-(1)^2\bigr) = \frac{3}{40} = 0.075, \nonumber\\
  a(0.075) &= \sqrt{4(1)(10)(0.075)/(1) + (1)^2} = \sqrt{4} = 2, \nonumber
\end{align}
with bankruptcy probability and expected penalty
\begin{align}
  \frac{B}{t} &= \frac{(0.075)}{(2)} = 0.0375, \qquad \gamma_1\frac{LB}{t} = \frac{(1)(10)(0.075)}{(2)} = 0.375: \nonumber
\end{align}
the cost the type $t=2$ pays to be recognized.

\step{2}{(a) $\gamma_\circ=0$} The fee (3.9) loses its benefit term:
\begin{align}
  M(B;t) &= \gamma_1\Bigl[\frac12 t - \frac{LB}{t}\Bigr], \qquad \frac{\partial M}{\partial B} = -\gamma_1\frac{L}{t} < 0. \nonumber
\end{align}
Debt only raises the expected penalty, so every type chooses the corner $B=0$. This agrees with (3.14) at $\gamma_\circ=0$:
\begin{align}
  B &= \frac{(0)}{4\gamma_1 L}\,(t^2-c^2) = 0 \quad\text{for every } t. \nonumber
\end{align}
All types pool at zero debt; the market cannot tell them apart and values every firm at the prior mean.
\answer{(a)} $B=0$ for all firms: no signalling, pooling.

\step{3}{(b) $\gamma_1=0$} The cost term disappears:
\begin{align}
  M(B;t) &= \frac12\gamma_\circ\,a(B), \qquad \frac{\partial M}{\partial B} = \frac12\gamma_\circ\,a'(B) > 0 \text{ for any increasing } a(\cdot). \nonumber
\end{align}
Whatever increasing schedule the market uses, every type wants the largest $B$ available, so no schedule can separate: the bad types copy the good ones at no cost. In (3.14) the coefficient explodes,
\begin{align}
  \lim_{\gamma_1\to 0^+}\frac{\gamma_\circ}{4\gamma_1 L}(t^2-c^2) &= +\infty \quad\text{for } t>c, \nonumber
\end{align}
which is the book's ``optimal level of debt is $\infty$''. Debt beyond $t$ means certain bankruptcy (the book's note 12), which a manager with $\gamma_1=0$ does not mind.
\answer{(b)} $B\to\infty$ for all firms (pooling at the maximum debt): debt carries no information.

\intuition A signal needs a reason to send it and a cost that the bad types find harder to bear. The weight $\gamma_\circ$ is the reason (the manager cares about how the market values the firm today); the weight $\gamma_1$ makes the bankruptcy penalty count, and since $LB/t$ is smaller for better firms, it is the type-dependent cost (single crossing). Drop either, and debt says nothing: zero debt for all if there is no reason, infinite debt for all if there is no cost.

\pitfalls
\begin{itemize}
  \item Plugging $\gamma_1=0$ into (3.14) as if it were a finite number: the formula divides by $\gamma_1$; the economic content is that every type wants as much debt as possible.
  \item Using the first-order condition (3.10) at $\gamma_\circ=0$: it would read $0=\gamma_1L/t$, which has no solution; the optimum is the corner $B=0$.
  \item Forgetting that separation needs the cost to fall with the type ($LB/t$ decreasing in $t$), not just a cost.
\end{itemize}

\bookans Amaro de Matos App.~C, p.~260: with $\gamma_1=0$ ``the optimal level of debt is $\infty$ to all firms and debt cannot be used as a signal''; with $\gamma_\circ=0$ ``the optimal level of debt is zero to all firms''; a signalling model needs the manager's return to cover both periods. Matches the book's answer (p.~260).
\refbib{Amaro de Matos \S3.2.1 ``Managers' Signaling because of Fee Schedules'', pp.~79--81, and notes 12--13 to Ch.~3, p.~280; lecture Topic~3, ``Signalling with debt (Ross 1977)''.}

%% BEGIN-FROM lista-topic3-asymmetric-information-solutions.tex EX 9 AS 9
\exsec{9. Adverse selection and ratings --- Tirole 6.6, p.~276}
%% BEGIN-STATEMENT 9
\stmtbegin
\textbf{Source:} Tirole, Ch. 6, Exercise 6.6, p.~276 (PDF 289) · Answer in book: yes, p.~598

A borrower has assets \emph{A} and must find financing for a fixed investment \emph{I}
\textgreater{} \emph{A}. As usual, the project yields \(R\) (success) or 0 (failure). The borrower
is protected by limited liability. The probability of success is \(p_H\) or \(p_L\), depending on
whether the borrower works or shirks, with \(\Delta p \equiv p_H - p_L > 0\). There is no private
benefit when working. The private benefit enjoyed by the borrower when shirking is either \(b\)
(with probability \(\alpha\)) or \(B\) (with probability \(1 - \alpha\)). At the date of
contracting, the borrower knows her private benefit, but the market (which is risk neutral and
charges a 0 average rate of interest) does not know it. Assume that \(p_L R + B < I\) (the project
is always inefficient if the borrower shirks) and that

\[p_H\left(R - \frac{B}{\Delta p}\right) < I - A < p_H\left(R - \frac{b}{\Delta p}\right) \qquad (1)\]

and

\[[\alpha p_H + (1 - \alpha)p_L]\left[R - \frac{b}{\Delta p}\right] < I - A. \qquad (2)\]

\begin{enumerate}
\def\labelenumi{(\roman{enumi})}
\item
  Interpret conditions (1) and (2) and show that there is no lending in equilibrium.
\item
  Suppose now that the borrower can at cost \(r(x) = rx\) (which is paid from the cash endowment
  \(A\)) purchase a signal with quality \(x \in [0, 1]\). (This quality can be interpreted as the
  reputation or the number of rating agencies that the borrower contracts with.) With probability
  \(x\), the signal reveals the borrower's type (\(b\) or \(B\)) perfectly; with probability
  \(1 - x\), the signal reveals nothing. The financial market observes both the quality \(x\) of the
  signal chosen by the borrower and the outcome of the signal (full or no information). The borrower
  then offers a contract that gives the borrower \(R_b\) and the lenders \(R - R_b\) in the case of
  success (so, a contract is the choice of an \(R_b \in [0, R]\)). The timing is summarized in
  Figure 6.5.
\end{enumerate}

\emph{Figure 6.5}

{\def\LTcaptype{none} % do not increment counter
\begin{longtable}[]{@{}
  >{\raggedright\arraybackslash}p{(\linewidth - 4\tabcolsep) * \real{0.3333}}
  >{\raggedright\arraybackslash}p{(\linewidth - 4\tabcolsep) * \real{0.3333}}
  >{\raggedright\arraybackslash}p{(\linewidth - 4\tabcolsep) * \real{0.3333}}@{}}
\toprule\noalign{}
\begin{minipage}[b]{\linewidth}\raggedright
Step 1
\end{minipage} & \begin{minipage}[b]{\linewidth}\raggedright
Step 2
\end{minipage} & \begin{minipage}[b]{\linewidth}\raggedright
Step 3
\end{minipage} \\
\midrule\noalign{}
\endhead
\bottomrule\noalign{}
\endlastfoot
Borrower chooses quality of signal \(x\) (this quality is observed by the capital market). &
Borrower's type revealed with probability \(x\). Nothing revealed with probability \(1 - x\). &
Borrower goes to the capital market. \\
\end{longtable}
}

Look for a pure strategy, \emph{separating} equilibrium, that is, an equilibrium in which the two
types pick different signal qualities.\textsuperscript{79}

\begin{itemize}
\tightlist
\item
  Argue that the bad borrower (borrower \(B\)) does not purchase a signal in a separating
  equilibrium.
\item
  Argue that the good borrower (borrower \(b\)) borrows under the same conditions regardless of the
  signal's realization, in a separating equilibrium.
\item
  Show that the good borrower chooses signal quality \(x \in (0, 1)\) given by
  \[A = x(A - rx) + (1 - x)\left[p_L\left(R - \frac{I - A + rx}{p_H}\right) + B\right].\]
\item
  Show that this separating equilibrium exists only if \(r\) is ``not too large.''
\end{itemize}

\begin{quote}
\begin{enumerate}
\def\labelenumi{\arabic{enumi}.}
\setcounter{enumi}{78}
\tightlist
\item
  One will assume that if the signal reveals the borrower's type, the investors put probability 1 on
  this type, even when they put weight 0 on the corresponding type after observing the quality of
  the signal.
\end{enumerate}
\end{quote}
\stmtend
%% END-STATEMENT

\asked (i) Say what conditions (1) and (2) mean economically, and prove that without ratings nobody lends. (ii) With a rating of quality $x$ (cost $rx$, reveals the type with probability $x$), look for a separating equilibrium and show: (ii-1) the bad type buys no rating; (ii-2) the good type borrows on the same terms whether or not the rating reveals her; (ii-3) her rating quality $\hat x\in(0,1)$ solves the displayed no-mimicking equation; (ii-4) such an equilibrium exists only if $r$ is not too large.

\approach Part (i) is the pooling/market-breakdown logic of lecture block 3.2 with types defined by the private benefit: (IC$_b$) sorts types into workers and shirkers at each $R_b$, the bad type can copy any funded contract, and condition (2) makes the pool unprofitable. Part (ii) is costly signalling through certification (Tirole \S6.3, Application 4; lecture block 3.7). In a separating equilibrium the good type buys the cheapest rating quality at which the bad type prefers her sure outside option to gambling that the rating reveals nothing; that no-mimicking constraint binds and pins down $\hat x$. Existence then requires checking that the good type is still funded, still works, can pay the fee and does not prefer to deviate (with pessimistic off-path beliefs).

\notation
\begin{itemize}
  \item $A>0$: borrower's cash; $I>A$: fixed investment; $R>0$: income in success, $0$ in failure; market rate $0$.
  \item $p_H>p_L\ge0$: success probabilities when working / shirking; $\Delta p\equiv p_H-p_L>0$.
  \item $b<B$: private benefits from shirking of the good type (probability $\alpha\in(0,1)$) and of the bad type (probability $1-\alpha$).
  \item $R_b\in[0,R]$: borrower's reward in success (0 in failure); $R-R_b$: investors' claim.
  \item $x\in[0,1]$: signal quality, cost $rx$ paid out of $A$, $r>0$; $\hat x$: the good type's equilibrium choice.
  \item $R_b^G(x)$: the good type's reward when recognised as good with net cash $A-rx$; $U_G$: her payoff.
  \item $c(x)$: bad type's payoff when treated as good with signal quality $x$; $c_0\equiv c(0)$.
  \item $\kappa\equiv p_L/p_H\in[0,1)$; $y\equiv r\hat x$: total signalling cost.
\end{itemize}

\assume (a) The good borrower's project has positive NPV when she works, $p_HR>I$ (otherwise there is
nothing to signal). (b) Following footnote 79, a revealing signal fixes the investors' belief; at a
signal quality not chosen by anyone in equilibrium, and with no revelation, investors believe they face
type $B$ (pessimistic off-path beliefs). (c) In part (i), contracts also pay $0$ in failure, as in part (ii);
a payment in failure would only weaken incentives and reduce what investors can be paid.

\step{1}{Interpreting (1) and (2)} A borrower with private benefit $\beta\in\{b,B\}$ works if and only if
$(\Delta p)R_b\ge\beta$, so the most she can promise investors while working is the pledgeable income
\begin{align}
  \mathcal{P}(\beta) = p_H\Bigl(R-\frac{\beta}{\Delta p}\Bigr)
     = p_HR - p_H\frac{\beta}{\Delta p}. \label{t3e9-pledge}
\end{align}
Condition (1) of the statement reads $\mathcal{P}(B)<I-A<\mathcal{P}(b)$: under symmetric information the
good type could raise $I-A$ and the bad type could not. Condition (2) of the statement says that a
contract that gives the minimum incentive stake of the good type, $R_b=b/\Delta p$, and is taken by both
types (the good one works, the bad one shirks) returns to investors, at the prior, less than $I-A$.
The assumption $p_LR+B<I$ says shirking destroys value even counting the private benefit.

\step{2}{The bad type always wants what the good type wants} Under a contract $R_b$, type $\beta$
chooses the better of working and shirking:
\begin{align}
  u_\beta(R_b) = \max\{p_HR_b,\; p_LR_b+\beta\}. \label{t3e9-u}
\end{align}
Since $B>b$, \eqref{t3e9-u} gives $u_B(R_b)\ge u_b(R_b)$ for every $R_b$, and both types' outside option is
to keep $A$. Hence any contract the good type accepts is also accepted by the bad type.

\step{3}{No lending in equilibrium (part (i))} Take any $R_b$ and split by range.

\emph{(a) $R_b\ge B/\Delta p$: both types work.} Investors receive
\begin{align}
  p_H(R-R_b) \le p_H\Bigl(R-\frac{B}{\Delta p}\Bigr)=\mathcal{P}(B) < I-A, \nonumber
\end{align}
by \eqref{t3e9-pledge} and condition (1): a loss whatever the investors' beliefs.

\emph{(b) $R_b<b/\Delta p$: both types shirk.} Investors receive $p_L(R-R_b)$ from either type. If they
break even, $p_L(R-R_b)\ge I-A$, then type $\beta$ obtains
\begin{align}
  p_LR_b+\beta &\le p_LR-(I-A)+\beta \nonumber\\
  &\le (p_LR+B-I)+A \nonumber\\
  &< A && \text{(using } p_LR+B<I), \nonumber
\end{align}
so neither type accepts: there is no lending in this range.

\emph{(c) $b/\Delta p\le R_b<B/\Delta p$: the good type works, the bad type shirks.} Because
$(\Delta p)R_b<B$, \eqref{t3e9-u} gives $u_B(R_b)=p_LR_b+B>p_HR_b=u_b(R_b)$. If the good type is funded at
$R_b$ she needs $p_HR_b\ge A$, so the bad type gets $u_B(R_b)>A$ there. By (a) and (b) (and the same
computation as (b) for a recognised bad type in this range) she cannot get more than $A$ from any contract
on which she is recognised. So she offers $R_b$ too, the investors' posterior on $R_b$ equals the prior,
and their revenue is
\begin{align}
  [\alpha p_H+(1-\alpha)p_L](R-R_b)
  \le [\alpha p_H+(1-\alpha)p_L]\Bigl(R-\frac{b}{\Delta p}\Bigr) < I-A, \nonumber
\end{align}
by condition (2). Every range yields a loss or a refusal: no lending occurs in equilibrium.
$\blacksquare$

\answer{(i)} (1): with symmetric information the good type could raise $I-A$ and the bad type could not. (2): the cheapest contract that makes the good type work, taken by both types at the prior (the bad type then shirks), does not repay $I-A$. Since the bad type copies any contract the good type is funded on, every contract either loses money or is refused: no lending.

\step{4}{The bad type buys no signal (part (ii), first bullet)} In a separating equilibrium the bad type's
quality $x_B\ne\hat x$ identifies her whatever the signal shows. With net cash $A-rx_B$ she needs
$I-A+rx_B\ge I-A$, so cases (a) and (b) of Step~3 apply \emph{a fortiori}: she is not funded on terms she
accepts and keeps $A-rx_B$. This is maximised at $x_B=0$:
\begin{align}
  U_B = A. \label{t3e9-UB}
\end{align}

\answer{(ii-1)} In a separating equilibrium the bad type is identified by her choice of $x$, is not funded on terms she accepts, and so any fee is wasted: she picks $x_B=0$ and gets $A$.

\step{5}{The good type borrows on the same terms in both events (second bullet)} At $\hat x$ the
investors infer type $b$ from the choice of $x$ alone; a revealing signal shows $b$, which is the same
belief. With the same belief the competitive market accepts the same contract. The good type has net
cash $A-r\hat x$, borrows $I-A+r\hat x$, and offers the smallest stake that investors accept:
\begin{align}
  p_H(R-R_b^G) &= I-A+rx, \label{t3e9-be}\\
  R-R_b^G &= \frac{I-A+rx}{p_H}, \nonumber\\
  R_b^G(x) &= R-\frac{I-A+rx}{p_H}. \label{t3e9-RbG}
\end{align}
She works if $(\Delta p)R_b^G\ge b$; substituting \eqref{t3e9-RbG}, dividing by $\Delta p$ and multiplying by $p_H$,
\begin{align}
  (\Delta p)\Bigl(R-\frac{I-A+rx}{p_H}\Bigr) \ge b
  \iff R-\frac{b}{\Delta p}\ge\frac{I-A+rx}{p_H}
  \iff I-A+rx \le p_H\Bigl(R-\frac{b}{\Delta p}\Bigr). \label{t3e9-icG}
\end{align}
Her payoff, substituting \eqref{t3e9-RbG}, is
\begin{align}
  U_G(x) = p_H\Bigl(R-\frac{I-A+rx}{p_H}\Bigr) = p_HR-(I-A+rx) = p_HR-I+A-rx, \label{t3e9-UG}
\end{align}
decreasing in $x$: she wants the cheapest signal that separates.

\answer{(ii-2)} At $\hat x$ the investors already infer type $b$ from the choice of $x$, and a revealing signal shows the same type, so in both events she borrows on $R_b^G=R-(I-A+r\hat x)/p_H$.

\step{6}{The mimicking constraint pins down $\hat x\in(0,1)$ (third bullet)} If the bad type chooses
$x$, the signal reveals her with probability $x$ (no funding, she keeps $A-rx$); with probability $1-x$
nothing is revealed and she is treated as good, receiving \eqref{t3e9-RbG}. She then shirks, because
\begin{align}
  (\Delta p)R_b^G < B \iff p_H\Bigl(R-\frac{B}{\Delta p}\Bigr) < I-A+rx, \nonumber
\end{align}
which holds by condition (1) and $rx\ge0$. Her payoff in that event is
\begin{align}
  c(x) = p_LR_b^G(x)+B = p_L\Bigl(R-\frac{I-A+rx}{p_H}\Bigr)+B
       = \underbrace{p_L\Bigl(R-\frac{I-A}{p_H}\Bigr)+B}_{c_0}-\frac{p_L}{p_H}\,rx = c_0-\kappa rx. \label{t3e9-c}
\end{align} Her mimicking payoff is $x(A-rx)+(1-x)c(x)$ (if $c(x)<A-rx$ she would
decline the loan and get $A-rx<A$, so the condition below covers that case too). She does not mimic iff
\begin{align}
  A \ge x(A-rx)+(1-x)\Bigl[p_L\Bigl(R-\frac{I-A+rx}{p_H}\Bigr)+B\Bigr]. \label{t3e9-nm}
\end{align}
Let $\Phi(x)$ be the right side of \eqref{t3e9-nm} minus $A$. Substituting \eqref{t3e9-c} and expanding,
\begin{align}
  \Phi(x) &= xA-rx^2+(1-x)(c_0-\kappa rx)-A \nonumber\\
          &= xA-rx^2+c_0-\kappa rx-xc_0+\kappa rx^2-A \nonumber\\
          &= -r(1-\kappa)x^2+(A-c_0-\kappa r)x+(c_0-A), \label{t3e9-Phi}
\end{align}
a strictly concave quadratic ($\kappa<1$). At $x=0$, the bad type shirks under $R_b^G(0)$, so
\begin{align}
  c_0 = p_LR_b^G(0)+B > p_HR_b^G(0) = p_HR-I+A > A, \nonumber
\end{align}
using \eqref{t3e9-UG} at $x=0$ and assumption (a): $\Phi(0)=c_0-A>0$. At $x=1$, \eqref{t3e9-Phi} gives
$\Phi(1)=-r(1-\kappa)+A-c_0-\kappa r+c_0-A=-r<0$. A concave function positive at 0 and negative at 1 has a
single root $\hat x\in(0,1)$, with $\Phi>0$ on $[0,\hat x)$ and $\Phi<0$ on $(\hat x,1]$. By
\eqref{t3e9-UG} the good type picks the smallest admissible $x$, i.e.\ $\hat x$, where \eqref{t3e9-nm} binds:
\begin{align}
  A = \hat x(A-r\hat x)+(1-\hat x)\Bigl[p_L\Bigl(R-\frac{I-A+r\hat x}{p_H}\Bigr)+B\Bigr]. \qquad\blacksquare \label{t3e9-xhat}
\end{align}

\answer{(ii-3)} $\hat x\in(0,1)$ is the unique root in $(0,1)$ of the bad type's indifference condition \eqref{t3e9-xhat}, the equation in the statement: the cheapest quality at which mimicking no longer pays.

\step{7}{The separating equilibrium exists only if $r$ is not too large (fourth bullet)} Write $y=r\hat x$
in \eqref{t3e9-xhat}. Using \eqref{t3e9-c},
\begin{align}
  A &= \hat xA-\hat xy+(1-\hat x)(c_0-\kappa y) \nonumber\\
  A-\hat xA &= -\hat xy+(1-\hat x)c_0-(1-\hat x)\kappa y \nonumber\\
  (1-\hat x)A-(1-\hat x)c_0 &= -y\,[\hat x+(1-\hat x)\kappa] \nonumber\\
  (1-\hat x)(c_0-A) &= y\,[\hat x+(1-\hat x)\kappa] \nonumber\\
  y &= \frac{(1-\hat x)(c_0-A)}{\hat x+(1-\hat x)\kappa}, \label{t3e9-y}\\
  r = \frac{y}{\hat x} &= \frac{(1-\hat x)(c_0-A)}{\hat x\,[\hat x+(1-\hat x)\kappa]}. \label{t3e9-r}
\end{align}
The right side of \eqref{t3e9-r} falls strictly from $+\infty$ to $0$ as $\hat x$ goes from $0$ to $1$, so
$\hat x$ is a strictly decreasing function of $r$ ($\hat x\to1$ as $r\to0$). The right side of
\eqref{t3e9-y} is strictly decreasing in $\hat x$, so the signalling cost $y(r)$ is strictly increasing in
$r$, from $0$ (as $r\to0$) to $(c_0-A)/\kappa$ (as $r\to\infty$, if $p_L>0$).

The candidate is an equilibrium only if three inequalities hold, each an upper bound on $y$:
\begin{itemize}
  \item[(E1)] the good type is funded and works, \eqref{t3e9-icG} at $\hat x$: $y\le p_H(R-b/\Delta p)-(I-A)$, positive by (1);
  \item[(E2)] the good type does not deviate. A deviation to $x'\ne\hat x$ is treated as bad unless the signal
     reveals $b$, so it yields at most $A-rx'+x'(p_HR-I)$. Subtracting this from \eqref{t3e9-UG} at $\hat x$,
     \begin{align}
       D(x') = (1-x')(p_HR-I)-r(\hat x-x'), \nonumber
     \end{align}
     which is positive for $x'\ge\hat x$ and linear on $[0,\hat x]$ with $D(\hat x)=(1-\hat x)(p_HR-I)>0$; so
     the binding deviation is $x'=0$: $y\le p_HR-I$;
  \item[(E3)] the fee is paid from cash: $y\le A$.
\end{itemize}
The bad type's deviations give her at most $A-rx'\le A$, so \eqref{t3e9-UB} is optimal. Let
$\bar y\equiv\min\{p_H(R-b/\Delta p)-(I-A),\,p_HR-I,\,A\}>0$. Since $y(r)$ rises from $0$, the
separating equilibrium exists if and only if $y(r)\le\bar y$, i.e.\ $r\le\bar r$, where $\bar r$ solves
$y(\bar r)=\bar y$ ($\bar r=+\infty$ if $(c_0-A)/\kappa\le\bar y$). Inverting \eqref{t3e9-y} at $y=\bar y$ (multiply out and collect the terms in $\bar x$):
\begin{align}
  \bar y\,[\bar x+(1-\bar x)\kappa] &= (1-\bar x)(c_0-A) \nonumber\\
  \bar y\kappa+\bar x\,\bar y(1-\kappa) &= (c_0-A)-\bar x(c_0-A) \nonumber\\
  \bar x\,[(c_0-A)+(1-\kappa)\bar y] &= c_0-A-\kappa\bar y \nonumber\\
  \bar x = \frac{c_0-A-\kappa\bar y}{c_0-A+(1-\kappa)\bar y}, &\qquad \bar r = \frac{\bar y}{\bar x}. \label{t3e9-rbar}
\end{align}

\answer{(ii-4)} The signalling cost $y=r\hat x$ rises with $r$, and the equilibrium needs $y\le\bar y=\min\{p_H(R-b/\Delta p)-(I-A),\;p_HR-I,\;A\}$ (good type works, does not prefer to deviate to $x=0$, can pay the fee). Hence it exists iff $r\le\bar r=\bar y/\bar x$ in \eqref{t3e9-rbar}.

\step{8}{Numerical illustration} Take $R=10$, $I=6$, $A=1$, $p_H=0.8$, $p_L=0.4$ ($\Delta p=0.4$), $b=0.4$,
$B=1.8$, $\alpha=0.3$, so $b/\Delta p=1$, $B/\Delta p=4.5$ and $I-A=5$. Condition (1): $0.8(10-4.5)=4.4<5<0.8(10-1)=7.2$. Condition (2):
$[0.3(0.8)+0.7(0.4)](10-1)=(0.24+0.28)(9)=(0.52)(9)=4.68<5$; and $p_LR+B=4+1.8=5.8<6$. Then
\begin{align}
  c_0 &= (0.4)\Bigl(10-\frac{5}{0.8}\Bigr)+1.8 = (0.4)(3.75)+1.8 = 1.5+1.8 = 3.3, \qquad \kappa=\frac{0.4}{0.8}=0.5. \nonumber
\end{align}
With $r=1$, \eqref{t3e9-Phi} has coefficients $-r(1-\kappa)=-0.5$, $A-c_0-\kappa r=1-3.3-0.5=-2.8$ and $c_0-A=2.3$:
\begin{align}
  -0.5x^2-2.8x+2.3 &= 0 \iff x^2+5.6x-4.6=0, \nonumber\\
  \hat x &= \frac{-5.6+\sqrt{5.6^2+4(4.6)}}{2} = \frac{-5.6+\sqrt{49.76}}{2} = -2.8+\sqrt{12.44} = 0.727038. \nonumber
\end{align}
Then $y=r\hat x=0.727038$, $U_G=p_HR-I+A-y=8-6+1-0.727038=2.272962$
and $R_b^G=10-(5+0.727038)/0.8=10-7.158798=2.841202\ge b/\Delta p=1$. Here
\begin{align}
  \bar y=\min\{7.2-5,\;8-6,\;1\}=\min\{2.2,\,2,\,1\}=1, \nonumber
\end{align}
so \eqref{t3e9-rbar} gives $\bar x=(2.3-0.5)/(2.3+0.5)=1.8/2.8=0.642857$ and $\bar r=1/0.642857=1.555556$: the rating equilibrium exists for
$r\le1.555556$.

\intuition Lending collapses because whatever contract attracts the good borrower attracts the bad
one, and condition (2) says the pool is unprofitable. A rating of quality $x$ lets the good type escape the
pool, but only partly: the bad type can buy the same rating and hope that it reveals nothing. The good type
therefore buys just enough quality that the gamble ``with probability $1-x$ I pass as good'' is not worth
the bad type's sure $A$. The cost $rx$ is pure dissipation, and when it eats the good type's surplus (or
cash, or incentive margin) separation fails.

\pitfalls
\begin{itemize}
  \item Forgetting that the fee $rx$ is paid out of $A$, so the good type must borrow $I-A+rx$, not $I-A$; this enters both her terms and the bad type's mimicking payoff.
  \item Writing the mimicking payoff as if the bad type worked under $R_b^G$: since $R_b^G<B/\Delta p$ she shirks, and her payoff is $p_LR_b^G+B$.
  \item Choosing $x=1$ (``buy the best rating''): the good type wants the \emph{cheapest} separating quality, so the no-mimicking constraint binds at $\hat x<1$.
\end{itemize}

\bookans Matches the book's answer (pp.~598--599): the bad type chooses $x=0$ and is not funded, the
good type's terms solve $p_H(R-R_b^G)=I-A+rx$, and the no-mimicking condition is the displayed equation,
binding at the equilibrium. The book does not spell out the ``$r$ not too large'' conditions; (E1)--(E3)
and \eqref{t3e9-rbar} are derived here.

\refbib{Tirole Ch.~6, \S6.3 Application 4 ``Certification'', pp.~249--251; Exercise 6.6, p.~276;
answer pp.~598--599. Lecture Topic~3, signalling block.}
%% END-FROM

%% BEGIN-FROM lista-topic3-asymmetric-information-solutions.tex EX 10 AS 10
\exsec{10. Maturity structure --- Tirole 6.10, p.~279}
%% BEGIN-STATEMENT 10
\stmtbegin
\textbf{Source:} Tirole, Ch. 6, Exercise 6.10, p.~279 (PDF 292) · Answer in book: yes, p.~600

At date 0 the entrepreneur has cash on hand \emph{A} and needs to finance an investment of fixed
size \(I\). At date 1, a deterministic income \(r\) accrues; a liquidity shock must be met in order
for the firm to continue. Liquidation yields nothing. The probability of success in the case of
continuation depends on a date-1 effort: for a good borrower, this probability is \(p_H\) or \(p_L\)
depending on whether she behaves (no private benefit) or misbehaves (private benefit \(B\));
similarly, for a bad borrower, it is \(q_H\) or \(q_L\). We assume that

\[p_H - p_L = q_H - q_L = \Delta p,\]

and so the incentive compatibility constraint in the case of continuation is the same for both types
of borrower:

\[(p_H - p_L)R_b = (q_H - q_L)R_b = (\Delta p)R_b \geqslant B,\]

where \(R_b\) is the borrower's reward in the case of success.

The borrower knows at the date of contracting whether she is a ``\(p\)-type'' or a ``\(q\)-type.''
Let

\[\rho_0^G \equiv p_H\left(R - \frac{B}{\Delta p}\right) \quad \text{and} \quad \rho_0^B = q_H\left(R - \frac{B}{\Delta p}\right)\]

denote the date-1 pledgeable incomes for the good and bad types.

The liquidity shock is deterministic and equal to \(\rho\). Information is asymmetric at date 0, but
the capital market learns the borrower's type perfectly at date 1, before the liquidity shock has to
be met. Assume that

\[\rho_0^G > \rho > \rho_0^B.\]

Suppose further that under symmetric information only the good borrower is creditworthy (provided
that she is incentivized to behave).

Assume that \(r < I - A < r + [\rho_0^G - \rho]\). Show that the good borrower can costlessly signal
her type.

\paragraph{Context needed to solve it}

\emph{From Tirole §6.3, Application 6 ``Short-Term Maturities'', pp.~254--255 (PDF 267--268) --- the
model of which the exercise is a variant (there the liquidity shock is random; in the exercise it is
deterministic and the asymmetry is about the success probabilities). It defines \(R\) and the
preferences the statement leaves implicit:}

Let us consider the following variant of the model of debt maturities set up in Section 5.2. At date
0 the entrepreneur has a project of fixed size \(I\), has wealth \(A\), and must borrow \(I - A\).
At date 1, the investment yields a deterministic and verifiable short-term (date-1) profit
\(r > 0\). With probability \(\lambda\), continuation requires reinvesting \(\rho\) (the liquidation
value is equal to 0); with probability \(1 - \lambda\), continuation requires no
reinvestment.\textsuperscript{51}

In the case of continuation, the firm at date 2 yields \(R\) in the case of success and 0 in the
case of failure. The probability of success is \(p_H\) in the case of good behavior and
\(p_L = p_H - \Delta p\) in the case of misbehavior (which yields private benefit \(B\)). Moral
hazard is introduced in order to create an entrepreneurial rent from continuation, or, equivalently,
a cost for the entrepreneur associated with early termination. Put differently, moral hazard
introduces a friction in the date-1 refinancing market. Let

\[\rho_1 \equiv p_H R \quad \text{and} \quad \rho_0 \equiv p_H\left(R - \frac{B}{\Delta p}\right)\]

denote the continuation NPV and pledgeable income, respectively.

{[}\ldots{]} the borrower can use retentions \((r - d)\) together with what can be raised in the
capital market \((\rho_0)\) to meet shock \(\rho\).

\begin{quote}
\begin{enumerate}
\def\labelenumi{\arabic{enumi}.}
\setcounter{enumi}{50}
\tightlist
\item
  As usual, the entrepreneur is risk neutral and protected by limited liability; and the investors
  are risk neutral and demand a 0 rate of return.
\end{enumerate}
\end{quote}

\emph{Figure 6.1 (timeline, left to right):} Date 0: Entrepreneur has wealth \(A\) and fixed-size
investment project \(I - A\). → Date 1: Short-run income \(r\). Reinvestment need \(\rho\) (prob.
\(\lambda\)) or 0 (prob. \(1 - \lambda\)). → (if reinvestment) Moral hazard (\(p = p_H\) or
\(p_L\)). → Outcome: success (\(R\)) or failure (0).
\stmtend
%% END-STATEMENT

\asked (a) Two types of entrepreneur differ only in their date-2 success probabilities; the market learns the type at date 1, before a deterministic liquidity shock $\rho$ must be paid. Only the good type can refinance the shock in the market ($\rho_0^G>\rho>\rho_0^B$) and only she is creditworthy. Show that she can separate herself from the bad type at no cost, i.e.\ while obtaining her symmetric-information payoff.

\approach This is signalling through the maturity of debt (Tirole \S6.3, Application 6, built on the liquidity model of Ch.~5; lecture block 3.7). A signal separates if it is cheap for the good type and expensive for the bad one. Short-term debt that takes all the date-1 income ($d=r$) leaves no hoarded liquidity, so the firm must go back to the market at date 1, after its type is known. The good type can refinance there; the bad type cannot and is liquidated. So the good type offers that contract, checks that it gives her the symmetric-information payoff (incentive and break-even constraints), and checks that the bad type does not want to copy it (no-mimicking, including renegotiation), with pessimistic off-path beliefs.

\notation
\begin{itemize}
  \item $A$: cash at date 0; $I$: fixed investment, $I-A>0$ borrowed; $r>0$: deterministic date-1 income.
  \item $\rho$: deterministic date-1 liquidity shock; liquidation value $0$.
  \item $p_H,p_L$ (good) and $q_H,q_L$ (bad): date-2 success probabilities when behaving / misbehaving; $p_H-p_L=q_H-q_L=\Delta p$.
  \item $B$: private benefit from misbehaving; $R$: date-2 income in success, 0 in failure; $R_b$: borrower's reward in success.
  \item $\rho_0^G=p_H(R-B/\Delta p)$, $\rho_0^B=q_H(R-B/\Delta p)$: date-1 pledgeable incomes, $\rho_0^G>\rho>\rho_0^B$.
  \item $d\in[0,r]$: short-term debt repaid at date 1 (retention $r-d$); $L$: date-0 investors' long-term claim on date-2 income (expected value).
  \item $R_n$, $R_l$: shares of $R$ promised to date-1 (new) and date-0 (old) investors.
\end{itemize}

\assume (a) As in Tirole's Application 6, a contract specifies $d$, a long-term claim and $R_b$; new claims
issued at date 1 may dilute the old long-term claim. (b) A continuation in which the borrower misbehaves
cannot be financed ($p_LR<\rho$ and $q_LR<\rho$), so date-1 investors fund only continuations with
$(\Delta p)R_b\ge B$. (c) Off-path beliefs: any date-0 offer other than the good type's equilibrium
contract is attributed to the bad type.

\step{1}{Symmetric-information benchmark} With the type known, the good borrower can always be continued
(the shock is deterministic and $\rho_0^G>\rho$, so the date-1 market covers it). Investors break even, so
her gross payoff is the project's value net of all outlays:
\begin{align}
  U_G^{SI} = A + r + p_HR - \rho - I. \label{t3e10-USI}
\end{align}
Funding is feasible because total pledgeable income covers the loan: investors can receive the date-1 income $r$ and, at date 2, the pledgeable income $\rho_0^G$ minus the $\rho$ that must be paid to whoever finances the shock, so they break even iff
\begin{align}
  r+(\rho_0^G-\rho)\ge I-A, \nonumber
\end{align}
which holds by the right inequality of the statement. The bad type is not creditworthy (statement), so her
symmetric-information payoff is to keep her cash:
\begin{align}
  U_B^{SI} = A. \label{t3e10-UBSI}
\end{align}

\step{2}{The signalling contract: short-term debt $d=r$} The good type offers: short-term debt $d=r$ (no
retention, no hoarded liquidity), a dilutable long-term claim worth
\begin{align}
  L = I-A-r > 0 \qquad (\text{by } r<I-A), \label{t3e10-L}
\end{align}
and a reward $R_b$ in success. Date-0 investors receive $d+L=r+(I-A-r)=I-A$: they break even. At date 1
the market learns she is good and she must raise $\rho$ by issuing a new claim $R_n$ with
$p_HR_n=\rho$. The old claim is $p_HR_l=L$. The borrower keeps the rest, so
\begin{align}
  p_HR_b &= p_HR-p_HR_n-p_HR_l \nonumber\\
         &= p_HR-\rho-L \nonumber\\
         &= p_HR-\rho-(I-A-r) && \text{(substituting \eqref{t3e10-L})} \nonumber\\
         &= A+r+p_HR-\rho-I = U_G^{SI} && \text{(by \eqref{t3e10-USI})}. \label{t3e10-UG}
\end{align}
She behaves if $(\Delta p)R_b\ge B$, i.e.\ (multiplying by $p_H/\Delta p$) $p_HR_b\ge p_HB/\Delta p$; by the second line of
\eqref{t3e10-UG} this is
\begin{align}
  p_HR-\rho-L &\ge p_H\frac{B}{\Delta p} \nonumber\\
  p_H\Bigl(R-\frac{B}{\Delta p}\Bigr) &\ge \rho+L \iff \rho_0^G\ge\rho+(I-A-r) \nonumber\\
  I-A &\le r+(\rho_0^G-\rho), \label{t3e10-ic}
\end{align}
which is the right inequality assumed in the statement. So the contract is feasible and gives the good
type exactly her symmetric-information payoff: signalling is costless.

\step{3}{The bad type does not mimic} Suppose the bad type offers the same contract. At date 1 all her
income $r$ goes to repay $d=r$. The market learns she is bad; even with the old claim fully diluted, the
most new investors can be promised while she behaves is $\rho_0^B<\rho$, and by assumption (b) a misbehaving
continuation cannot be funded either. The project is liquidated and yields $0$, so
\begin{align}
  U_B^{\text{mimic}} = 0 \le A = U_B^{SI}. \label{t3e10-mimic}
\end{align}
Could the old investors waive $d$ to let her continue? Then the firm would keep $r$ and would need
$\rho-r$ from the market; out of the bad type's date-2 pledgeable income $\rho_0^B$, the new money takes $\rho-r$, so the old investors could get at most
\begin{align}
  \rho_0^B-(\rho-r)=r-(\rho-\rho_0^B)<r, \nonumber
\end{align}
less than collecting $d=r$ (because $\rho>\rho_0^B$). So they do not renegotiate, and \eqref{t3e10-mimic} stands.

\step{4}{The separating equilibrium} The good type offers the Step~2 contract and the bad type stays out.
Investors believe the Step~2 contract comes from the good type (correct, by Step~3) and fund it; any other
offer is attributed to the bad type and, since she is not creditworthy, is not funded. The good type
cannot gain by deviating: she already gets $U_G^{SI}$ in \eqref{t3e10-UG}, and other offers are not funded.
The bad type gets $A$, by \eqref{t3e10-mimic} at least as much as by mimicking. Hence the good borrower
signals her type at zero cost. $\blacksquare$

Signalling works because liquidity and refinancing are worth different amounts to the two types. A
contract that hoarded enough liquidity, $r-d\ge\rho-\rho_0^B$, would let a mimicking bad type meet the
shock with retentions plus new money, so she would get a rent instead of 0. The good type gives up
nothing by not hoarding, because the date-1 market will see that she is good and lend to her.

\answer{(a)} The good type offers short-term debt $d=r$ (no hoarded liquidity) plus a dilutable long-term claim worth $L=I-A-r$. Date-0 investors break even, she refinances $\rho$ at date 1 once revealed good and obtains $p_HR_b=A+r+p_HR-\rho-I=U_G^{SI}$, so the signal costs her nothing. A mimicking bad type would be liquidated at date 1 (payoff $0<A$), so she stays out: a costless separating equilibrium.

\step{5}{Numerical illustration} Take $p_H=0.8$, $q_H=0.5$, $\Delta p=0.4$, $B=1.2$ ($B/\Delta p=1.2/0.4=3$),
$R=10$, $\rho=4.5$, $r=1$, $I-A=1.6$. Then $\rho_0^G=0.8(10-3)=0.8(7)=5.6$ and $\rho_0^B=0.5(7)=3.5$, so
$\rho_0^G>\rho>\rho_0^B$; $r=1<I-A=1.6<r+\rho_0^G-\rho=1+5.6-4.5=2.1$; and the bad type fails even with full hoarding,
$r+\rho_0^B-\rho=1+3.5-4.5=0<1.6$. By \eqref{t3e10-L}, $L=1.6-1=0.6$. Then
\begin{align}
  R_n &= \frac{\rho}{p_H} = \frac{(4.5)}{(0.8)} = 5.625, \qquad R_l = \frac{L}{p_H} = \frac{(0.6)}{(0.8)} = 0.75, \nonumber\\
  p_HR_b &= r+p_HR-\rho-(I-A) = 1+8-4.5-1.6 = 2.9, \qquad R_b = \frac{2.9}{0.8} = 3.625\ge\frac{B}{\Delta p}=3, \nonumber
\end{align}
and $R_b+R_n+R_l=3.625+5.625+0.75=10=R$.

\intuition Short-term debt drains the firm's cash at date 1 and forces it back to the capital market just
after the market has learned the firm's type. The good type is not afraid of that meeting; the bad type is,
because the market will refuse to fund her shock. Short maturity is an action whose cost falls only on
the type that has something to hide, which is why it separates. Here it costs the good type nothing only
because the shock is certain and small enough to refinance. With a random shock that may exceed
$\rho_0^G$, she would want to hoard liquidity even under symmetric information, and the signal would cost
her some continuation (Application 6).

\pitfalls
\begin{itemize}
  \item Thinking the good type should hoard liquidity, as in Ch.~5: with a certain shock and a market that will know she is good, hoarding is unnecessary, and it would let the bad type survive by mimicking.
  \item Forgetting the date-1 claim structure: the old long-term claim is diluted by the new money raised for $\rho$, and the incentive constraint must hold \emph{after} both claims are paid ($\rho+L\le\rho_0^G$).
  \item Ignoring renegotiation: one must check that old investors would not waive $d$ to keep a revealed bad type alive.
\end{itemize}

\bookans Matches the book's answer (p.~600): separation by not hoarding liquidity, $d=r$, the good type
refinancing at date 1 and diluting existing claims, the bad type liquidated. The long-term claim
$L=I-A-r$, the renegotiation check and assumption (b) are added here; the book leaves them implicit.

\refbib{Tirole Ch.~6, \S6.3 Application 6 ``Short-Term Maturities'', pp.~254--257; Ch.~5, \S5.2
(liquidity hoarding and short-term debt); Exercise 6.10, p.~279; answer p.~600.}
%% END-FROM

% =============================================================================
\exsec{11. Herd behavior --- Tirole 6.9, pp.~278--279}
%% BEGIN-STATEMENT 11
\stmtbegin
\textbf{Source:} Tirole, Ch. 6, Exercise 6.9, pp.~278--279 (PDF 291) · Answer in book: yes, p.~600

It is often argued that the managers of industrial companies, banks, or mutual funds are prone to
herd.\textsuperscript{81} They engage in similar investments with sometimes little evidence that
their strategy is the most profitable. An economic agent may indeed select a popular strategy
against her own information that another strategy may be more profitable. A number of contributions
have demonstrated that herding behavior may actually be individually rational even though it is
often collectively inefficient. The literature on herding behavior starts with the seminal
contributions of Banerjee (1992), Bikhchandani et al.~(1992), Scharfstein and Stein (1990), and
Welch (1992); see Bikhchandani and Sharma (2001) for a survey of applications of this literature to
financial markets.

There are several variants of the following basic argument. Consider first a sequence of agents
\(i = 1, 2, \dots\) choosing sequentially between strategies A and B. Agents receive their own
signals; they observe previous decisions but not the others' signals. Suppose that agents 1 and 2
have, on the basis of their own information, selected A. Agent 3, observing the first two choices,
may well then select strategy A even if her own signal favors the choice of B. Agent 4, not knowing
agent 3's motivation to choose A, may then also choose strategy A even if his own signal points
toward the choice of B. And so forth. It may therefore be the case that all agents choose A, even
though the cumulative evidence, if it were shared, would indicate that B is the best choice.

The literature also analyzes herd behavior in situations in which agents have principals (that is,
they are not full residual claimants for the consequences of their choices). In particular, such
agents may adopt herd behaviors because of \emph{reputational concerns} (see Chapter 7). Suppose,
for instance, that a manager's job is rather secure; herding with the managers of other firms is
then likely to be attractive to the manager: if the strategy fails, the manager has the excuse that
other managers also got it wrong (``it was hard to predict''). Choosing an unpopular strategy, even
if one's information points in that direction, is risky, as there will be no excuse if it fails. The
literature on herd behavior has also investigated the use of benchmarking by principals in
\emph{explicit incentives} (compensation contracts) rather than in implicit ones (career concerns).

Let us build an example of herding behavior in the context of the privately-known-prospects model of
Section 6.2. There are two entrepreneurs, \(i = 1, 2\), operating in different markets, but whose
optimal strategy is correlated. There are two periods, \(t = 1, 2\). Entrepreneur \(i\) can raise
funds only at date \(t = i\) (so they secure funding sequentially). A project yields \(R\) when it
succeeds and 0 when it fails. The entrepreneurs are risk neutral and protected by limited liability;
the investors are risk neutral and demand a rate of return equal to 0. The entrepreneurs have no net
worth or cash initially.

The two entrepreneurs each have to choose between strategy A and B. Strategies differ in their
probability of success. A borrowing contract with investors specifies both the managerial
compensation \(R_b\) in the case of success (and 0 in the case of failure) and the strategy that the
entrepreneur will select.\textsuperscript{82} Crucially, entrepreneur 2 and her potential investors
observe the date-1 financing contract for entrepreneur 1. Entrepreneurs, but not investors, learn
the state of nature.

Consider the following stochastic structure.

\emph{Unfavorable environment (probability \(1 - \alpha\)).} The probabilities of success are, with
equal probabilities, \((q, 0)\) for one project and \((0, q)\) for the other, where the first
element is entrepreneur 1's probability of success and the second entrepreneur 2's. So entrepreneurs
necessarily choose different projects if they apply for funding.

\emph{Favorable environment (probability \(\alpha\)).} With probability \(\theta\), the best project
is the same for both and has probability of success \(p\); the worst project for both has
probability of success \(r\), where

\[p > \max\{q, r\}.\]

With probability \(1 - \theta\), the two entrepreneurs' best strategies differ: the probabilities of
success are \((p, r)\) and \((r, p)\), respectively, for entrepreneur 1's and entrepreneur 2's best
strategy (which are A or B with equal probabilities). Thus \(\theta\) is the probability of
correlation of the best strategies in a favorable environment; this probability is equal to 0 in the
unfavorable environment.

Let \(m \equiv \alpha p + (1 - \alpha)q\) and assume that

\[qR > I.\]

Show that funding and herding (with probability \(\alpha(1 - \theta)\), entrepreneur 2 chooses
entrepreneur 1's best strategy even though it does not maximize her probability of success) is an
equilibrium behavior as long as

\[r\left[R - \frac{I}{\theta p + (1 - \theta)r}\right] \geqslant p\left[R - \frac{I}{q}\right].\]

Note that entrepreneur 2 is on average worse off than in an hypothetical situation in which
investors did not observe the strategy of entrepreneur 1 (or that in which the optimal strategies
were uncorrelated).

\begin{quote}
\begin{enumerate}
\def\labelenumi{\arabic{enumi}.}
\setcounter{enumi}{80}
\item
  One of the first empirical papers on herding behavior is Lakonishok et al.~(1992). The large
  empirical literature on the topic includes Chevalier and Ellison (1999).
\item
  One will assume therefore that the first entrepreneur cannot condition her financing contract on
  the later choice of strategy by the second entrepreneur.
\end{enumerate}
\end{quote}

\paragraph{Context needed to solve it}

\emph{From Tirole §6.2 (Implications of the Lemons Problem and of Market Breakdown), pp.~241--242
(PDF 254--255):}

\emph{Privately-known-prospects model.} A borrower/entrepreneur has no funds (\(A = 0\)) to finance
a project costing \(I\). The project yields \(R\) in the case of success and 0 in the case of
failure. The borrower and the lenders are risk neutral, and the borrower is protected by limited
liability. The interest rate in the economy is normalized at 0.

The borrower can be one of two types. A good borrower has a probability of success equal to \(p\). A
bad borrower has a probability of success \(q\). Assume that \(p > q\) and that \(pR > I\) (at least
the good type is creditworthy). There are two subcases, which we will treat separately: either
\(pR > I > qR\) (only the good type is creditworthy), or \(pR > qR > I\) (both types are
creditworthy).

The borrower has private information about her type. The capital market, which is competitive and
demands an expected rate of return equal to 0, puts probabilities \(\alpha\) and \(1 - \alpha\) on
the borrower being a good or a bad type, respectively. Under asymmetric information, the capital
market does not know whether it faces a ``\(p\)-borrower'' (a good borrower) or a ``\(q\)-borrower''
(a bad borrower). Let

\[m \equiv \alpha p + (1 - \alpha) q\]

denote the investors' prior probability of success.

Note that we have left out for the moment moral hazard in the definition of the
privately-known-prospects model. The coexistence of moral hazard with adverse selection is not
necessary for most applications (for which one can therefore ignore private benefits, \(B = 0\) in
the notation of the book, and thereby remove the moral-hazard component) since adverse selection by
itself creates an agency cost and a concomitant credit rationing, and triggers a number of
interesting institutional responses. Ignoring moral hazard therefore simplifies the presentation.
{[}\ldots{]} Note also that we assume that the entrepreneur has no cash on hand (\(A = 0\)), and so
she cannot signal her trust in the project by investing her personal wealth into it. {[}\ldots{]}
\stmtend
%% END-STATEMENT

\asked Show that the following is an equilibrium: both entrepreneurs are funded, and in the favourable environment with uncorrelated best strategies (probability $\alpha(1-\theta)$) entrepreneur~2 copies entrepreneur~1's strategy although her own best strategy is the other one, provided $r[R-I/(\theta p+(1-\theta)r)]\ge p[R-I/q]$. Then show that entrepreneur~2 is worse off on average than if investors could not see entrepreneur~1's strategy.

\approach This is the privately-known-prospects model of Tirole \S6.2 with pooling contracts: competitive investors set the entrepreneur's reward $R_b$ so that they break even given what they infer. Entrepreneur~2's choice of strategy is observed, so it works as a signal of the environment: copying entrepreneur~1 is read as ``favourable environment'', deviating as ``unfavourable''. Herding is an equilibrium when the better financing terms of the ``favourable'' pool outweigh the lower success probability of the copied strategy. Because investors break even, the entrepreneur bears the efficiency loss.

\notation
\begin{itemize}
  \item $R$: income on success; $I$: investment; no cash ($A=0$); interest rate 0.
  \item Unfavourable environment (prob.\ $1-\alpha$): success probabilities $(q,0)$ and $(0,q)$ for the two strategies; favourable (prob.\ $\alpha$): best strategies coincide with prob.\ $\theta$ (success $p$ for both, $r$ for the other strategy), differ with prob.\ $1-\theta$ ($(p,r)$ for entrepreneur~1's best, $(r,p)$ for entrepreneur~2's best). $p>\max\{q,r\}$, $qR>I$.
  \item $m\equiv\alpha p+(1-\alpha)q$; $R_b^1$: entrepreneur~1's reward; $R_b^F$: entrepreneur~2's reward when she copies (``favourable'' pool); $R_b^U$: her reward when she picks the other strategy (``unfavourable'' pool).
  \item $\pi_F\equiv\theta p+(1-\theta)r$: entrepreneur~2's success probability when she copies in the favourable environment.
\end{itemize}

\step{1}{Entrepreneur 1} He always picks his best strategy, so his contract reveals nothing about the environment; investors expect success probability $m$ and break even:
\begin{align}
  m\,(R-R_b^1) &= I \quad\Longrightarrow\quad R_b^1 = R - \frac{I}{m}. \label{r5e11-R1}
\end{align}
$R_b^1>0$ because $mR>qR>I$ ($m>q$ since $p>q$).

\step{2}{Entrepreneur 2's choices by state} In the unfavourable environment entrepreneur~1's project has success probability 0 for her, so she must pick the other strategy. In the favourable environment, under the candidate equilibrium, she always copies entrepreneur~1. Investors therefore read ``same strategy'' as favourable and ``different strategy'' as unfavourable.

\step{3}{Rewards in each pool} Same strategy: the pool is the favourable environment, in which copying succeeds with $p$ (prob.\ $\theta$) or $r$ (prob.\ $1-\theta$):
\begin{align}
  \pi_F\,(R-R_b^F) &= I \quad\Longrightarrow\quad R_b^F = R - \frac{I}{\theta p+(1-\theta)r}. \label{r5e11-RF}
\end{align}
Different strategy: the pool is the unfavourable environment, success probability $q$:
\begin{align}
  q\,(R-R_b^U) &= I \quad\Longrightarrow\quad R_b^U = R - \frac{I}{q}, \label{r5e11-RU}
\end{align}
and $R_b^U>0$ by $qR>I$, so the unfavourable type is funded.

\step{4}{Entrepreneur 2's incentive to herd} In the uncorrelated favourable state, copying gives success $r$ with reward $R_b^F$; using her own best strategy gives success $p$ but places her in the ``unfavourable'' pool, reward $R_b^U$. She herds iff
\begin{align}
  r\,R_b^F &\ge p\,R_b^U \nonumber\\
  r\Bigl[R - \frac{I}{\theta p+(1-\theta)r}\Bigr] &\ge p\Bigl[R - \frac{I}{q}\Bigr], \label{r5e11-herd}
\end{align}
substituting \eqref{r5e11-RF} and \eqref{r5e11-RU}: the stated condition. In the correlated favourable state copying gives success $p$ with $R_b^F$, against $r$ with $R_b^U$; since \eqref{r5e11-herd} with $r<p$ implies $R_b^F>R_b^U$, we get $pR_b^F>rR_b^U$ and she copies. In the unfavourable state she cannot copy profitably (success 0). Investors break even in each pool by construction, so beliefs are consistent and the candidate is an equilibrium.

\step{5}{Necessary implication} Because $r<p$, \eqref{r5e11-herd} needs $R_b^F>R_b^U$, that is,
\begin{align}
  \frac{I}{\theta p+(1-\theta)r} &< \frac{I}{q} \iff \theta p+(1-\theta)r > q: \nonumber
\end{align}
copying must be good enough news about the environment.

\step{6}{Entrepreneur 2 is worse off} With investors breaking even, her expected utility is the expected NPV of her project. Under herding,
\begin{align}
  U_2^{\text{herd}} &= (1-\alpha)\,q\,R_b^U + \alpha\,\pi_F\,R_b^F \nonumber\\
   &= (1-\alpha)(qR-I) + \alpha(\pi_F R - I) \nonumber\\
   &= \bigl[(1-\alpha)q + \alpha\theta p + \alpha(1-\theta)r\bigr]R - I. \label{r5e11-Uh}
\end{align}
If investors could not see entrepreneur~1's choice (or strategies were uncorrelated), she would always pick her best strategy and investors would price at $m$:
\begin{align}
  U_2^{\text{no obs}} &= mR - I = \bigl[\alpha p + (1-\alpha)q\bigr]R - I. \label{r5e11-Un}
\end{align}
Subtracting \eqref{r5e11-Uh} from \eqref{r5e11-Un}:
\begin{align}
  U_2^{\text{no obs}} - U_2^{\text{herd}} &= \alpha pR - \alpha\theta pR - \alpha(1-\theta)rR = \alpha(1-\theta)(p-r)R > 0. \label{r5e11-loss}
\end{align}

\step{7}{Numbers} Take $R=10$, $I=5$, $q=0.55$, $p=0.9$, $r=0.6$, $\theta=0.5$, $\alpha=0.5$ ($qR=5.5>5$, $p>\max\{q,r\}$). Then
\begin{align}
  m &= (0.5)(0.9)+(0.5)(0.55) = 0.725, \qquad R_b^1 = 10 - \frac{(5)}{(0.725)} = 3.1034, \nonumber\\
  \pi_F &= (0.5)(0.9)+(0.5)(0.6) = 0.75, \qquad R_b^F = 10 - \frac{(5)}{(0.75)} = 3.3333, \nonumber\\
  R_b^U &= 10 - \frac{(5)}{(0.55)} = 0.9091, \nonumber\\
  rR_b^F &= (0.6)(3.3333) = 2.0000 \;\ge\; pR_b^U = (0.9)(0.9091) = 0.8182, \nonumber
\end{align}
so herding holds. Utilities: $U_2^{\text{herd}} = (0.5)(0.55)(0.9091)+(0.5)(0.75)(3.3333) = 0.25+1.25 = 1.5$; $U_2^{\text{no obs}} = (0.725)(10)-5 = 2.25$; loss $=(0.5)(0.5)(0.9-0.6)(10) = 0.75 = 2.25-1.5$.
\answer{\unskip} Herding with funding is an equilibrium under \eqref{r5e11-herd}; it requires $\theta p+(1-\theta)r>q$, and it costs entrepreneur~2 $\alpha(1-\theta)(p-r)R$ in expected utility \eqref{r5e11-loss}.

\intuition Entrepreneur~2's choice is public, and in the unfavourable environment the two projects must differ. So ``doing something different'' is the hallmark of bad environments, and investors charge for it. Choosing the popular strategy buys the reputation of the favourable pool, and she pays for it by giving up her own better project. Since investors break even, the lost NPV, $\alpha(1-\theta)(p-r)R$, comes out of her pocket: herding is individually rational and collectively inefficient.

\pitfalls
\begin{itemize}
  \item Pricing entrepreneur~2's ``copy'' contract at $p$ instead of $\pi_F=\theta p+(1-\theta)r$: the pool includes the herders whose copied project only succeeds with $r$.
  \item Comparing $R_b^F$ with $R_b^U$ instead of the \emph{expected} rewards $rR_b^F$ and $pR_b^U$: the herder trades a lower success probability for a higher reward.
  \item Thinking investors lose from herding: they break even in every pool; the entrepreneur bears the whole cost.
\end{itemize}

\bookans Tirole p.~600: $m(R-R_b^1)=I$; $q(R-R_b^U)=I$ and $[\theta p+(1-\theta)r](R-R_b^F)=I$; herding requires $rR_b^F\ge pR_b^U$, which needs $\theta p+(1-\theta)r>q$ and holds e.g.\ ``if the project is hardly creditworthy in the unfavorable environment ($qR\simeq I$) and $r$ is not too small''. Matches the book's answer (p.~600); the book does not compute the welfare loss, which is \eqref{r5e11-loss}.
\refbib{Tirole \S6.2 ``Implications of the Lemons Problem and of Market Breakdown'' (privately-known-prospects model), pp.~241--243; lecture Topic~3, ``Pooling and cross-subsidies''.}

% =============================================================================
\exsec{12. Myers--Majluf with issuing costs --- Amaro de Matos 3.2.6 \#5, p.~95}
%% BEGIN-STATEMENT 12
\stmtbegin
\textbf{Source:} Amaro de Matos, Ch. 3, Exercise 3.2.6 \#5, p.~95 (PDF 113) · Answer in book: no
(p.~261 only refers to footnote 12, p.~203, of Myers and Majluf 1984)

Myers and Majluf (1984) assume that there are no issuing costs and prove that \emph{P}′ \textless{}
\emph{P}. What happens to the investment decision of the management and to the conclusion that
\(P' < P\) in the presence of issuing costs?

\paragraph{Context needed to solve it}

\emph{From Amaro de Matos §3.2, p.~78 (PDF 96):}

\textbf{Assumption 3.2.1.} \emph{Managers have better information about the value of the firm than
other agents in the economy.}

\emph{From Amaro de Matos §3.2.3 (Signals Conditioned by Investment Opportunities), pp.~86--88 (PDF
104--106):}

To see how all these conclusions are derived, the model must be formalized. Let \(I\) be the
required investment and \(S\) denote the \emph{slack} of the firm (i.e., large holdings of cash or
marketable securities) or the ability to issue default-risk-free debt.

If \(S \geq I\), it will not be necessary to issue any equity at all in order to finance the
investment. The only interesting case happens when \(S < I\), where the required amount of equity to
be issued is \(E = I - S\).

As opposed to the other models previously discussed, this is a three-date model with three distinct
times---say, time 0, time 1, and time 2. Let us introduce the following notation. At time 0, the
possible value of the assets of the firm at time 1 is described by a random variable
\(\widetilde{A}\). Similarly, let \(\widetilde{B}\) denote the possible net present value of the
investment opportunity at time 1.

At time 0, the information is symmetric. All the agents, managers and market, have the same
information---namely, the distributions of \(\widetilde{A}\) and \(\widetilde{B}\) and the value of
\(S\).

At time 1, the asymmetry of information between managers and outsider investors is specified.
Assumption 3.2.1 takes the form

\textbf{Assumption 3.2.5.} \emph{At time 1, the management knows the realizations of
\(\widetilde{A}\) and \(\widetilde{B}\), denoted respectively by \(a\) and \(b\), which are unknown
to the outsider investors. Negative values of \(a\) and \(b\) are ruled out.}

At this stage, the only thing that the outsiders know is that the amount of equity to be issued is
\(E = \max[0, I - S]\). Also, at time 1, let \(P\) be the market value of the shares of old
shareholders if the stock is not issued and \(P'\) be the price of the same shares if it is issued.
At time 2, finally, all the information is disclosed.

Let \(V^\circ\) be the value of the firm for the old shareholders. If the management, knowing the
values of \(a\) and \(b\), decides not to issue equity, then the company is still owned by the same
old shareholders with assets of value \(a\) and slack \(S\). Thus

\[V^\circ = a + S.\]

If, on the other hand, knowing the value of \(a\) and \(b\), the management decides to issue equity
\(E\), then the ownership of the company is shared between the old and new shareholders. The total
value of the company will include the issued equity \(E\) and the net present value of the
investment opportunity apart from the assets of value \(a\) and slack \(S\). But the old
shareholders will have only a proportion \(P'/(P' + E)\) of all of this. Thus, in this case

\[V^\circ = \frac{P'}{P' + E}(E + S + a + b).\]

To explain the decision of the managers, the following assumption is required:

\textbf{Assumption 3.2.6.} \emph{The management acts in the best interest of the old shareholders.
Moreover, old shareholders are assumed to be passive---that is, they do not adjust their portfolios
in response to the issue--invest decision of the firm.}

If that is so, the decision to issue new equity is taken only if the old shareholders are better off
in such case---that is, only if

\[\frac{P'}{P' + E}(E + S + a + b) \geq a + S.\]

This inequality may be rewritten as

\[b \geq -E + \frac{E}{P'}(S + a) \qquad (3.20)\]

and should naturally be restricted by Assumption 3.2.5 to the region \(a \geq 0\), \(b \geq 0\).

Hence, inequality (3.20) divides the domain of the joint probability distribution of
\(\widetilde{A}\) and \(\widetilde{B}\) into two disjoint regions identified in Figure 3.1: a region
\(M'\) in which all pairs \((a, b)\) satisfy the inequality and the firm issues new equity and
invests in the new project, and a region \(M\) in which no pair \((a, b)\) satisfies the inequality
and there is no investment. In the first region, the amount of equity issued is \(E = I - S\),
whereas in the second, it is zero.

\emph{Figure 3.1 (The Regions of Investment and No-Investment Decisions According to Myers and
Majluf's Model):} horizontal axis \(a\), vertical axis \(b\), both starting at 0. An upward-sloping
straight line labelled \(b = -E + \frac{E}{P'}(S + a)\) starts on the horizontal axis at a positive
value of \(a\) (where \(b = 0\)). The area above/left of the line is labelled ``region \(M'\)''; the
area below/right of the line is labelled ``region \(M\)''.

It is simple to see that the price \(P\) for the shares when \((a, b) \in M\) is greater than the
price \(P'\) when \((a, b) \in M'\). The equilibrium prices should be given by

\[P = S + E(\widetilde{A} \mid M) \qquad (3.21)\]

and

\[P' = S + E(\widetilde{A} + \widetilde{B} \mid M').\]

For any \((a, b) \in M\), the inequality (3.20) is not satisfied and hence

\[b < -E + \frac{E}{P'}(S + a) \quad \text{or} \quad a > P'(1 + b/E) - S.\]

Since \(b/E \geq 0\) from Assumption 3.2.5, it follows that for any \((a, b) \in M\),
\(a > P' - S\). Thus \(E(\widetilde{A} \mid M) > P' - S\) and (3.21) gives directly \(P > P'\).
\stmtend
%% END-STATEMENT

\asked (a) How the management's issue--invest rule changes when issuing equity costs the firm something. (b) Whether the result $P'<P$ (the shares are worth less when the firm issues than when it does not) survives.

\approach Redo the Myers--Majluf comparison of \S3.2.3 for old shareholders, adding an issue cost $c$ that new investors' money must cover. The algebra that gave (3.20) gives a new boundary of the investment region; the proof that $P>P'$ used one inequality, $a>P'-S$ on $M$, which we recheck. Finally a counterexample shows that the conclusion can fail. Concept: adverse selection in equity issues; issue costs make the issue decision depend on the project's NPV $b$, so the issue can become good news.

\notation
\begin{itemize}
  \item $S$: slack; $I$: investment; $E=I-S>0$: net amount the firm needs; $c\ge 0$: issue cost (paid to underwriters, so new investors pay $E+c$ for the firm to receive $E$).
  \item $a,b\ge0$: realized value of assets in place and NPV of the project, known to managers only; $M'$ ($M$): region of $(a,b)$ where the firm issues and invests (does not).
  \item $P'$: market value of the old shares when the firm issues; $P$: when it does not; $V^\circ$: intrinsic value of the old shareholders' stake.
\end{itemize}

\step{1}{Old shareholders' value} Without the issue, $V^\circ=a+S$. With the issue, new investors pay $E+c$ for a fair stake given the market's information; total market value after the issue is $P'+E+c$, so they own the fraction $(E+c)/(P'+E+c)$. The firm's true value is $E+S+a+b$ (the $c$ has left the firm). Hence
\begin{align}
  V^\circ_{\text{issue}} &= \frac{P'}{P'+E+c}\,(E+S+a+b). \label{r5e12-V}
\end{align}

\step{2}{(a) New investment rule} Management (acting for passive old shareholders, Assumption 3.2.6) issues iff \eqref{r5e12-V} is at least $a+S$:
\begin{align}
  \frac{P'}{P'+E+c}(E+S+a+b) &\ge a+S \nonumber\\
  P'(E+S+a+b) &\ge (a+S)(P'+E+c) \nonumber\\
  P'E + P'b &\ge (a+S)(E+c) \nonumber\\
  b &\ge -E + \frac{E+c}{P'}\,(S+a). \label{r5e12-rule}
\end{align}
At $c=0$ this is (3.20). For a given $P'$ the boundary moves up by $c(S+a)/P'$ and becomes steeper (slope $(E+c)/P'$ instead of $E/P'$); it crosses $b=0$ at $a=P'E/(E+c)-S<P'-S$. So the no-issue region $M$ grows: more positive-NPV projects are passed up. Two forces are now mixed: with symmetric information ($P'=S+a+b-c$), \eqref{r5e12-rule} reduces to
\begin{align}
  (S+a+b-c)(E+b) &\ge (S+a)(E+c) \nonumber\\
  (S+a)b + (E+b)(b-c) - (S+a)c &\ge 0 \nonumber\\
  (S+a+E+b)(b-c) &\ge 0 \iff b \ge c, \nonumber
\end{align}
the efficient rule net of costs; the information problem adds underinvestment on top of it.
\answer{(a)} Issue and invest iff $b\ge (E+c)(S+a)/P' - E$: for a given price the investment region shrinks, and even under symmetric information only projects with $b\ge c$ are taken.

\step{3}{(b) Where the proof of $P>P'$ breaks} For $(a,b)\in M$, \eqref{r5e12-rule} fails:
\begin{align}
  (S+a)(E+c) &> P'(E+b) \;\ge\; P'E \quad (b\ge0), \nonumber\\
  a &> \frac{P'E}{E+c} - S. \nonumber
\end{align}
Taking expectations over $M$ and using $P=S+E(\widetilde A\mid M)$, eq.~(3.21),
\begin{align}
  P &> \frac{E}{E+c}\,P' \iff P' < \frac{E+c}{E}\,P. \label{r5e12-bound}
\end{align}
At $c=0$ this is $P'<P$; with $c>0$ it only bounds $P'$ by $P(E+c)/E>P$, so $P'>P$ is not ruled out.

\step{4}{A counterexample with $P'>P$} Let assets in place be known, $S+a=100$, so only $b$ is private, with $b\sim U[0,10]$; $E=10$, $c=1$. By \eqref{r5e12-rule} the firm issues iff $b\ge b^*$ with
\begin{align}
  b^* &= \frac{(E+c)(S+a)}{P'} - E = \frac{(11)(100)}{P'} - 10. \nonumber
\end{align}
The market prices the old shares at the issuers' average: $P' = (S+a) + E(b\mid b\ge b^*) - c = 100 + \tfrac{b^*+10}{2} - 1 = 104 + \tfrac{b^*}{2}$. Substituting into $b^*$:
\begin{align}
  (b^*+10)\Bigl(104+\frac{b^*}{2}\Bigr) &= 1100 \nonumber\\
  \frac{(b^*)^2}{2} + 109\,b^* + 1040 &= 1100 \nonumber\\
  (b^*)^2 + 218\,b^* - 120 &= 0 \nonumber\\
  b^* &= \frac{-218+\sqrt{(218)^2+4(120)}}{2} = \frac{-218+\sqrt{48{,}004}}{2} = 0.5491. \nonumber
\end{align}
Then $P' = 104 + (0.5491)/2 = 104.27$, while firms that do not issue are worth $P=S+a=100$. So $P'>P$: the issue is good news. The bound \eqref{r5e12-bound} holds, $104.27<(11/10)(100)=110$. Without the cost ($c=0$) every firm with $b\ge0$ would issue ($b^*=1000/105-10<0$), $M$ would be empty and the comparison would be void.
\answer{(b)} The conclusion $P'<P$ no longer follows: with issue costs only firms with a good enough project issue, the issue reveals information about $b$ as well as about $a$, and $P'$ can exceed $P$ (above, $104.27>100$).

\intuition Without costs, the only reason not to issue is that assets in place are undervalued, so silence is good news and an issue is bad news. Issue costs add a second reason not to issue, a poor project, so silence now also means ``no good project'' and an issue also means ``good project''. Which effect dominates depends on how much of the uncertainty is about $a$ and how much about $b$. Costs also make internal funds (slack) more valuable, strengthening the pecking order.

\pitfalls
\begin{itemize}
  \item Letting the issue cost reduce the firm's value but not the new investors' payment: new money is $E+c$, of which only $E$ reaches the project.
  \item Treating \eqref{r5e12-rule} as a rule with $P'$ fixed: $P'$ is itself an equilibrium object that depends on who issues (Step~4 solves the fixed point).
  \item Concluding that issue costs make $P'>P$ always: they only remove the proof; with mostly $a$-uncertainty $P'<P$ still holds.
\end{itemize}

\bookans Amaro de Matos App.~C, p.~261, only says ``See footnote 12, page 203, of Myers and Majluf (1984)'' (not available locally); no worked answer to compare with.
\refbib{Amaro de Matos \S3.2.3 ``Signals Conditioned by Investment Opportunities'', pp.~86--88, eqs.~(3.20)--(3.21); lecture Topic~3, ``Myers--Majluf and the pecking order''.}

%% BEGIN-FROM lista-topic3-asymmetric-information-solutions.tex EX 15 AS 13
\exsec{13. Assets in place and the negative stock-price reaction with a continuum of types --- Tirole 6.5, p.~275}
%% BEGIN-STATEMENT 13
\stmtbegin
\textbf{Source:} Tirole, Ch. 6, Exercise 6.5, p.~275 (PDF 288) · Answer in book: yes, p.~597

Consider the privately-known-prospects model of Application 2 in Section 6.2.2, but with a continuum
of types. The entrepreneur already owns a project, which with probability \(p\) yields profit \(R\)
and probability \(1 - p\) profit 0. The probability \(p\) is private information of the borrower.
From the point of view of the investors, \(p\) is drawn from cumulative distribution \(F(p)\) with
continuous density \(f(p) > 0\) on some interval \([\underline{p}, \bar{p}]\). Assume that the
distribution has monotone hazard rates:

\[\frac{f(p)}{F(p)} \quad \text{is decreasing in } p\]

and

\[\frac{f(p)}{1 - F(p)} \quad \text{is increasing in } p.\]

(This assumption, which is satisfied by most usual distributions, is known to imply that the
truncated means \(m^-(p)\) and \(m^+(p)\) have slope less than 1:

\[0 < (m^-(p))' \equiv \frac{d}{dp}[E(\tilde{p} \mid \tilde{p} \leqslant p)] \leqslant 1\]

and

\[0 < (m^+(p))' \equiv \frac{d}{dp}[E(\tilde{p} \mid \tilde{p} \geqslant p)] \leqslant 1\]

(see, for example, An 1998).)

The model is otherwise as in Section 6.2.2. A seasoned offering may be motivated by a profitable
deepening investment: at cost \(I\), the probability of success can be raised by an amount \(\tau\)
such that

\[\tau R > I\]

(of course, we need to assume that \(\bar{p} + \tau \leqslant 1\)). The entrepreneur has no cash on
hand, is risk neutral, and is protected by limited liability. The investors are risk neutral and
demand a rate of return equal to 0.

\begin{enumerate}
\def\labelenumi{(\roman{enumi})}
\item
  Show that in any equilibrium, only types \(p \leqslant p^*\), for some cutoff \(p^*\), raise funds
  and finance the deepening investment.
\item
  Show that \(p^* > \underline{p}\) and that if \(p^* < \bar{p}\), then
\end{enumerate}

\[\frac{\tau R}{I} = \frac{p^* + \tau}{m^-(p^*) + \tau}.\]

Show that if the benefits from investment are ``not too large,'' in that

\[\frac{\tau R}{I} < \frac{\bar{p} + \tau}{E[p] + \tau},\]

then indeed \(p^* < \bar{p}\).

Show that if there are multiple equilibria, the one with the highest cutoff \(p^*\) Pareto-dominates
(is better for all types than) the other equilibria.

\begin{enumerate}
\def\labelenumi{(\roman{enumi})}
\setcounter{enumi}{2}
\item
  Is there a negative stock price reaction upon announcement of an equity issue?
\item
  Focusing on an interior Pareto-dominant equilibrium, show that, when \(\tau\) increases, the
  volume of equity issues increases.
\end{enumerate}

\paragraph{Context needed to solve it}

\emph{From Tirole §6.2 (privately-known-prospects model), p.~241 (PDF 254):}

\emph{Privately-known-prospects model.} A borrower/entrepreneur has no funds (\(A = 0\)) to finance
a project costing \(I\). The project yields \(R\) in the case of success and 0 in the case of
failure. The borrower and the lenders are risk neutral, and the borrower is protected by limited
liability. The interest rate in the economy is normalized at 0.

The borrower can be one of two types. A good borrower has a probability of success equal to \(p\). A
bad borrower has a probability of success \(q\). Assume that \(p > q\) and that \(pR > I\) (at least
the good type is creditworthy). {[}\ldots{]}

The borrower has private information about her type. The capital market, which is competitive and
demands an expected rate of return equal to 0, puts probabilities \(\alpha\) and \(1 - \alpha\) on
the borrower being a good or a bad type, respectively. Under asymmetric information, the capital
market does not know whether it faces a ``\(p\)-borrower'' (a good borrower) or a ``\(q\)-borrower''
(a bad borrower). Let

\[m \equiv \alpha p + (1 - \alpha) q\]

denote the investors' prior probability of success.

\emph{From Tirole §6.2.2, Application 2: Assets in Place, the Negative Stock Price Reaction, and the
Going-Public Decision, pp.~244--245 (PDF 257--258):}

Let us next suppose that the entrepreneur already owns a project that, without further investment,
will succeed with probability \(p\) or \(q\), yielding profit \(R\). As before, the entrepreneur
knows the probability of success while the investors put probability \(\alpha\) on \(p\) and
\((1 - \alpha)\) on \(q\). Thus, in the absence of further information (and so the investors'
expectation of the probability of success is \(m\)), the assets in place are undervalued
(respectively, overvalued) if the true probability of success is \(p\) (respectively, \(q\)).

For computational simplicity, we will assume that the entrepreneur initially owns all shares. But
nothing is altered if she owns only a fraction of the shares. By ``stock price reaction upon the
announcement of an equity issue,'' we mean the difference between the total value of shares (whoever
owns them) before and after the announcement. This notion corresponds to the approach taken by event
studies in empirical work.

An equity offering may be motivated by a profitable ``deepening investment'' (more generally, the
key feature is that one cannot contract on the cash flow generated by this investment separately
from that generated by assets in place: the incomes generated by the two are intertwined or
fungible\textsuperscript{22}). At cost \(I\), the probability of success can be raised by an amount
\(\tau\) such that

\[\tau R > I.\]

That is, investing is efficient for both types of borrowers. Note that we assume for the moment that
the increase in profitability is uniform across types: the probability of success becomes
\(p + \tau\) for a good borrower and \(q + \tau\) for a bad one.

The entrepreneur, however, has no cash on hand. Accordingly, the full amount \(I\) must be raised
from investors. The entrepreneur must therefore issue new shares, thereby reducing the fraction of
shares she owns.

A key insight is that relinquishing shares to investors is relatively less costly to the borrower
with overvalued assets in place (the bad borrower) than to the borrower with undervalued assets in
place (the good borrower). Thus, if the good borrower conducts an equity offering, so does the bad
one.

Let us therefore investigate the possibility of an (efficient) pooling equilibrium. The entrepreneur
must offer a stake \(R_l\) in success to the investors such that

\[[\alpha(p + \tau) + (1 - \alpha)(q + \tau)]R_l = I \iff (m + \tau)R_l = I,\]

where, as earlier, \(m \equiv \alpha p + (1 - \alpha)q\) is the prior mean probability of success.
There exists a unique \(R_l\), \(0 < R_l < R\), satisfying this condition.

The good borrower, though, can guarantee herself \(pR\) by not diluting her
stake.\textsuperscript{23} Thus, she is willing to issue new shares only if

\[(p + \tau)(R - R_l) \geqslant pR \iff \tau R \geqslant \frac{p + \tau}{m + \tau} I. \qquad (6.1)\]

{[}\ldots{]} We are thus led to consider two situations:

\emph{Pooling equilibrium.} If condition (6.1) holds, then both types conduct an equity
offering.\textsuperscript{24} If the accrual of this deepening investment is
anticipated,\textsuperscript{25} the total value of shares \emph{before and after} the seasoned
equity offering is

\[(m + \tau)R - I.\]

There is no stock price reaction to the offering, which is perfectly anticipated and uninformative.

\emph{Separating equilibrium.} More interestingly, suppose that condition (6.1) is violated. The
good borrower then does not raise funds. The bad borrower still does, but under market conditions
that are not as favorable as in a pooling equilibrium. Because the investors know that the equity
offering reveals overvalued assets, they demand a higher stake \(R_l^B > R_l\) such that

\[(q + \tau)R_l^B = I.\]

The good borrower does not want to raise funds because

\[(p + \tau)(R - R_l^B) < pR \iff \tau R < \frac{p + \tau}{q + \tau} I, \qquad (6.2)\]

which holds if condition (6.1) is violated.

The announcement of a seasoned equity offering then leads to a \emph{negative stock price reaction}.
The preannouncement total value of shares is\textsuperscript{26}

\[V_0 = \alpha[pR] + (1 - \alpha)[(q + \tau)R - I].\]

After the announcement, it becomes

\[V_1 = (q + \tau)R - I.\]

Hence,

\[V_0 > V_1 \iff pR > (q + \tau)R - I.\]

But we know that

\[pR > (p + \tau)\left(R - \frac{I}{q + \tau}\right),\]

and so \emph{a fortiori} \(V_0 > V_1\).

\begin{quote}
\begin{enumerate}
\def\labelenumi{\arabic{enumi}.}
\setcounter{enumi}{21}
\item
  Otherwise, it would be optimal for the good borrower to engage in project finance so as to avoid
  having to cross-subsidize the bad one.
\item
  That this ``reservation utility'' depends on the borrower's type is the essential difference with
  the barebones model. Here, in the jargon of incentive theory, ``reservation utilities are
  type-contingent.'' {[}\ldots{]}
\item
  This pooling equilibrium is not unique whenever
  \((p + \tau)I/(m + \tau) \leqslant \tau R \leqslant (p + \tau)I/(q + \tau)\); indeed, if investors
  believe that an equity offering comes from a bad borrower and
  \(\tau R \leqslant (p + \tau)I/(q + \tau)\), then the good type indeed prefers not to raise funds.
  However, the pooling equilibrium is the \emph{Pareto-dominant} equilibrium (it is the best
  equilibrium for both the good and the bad borrower), and so we will focus on it. Note also that if
  condition (6.1) holds, then the bad borrower definitely prefers to raise funds since the analogous
  condition for her is \(\tau R \geqslant \frac{q + \tau}{m + \tau} I\), which is always satisfied.
\item
  Otherwise, the news of the existence of an investment opportunity by itself raises the value of
  shares.
\item
  As in the previous footnote, note that we assume that the investment opportunity is perfectly
  anticipated by the capital market. Otherwise, the issue of new securities could convey good news
  about the firm's opportunity set and the concomitant boost in share price might dominate the
  effect unveiled here.
\end{enumerate}
\end{quote}
\stmtend
%% END-STATEMENT

\asked (i) Show that in any equilibrium the firms that issue equity to fund the deepening investment are exactly the types below some cutoff $p^*$. (ii) Show that $p^*>\underline p$; that an interior cutoff satisfies $\tau R/I=(p^*+\tau)/(m^-(p^*)+\tau)$; that $p^*<\bar p$ when the gains from investing are not too large; and that the equilibrium with the highest cutoff is best for every type. (iii) Say whether an issue announcement lowers the stock price. (iv) At an interior Pareto-dominant equilibrium, show that the volume of equity issues rises with $\tau$.

\approach This is Tirole's Application 2 (assets in place and the negative stock-price reaction, \S6.2.2; Myers--Majluf in lecture block 3.5) with a continuum of types. An issuer pools with all other issuers at a single stake $R_l$, so issuing dilutes her assets in place at the pooled price; the dilution costs more the better her assets ($p$), which yields a cutoff rule. Investors' break-even over the issuing interval $[\underline p,p^*]$ gives $R_l$; indifference of the cutoff type gives the equilibrium condition. The price reaction compares total share value before and after the market learns $p\le p^*$; the comparative statics in $\tau$ use the implicit function theorem.

\notation
\begin{itemize}
  \item $p\in[\underline p,\bar p]$: privately known success probability of assets in place, c.d.f.\ $F$, density $f>0$; $E[p]$: prior mean.
  \item $R$: profit in success, 0 in failure; $I$: cost of the deepening investment, which raises $p$ to $p+\tau$, $\tau R>I$, $\bar p+\tau\le1$.
  \item $R_l$: investors' stake in success (they get 0 in failure); $p^*$: issue cutoff.
  \item $m^-(p)=E[\tilde p\mid\tilde p\le p]$, $m^+(p)=E[\tilde p\mid\tilde p\ge p]$; $0<(m^\pm)'\le1$.
  \item $V_0$, $V_1$: total value of shares before / after the issue announcement; $H(p^*,\tau)$: equilibrium function of (iv).
\end{itemize}

\step{1}{Issuers form a lower interval (part (i))} Not issuing gives the entrepreneur $pR$. Issuing for a
stake $R_l$ gives $(p+\tau)(R-R_l)$, so type $p$ issues iff
\begin{align}
  (p+\tau)(R-R_l) &\ge pR \nonumber\\
  pR+\tau R-(p+\tau)R_l &\ge pR \nonumber\\
  \tau R &\ge (p+\tau)R_l. \label{t3e15-issue}
\end{align}
In any equilibrium all issuers use the same stake: every type's payoff falls with $R_l$, so if two stakes
were accepted every issuer would pick the smaller. The left side of \eqref{t3e15-issue} does not depend on $p$ and the right side rises with $p$ (slope $R_l>0$),
so if type $p$ issues, every type below $p$ issues too: the issuers are the types $p\le p^*$ for some cutoff $p^*\in[\underline p,\bar p]$. $\blacksquare$

\answer{(i)} Type $p$ issues iff $\tau R\ge(p+\tau)R_l$; the cost of dilution $(p+\tau)R_l$ rises with $p$, so issuers are the types $p\le p^*$.

\step{2}{Break-even stake and the interior cutoff (part (ii))} Investors break even on the issuers, whose average success probability after investing is $E(\tilde p+\tau\mid\tilde p\le p^*)=m^-(p^*)+\tau$:
\begin{align}
  \bigl[E(\tilde p\mid\tilde p\le p^*)+\tau\bigr]R_l &= I \nonumber\\
  R_l &= \frac{I}{m^-(p^*)+\tau}. \label{t3e15-Rl}
\end{align}
If $p^*<\bar p$, the cutoff type is indifferent in \eqref{t3e15-issue} (types just below issue, types just above do not, and payoffs are continuous in $p$). Substituting \eqref{t3e15-Rl} and dividing by $I$,
\begin{align}
  \tau R &= (p^*+\tau)\frac{I}{m^-(p^*)+\tau} \nonumber\\
  \frac{\tau R}{I} &= \frac{p^*+\tau}{m^-(p^*)+\tau}. \label{t3e15-cut}
\end{align}

\step{3}{$p^*>\underline p$} Whatever the investors believe, they put mass only on $[\underline p,\bar p]$,
so the stake they require is at most $I/(\underline p+\tau)$. Under that stake, \eqref{t3e15-issue} holds
for type $\underline p$ with strict inequality:
\begin{align}
  (\underline p+\tau)\frac{I}{\underline p+\tau} = I < \tau R. \nonumber
\end{align}
By continuity it also holds strictly for a positive-measure set of types above $\underline p$. So the
cutoff cannot be $\underline p$ (equivalently, $m^-(\underline p)=\underline p$ makes the right side of
\eqref{t3e15-cut} equal 1, below $\tau R/I$).

\step{4}{$p^*<\bar p$ when benefits are not too large} If $p^*=\bar p$ (everyone issues), then
$m^-(\bar p)=E[p]$ and \eqref{t3e15-Rl} gives $R_l=I/(E[p]+\tau)$. Type $\bar p$ would then issue only if
\begin{align}
  \tau R \ge (\bar p+\tau)\frac{I}{E[p]+\tau}
  \iff \frac{\tau R}{I} \ge \frac{\bar p+\tau}{E[p]+\tau}, \nonumber
\end{align}
which the assumed condition $\tau R/I<(\bar p+\tau)/(E[p]+\tau)$ rules out. Types near $\bar p$ therefore
strictly prefer not to issue, a contradiction, so $p^*<\bar p$. An interior solution exists:
$\varphi(p)\equiv(p+\tau)/(m^-(p)+\tau)$ is continuous, $\varphi(\underline p)=1<\tau R/I$ and
$\varphi(\bar p)>\tau R/I$, so by the intermediate value theorem \eqref{t3e15-cut} has a root in
$(\underline p,\bar p)$.

\step{5}{The highest cutoff Pareto-dominates} Take equilibria with cutoffs $p_1^*<p_2^*$. Since
$(m^-)'>0$, \eqref{t3e15-Rl} gives $R_{l2}<R_{l1}$. Types $p\le p_1^*$ issue in both and pay less in the
second. Types $p\in(p_1^*,p_2^*]$ get $pR$ in the first and choose to issue in the second, so by
\eqref{t3e15-issue} they get $(p+\tau)(R-R_{l2})\ge pR$. Types above $p_2^*$ get $pR$ in both. Investors
break even in both. Hence the highest-cutoff equilibrium is weakly better for every type. $\blacksquare$
(The condition $(m^-)'\le1$ does not by itself make the root unique; for a uniform $F$, $(m^-)'=\frac12$
and the root is unique.)

\answer{(ii)} $p^*>\underline p$ (the lowest type always gains from issuing); an interior cutoff solves $\tau R/I=(p^*+\tau)/(m^-(p^*)+\tau)$ with $R_l=I/(m^-(p^*)+\tau)$; $\tau R/I<(\bar p+\tau)/(E[p]+\tau)$ rules out $p^*=\bar p$; and a higher cutoff means a lower $R_l$, so the highest-cutoff equilibrium Pareto-dominates.

\step{6}{Negative stock-price reaction (part (iii))} The opportunity is anticipated. Before the
announcement, total share value is expected value net of the funds investors will contribute (they break
even): an issuer ($p\le p^*$, probability $F(p^*)$) is worth $(p+\tau)R-I$, a non-issuer $pR$. Averaging, and using $E[p]=F(p^*)m^-(p^*)+[1-F(p^*)]m^+(p^*)$,
\begin{align}
  V_0 &= F(p^*)\bigl[(m^-(p^*)+\tau)R-I\bigr]+[1-F(p^*)]m^+(p^*)R \label{t3e15-V0}\\
      &= \bigl\{F(p^*)m^-(p^*)+[1-F(p^*)]m^+(p^*)\bigr\}R+F(p^*)(\tau R-I) \nonumber\\
      &= E[p]R+F(p^*)(\tau R-I). \nonumber
\end{align}
After the announcement the market knows $p\le p^*$:
\begin{align}
  V_1 = [m^-(p^*)+\tau]R-I. \label{t3e15-V1}
\end{align}
Subtracting \eqref{t3e15-V1} from \eqref{t3e15-V0},
\begin{align}
  V_0-V_1 = [1-F(p^*)]\Bigl\{m^+(p^*)R-\bigl[(m^-(p^*)+\tau)R-I\bigr]\Bigr\}. \label{t3e15-dV}
\end{align}
At an interior equilibrium, \eqref{t3e15-cut} gives $I=\tau R(m^-+\tau)/(p^*+\tau)$, so
\begin{align}
  (m^-+\tau)R-I &= (m^-+\tau)R\Bigl[1-\frac{\tau}{p^*+\tau}\Bigr] = (m^-+\tau)R\,\frac{p^*}{p^*+\tau}, \nonumber\\
  V_0-V_1 &= [1-F(p^*)]R\Bigl[m^+(p^*)-p^*\,\frac{m^-(p^*)+\tau}{p^*+\tau}\Bigr]. \label{t3e15-dV2}
\end{align}
Since $m^+(p^*)>p^*$ and $(m^-(p^*)+\tau)/(p^*+\tau)<1$, the bracket is positive and $1-F(p^*)>0$: the
announcement of an equity issue lowers the stock price. If everyone issues ($p^*=\bar p$, pooling), the
factor $1-F(p^*)=0$ and there is no reaction.

\answer{(iii)} Yes, at an interior equilibrium: $V_0-V_1=[1-F(p^*)]R\,[m^+(p^*)-p^*(m^-(p^*)+\tau)/(p^*+\tau)]>0$, because the announcement tells the market that $p\le p^*$. With full pooling there is no reaction.

\step{7}{Comparative statics in $\tau$ (part (iv))} Define
\begin{align}
  H(p^*,\tau) \equiv \frac{\tau R}{I}\bigl[m^-(p^*)+\tau\bigr]-(p^*+\tau), \label{t3e15-H}
\end{align}
so that \eqref{t3e15-cut}, multiplied by $m^-(p^*)+\tau$ and rearranged, is $H=0$. Then $H(\underline p,\tau)=(\tau R/I)(\underline p+\tau)-(\underline p+\tau)=(\tau R/I-1)(\underline p+\tau)>0$ (using $m^-(\underline p)=\underline p$) and
$H(\bar p,\tau)<0$ under the condition of (ii). At the highest root (the Pareto-dominant one) $H$ crosses
from positive to negative, so $H_{p^*}=(\tau R/I)(m^-)'-1<0$ (taking the crossing to be transversal).
From $H=0$, $(\tau R/I)(m^-+\tau)=p^*+\tau$, so
\begin{align}
  \frac{\tau R}{I}-1 = \frac{p^*+\tau}{m^-+\tau}-1 = \frac{p^*-m^-}{m^-+\tau}. \nonumber
\end{align}
Differentiating \eqref{t3e15-H} in $\tau$ (product rule on $(\tau R/I)(m^-+\tau)$) and substituting this,
\begin{align}
  H_\tau &= \frac RI(m^-+\tau)+\frac{\tau R}{I}-1 \nonumber\\
         &= \frac RI(m^-+\tau)+\frac{p^*-m^-(p^*)}{m^-(p^*)+\tau} > 0. \label{t3e15-Ht}
\end{align}
By the implicit function theorem,
\begin{align}
  \frac{dp^*}{d\tau} = -\frac{H_\tau}{H_{p^*}} > 0, \label{t3e15-dp}
\end{align}
and the volume of equity issues, $F(p^*)I$ (the issuers are the types below $p^*$), rises with $\tau$, since
$d[F(p^*)I]/d\tau=If(p^*)\,dp^*/d\tau>0$. $\blacksquare$ (The sign of $H_\tau$: $m^-+\tau>0$ and $p^*>m^-(p^*)$, since a truncated mean lies below the truncation point.)

\answer{(iv)} $dp^*/d\tau=-H_\tau/H_{p^*}>0$, so the issue volume $F(p^*)I$ rises with $\tau$ at rate $If(p^*)\,dp^*/d\tau>0$.

\step{8}{Numerical illustration (uniform $F$)} Take $p\sim U[0.1,0.5]$, $\tau=0.3$, $R=10$, $I=2.5$:
$\tau R=3>I$, $E[p]=0.3$, and $\tau R/I=1.2<(0.8)/(0.6)=1.333333$. With $m^-(p)=(p+0.1)/2$,
\eqref{t3e15-cut} reads $1.2[(p^*+0.1)/2+0.3]=p^*+0.3$, so $0.6p^*+0.06+0.36=p^*+0.3$, i.e.\ $0.42-0.3=0.4p^*$ and $p^*=0.12/0.4=0.3$. Then
$R_l=2.5/(0.2+0.3)=5$, and the cutoff type is indifferent: $(0.6)(5)=3=(0.3)(10)$. Also $F(p^*)=0.5$,
$m^+(p^*)=0.4$, $V_0=3+0.5(0.5)=3.25$, $V_1=(0.5)(10)-2.5=2.5$, and $V_0-V_1=0.75$, which matches
\eqref{t3e15-dV2}: $(0.5)(10)[0.4-0.3(0.5)/0.6]=0.75$. For (iv), $H_{p^*}=1.2(0.5)-1=-0.4$ and
$H_\tau=4(0.5)+0.1/0.5=2.2$, so $dp^*/d\tau=5.5$ and the issue volume $F(p^*)I=1.25$ rises at rate
$2.5\times2.5\times5.5=34.375$.

\intuition Issuing equity dilutes the entrepreneur's assets in place at the pooled price. Types whose
assets are worth more (high $p$) lose more from that, so only types below a cutoff issue, and an issue
announcement tells the market the firm is in the lower part of the distribution: the price drops.
Equilibria with higher cutoffs carry less stigma (a lower $R_l$) and are better for everyone. A more
valuable investment ($\tau$ up) makes higher types willing to accept the dilution, so the cutoff and
the volume of issues rise.

\pitfalls
\begin{itemize}
  \item Pricing the issue at the prior mean $E[p]$: investors know that only types $p\le p^*$ issue, so the stake uses $m^-(p^*)$.
  \item Measuring the price reaction on the entrepreneur's stake only: as in Tirole, the reaction is the change in the \emph{total} value of shares, net of the funds investors contribute.
  \item Writing the issue volume as $[1-F(p^*)]I$ (the book's typo): the issuers are the types \emph{below} the cutoff, so the volume is $F(p^*)I$.
\end{itemize}

\bookans Matches the book's answer (pp.~597--598): cutoff rule, $R_l=I/(m^-(p^*)+\tau)$, the interior
condition, $p^*>\underline p$, the Pareto ranking, $V_0-V_1>0$ and $dp^*/d\tau>0$. Two typos in the
book: its stock-price item is labelled ``(ii)'' instead of (iii), and it gives the volume of issues as
``$[1-F(p^*)]I$'', whereas the issuers are the types $p\le p^*$, so the volume is $F(p^*)I$. The conclusion
(volume rises with $\tau$) is unaffected.

\refbib{Tirole Ch.~6, \S6.2.2 Application 2 ``Assets in Place, the Negative Stock Price Reaction, and the
Going-Public Decision'', pp.~244--245; Exercise 6.5, p.~275; answer pp.~597--598. Lecture Topic~3,
pecking-order block.}
%% END-FROM

%% BEGIN-FROM lista-topic3-asymmetric-information-solutions.tex EX 12 AS 14
\exsec{14. Pecking order with variable investment --- Tirole 6.8, p.~278}
%% BEGIN-STATEMENT 14
\stmtbegin
\textbf{Source:} Tirole, Ch. 6, Exercise 6.8, p.~278 (PDF 291) · Answer in book: yes, p.~599

Consider the privately-known-prospects model with risk neutrality and variable investment. For
investment \(I\), the realized income is either \(R^S I\) (in the case of success) or \(R^F I\) (in
the case of failure), where \(R^S > R^F \geqslant 0\). A good borrower has probability \(p_H\) of
success when working and \(p_L\) when shirking; similarly, a bad borrower has probability \(q_H\) of
success when working and \(q_L\) when shirking, where \(p_H - p_L = \Delta p = q_H - q_L\), for
simplicity. The entrepreneur's private benefit is 0 when working and \(BI\) when shirking. The
entrepreneur is risk neutral and protected by limited liability; the investors are risk neutral and
demand a rate of return equal to 0.

\begin{enumerate}
\def\labelenumi{(\roman{enumi})}
\item
  Let \(\tilde{U}_b^{SI}\) denote the bad borrower's gross utility under symmetric
  information.\textsuperscript{80} Consider the problem of maximizing the good borrower's utility
  subject to the investors' breaking even on that borrower, to the mimicking constraint that the
  good borrower's terms not be preferred by the bad borrower to her symmetric-information terms, and
  to the no-shirking constraint. Let \(\{R_b^S, R_b^F\}\) denote the (nonnegative) rewards of the
  good borrower in the cases of success and failure. Write the separating program.
\item
  Show that \(R_b^F = 0\).
\item
  (Only if you have read the supplementary section.) Show that the separating outcome is the only
  perfect Bayesian equilibrium of the issuance game if and only if \(\alpha \leqslant \alpha^*\) for
  some threshold \(\alpha^*\).
\end{enumerate}

\begin{quote}
\begin{enumerate}
\def\labelenumi{\arabic{enumi}.}
\setcounter{enumi}{79}
\tightlist
\item
  This utility was derived in Section 3.4.2. It is equal to
  \[\left[1 + \frac{q_H R - 1}{1 - q_H(R - B/\Delta p)}\right]A\] if \(q_H R \geqslant 1\), and to
  \(A\) otherwise.
\end{enumerate}
\end{quote}

\paragraph{Context needed to solve it}

\emph{From Tirole §6.2 (privately-known-prospects model), p.~241 (PDF 254):}

\emph{Privately-known-prospects model.} A borrower/entrepreneur has no funds (\(A = 0\)) to finance
a project costing \(I\). The project yields \(R\) in the case of success and 0 in the case of
failure. The borrower and the lenders are risk neutral, and the borrower is protected by limited
liability. The interest rate in the economy is normalized at 0.

The borrower can be one of two types. A good borrower has a probability of success equal to \(p\). A
bad borrower has a probability of success \(q\). {[}\ldots{]}

The borrower has private information about her type. The capital market, which is competitive and
demands an expected rate of return equal to 0, puts probabilities \(\alpha\) and \(1 - \alpha\) on
the borrower being a good or a bad type, respectively. Under asymmetric information, the capital
market does not know whether it faces a ``\(p\)-borrower'' (a good borrower) or a ``\(q\)-borrower''
(a bad borrower). {[}\ldots{]}

\emph{From Tirole §3.4.1--3.4.2 (the continuous-investment model, where footnote 80's formula is
derived), pp.~127--128 (PDF 140--141):}

\textbf{3.4.1 The Continuous-Investment Model.} The continuous-investment model of this section is
the polar opposite of the fixed-investment one. The fixed-investment model depicts a situation in
which returns are sharply decreasing beyond some investment level. In contrast, we now assume that
there are constant returns to scale in the investment technology. An investment
\(I \in [0, \infty)\) yields income \(RI\), proportional to \(I\), in the case of success, and 0 in
the case of failure. The borrower's private benefit from misbehaving is also taken to be
proportional to investment. As before, the borrower has a choice between behaving, in which case she
derives no private benefit and the probability of success is \(p_H\), and misbehaving, that is,
enjoying private benefit \(BI\), and reducing the probability of success to
\(p_L = p_H - \Delta p < p_H\). {[}\ldots{]}

The borrower initially has cash \(A\), and must therefore borrow \(I - A\) to finance a project of
size \(I\). A loan agreement specifies that the lenders (who as before are assumed to make no
profit) and the borrower receive 0 each in the case of failure, and \(R_l\) and \(R_b\),
respectively, in the case of success, where \(R_l + R_b = RI\).

As in Section 3.2, we assume that investment has positive NPV (net present value), here per unit of
investment, if the borrower behaves,

\[p_H R > 1, \qquad (3.7)\]

but negative NPV otherwise,

\[1 > p_L R + B, \qquad (3.8)\]

so that unless one can control the agency problem the investment cannot be funded. We also make an
assumption that guarantees that the equilibrium investment is finite:

\[p_H R < 1 + \frac{p_H B}{\Delta p}. \qquad (3.9)\]

{[}\ldots{]} Finally, we keep assuming that the capital market is competitive. {[}\ldots{]}

\textbf{3.4.2 The Lenders' Credit Analysis.} Following the steps of Section 3.2, the incentive
compatibility and the breakeven conditions are

\[(\Delta p) R_b \geqslant BI \qquad (\text{IC}_b)\]

and

\[p_H(RI - R_b) \geqslant I - A. \qquad (\text{IR}_l)\]

In equilibrium, competitive lenders make no profit on the contract that is most advantageous for the
borrower; the borrower's net utility is therefore equal to the social surplus brought about by the
investment:

\[U_b = (p_H R - 1)I. \qquad (3.10)\]

From (3.10) it is optimal for the borrower to invest as much as possible. The upper bound on
investment and in turn her borrowing capacity (``outside financing capacity'' or ``debt capacity'')
are determined by constraints (IC\(_b\)) and (IR\(_l\)). Substituting (IC\(_b\)) into (IR\(_l\)), we
obtain

\[I \leqslant kA, \qquad (3.11)\]

where

\[k = \frac{1}{1 - p_H(R - B/\Delta p)} > 1. \qquad (3.12)\]

{[}\ldots{]} Another important concept {[}\ldots{]} is the shadow value \(v\) of equity (here cash).
The entrepreneur derives gross utility \(v > 1\) from one more unit of equity. Letting
\(U_b^g \equiv A + U_b\) denote the borrower's gross utility, and using (3.10) and (3.11), we have

\[U_b^g \equiv vA, \qquad (3.14)\]

where the shadow value of equity is

\[v = \frac{p_H B/\Delta p}{1 - p_H(R - B/\Delta p)} > 1. \qquad (3.15)\]

\emph{From Tirole, Supplementary Section §6.4 (Contract Design by an Informed Party: An
Introduction), pp.~264--268 (PDF 277--281) --- the ``issuance game,'' the separating program and the
uniqueness threshold \(\alpha^*\) used in part (iii):}

A borrower who attempts to raise funds from lenders has private information about some
characteristic (private benefit, value of assets in place, prospects of the firm, value of
collateral) that affects the lenders' payoff. The borrower may have type \(b\) or \(\tilde{b}\).
While the borrower knows her type, the lenders only know that this type is \(b\) with probability
\(\alpha\) and \(\tilde{b}\) with probability \(1 - \alpha\).

\emph{Figure 6.4 (timeline, left to right):} Borrower offers a contract → Investors accept/refuse →
Borrower picks action(s).

Let us, abstractly, denote the contractual terms faced by the borrower by \(c\). Let \(U_b(c)\) and
\(\tilde{U}_b(c)\) denote the two types' net utilities for arbitrary contractual terms
\(c\).\textsuperscript{68} Let \(U_l(c)\) and \(\tilde{U}_l(c)\) denote the investors' expected
profit when contractual terms are \(c\) and the borrower turns out to have type \(b\) and
\(\tilde{b}\), respectively.

\emph{Example (privately-known-prospects).} In Section 6.2, the borrower had possible types
\(b = p\) and \(\tilde{b} = q\). The contractual terms \(c\) were just the borrower's reward
\(R_b^S\) in the case of success. More generally, they also contain the probability of investment,
\(x\), her reward in the case of failure, \(R_b^F\), and in the absence of investment, \(R_b^0\),
even though the latter in equilibrium can be taken to be equal to 1, 0, and 0, respectively (see
Section 6.5). We have

\[U_b(c) = x[pR_b^S + (1 - p)R_b^F] + (1 - x)R_b^0,\]
\[\tilde{U}_b(c) = x[qR_b^S + (1 - q)R_b^F] + (1 - x)R_b^0,\]
\[U_l(c) = x[p(R - R_b^S) - (1 - p)R_b^F] - (1 - x)R_b^0,\]
\[\tilde{U}_l(c) = x[q(R - R_b^S) - (1 - q)R_b^F] - (1 - x)R_b^0.\]

{[}\ldots{]} Figure 6.4 describes the timing of the \emph{issuance game}. As earlier, we assume that
the borrower designs the issue and offers the associated claims to a competitive capital market.
Investors purchase the claims if and only if they expect a nonnegative profit. Lastly, the borrower
chooses some action(s).

A few clarifications are in order. First, we allow for post-contracting actions by the borrower in
order to accommodate situations in which the borrower can waste resources (as in Chapters 3--5).
Second, we said that investors subscribe ``if and only if they expect a nonnegative profit.'' The
expectation should be taken relative to the updated beliefs, that is, the investors' beliefs after
they observe the contract offer and thus possibly learn something about the borrower. Third, we will
analyze the perfect Bayesian equilibrium (or equilibria) of the issuance game.\textsuperscript{70}
Fourth, a ``contract'' can in principle be anything that the borrower sees fit to design. However,
for the purpose of the analysis, it can be shown that there is no loss of generality in assuming
that the borrower offers an ``option contract'' \((c, \tilde{c})\), that is, as many contractual
terms as there are possible types. {[}\ldots{]} It can further be shown that there is no loss of
generality in assuming that the option contract is ``incentive compatible,'' that is, that type
\(b\) prefers contractual terms \(c\) to \(\tilde{c}\) and type \(\tilde{b}\) prefers \(\tilde{c}\)
to \(c\) {[}\ldots{]}.

{[}\ldots{]} Let an \emph{allocation} be a pair of (possibly identical) type-contingent contractual
terms \((c, \tilde{c})\). That is, it defines contractual terms \(c\) for type \(b\) and
\(\tilde{c}\) for type \(\tilde{b}\). (Note that an option contract defines an allocation.)

\emph{Definition.} An allocation \((c, \tilde{c})\) is \emph{incentive compatible} if type \(b\)
prefers \(c\) to \(\tilde{c}\) and type \(\tilde{b}\) prefers \(\tilde{c}\) to \(c\):

\[U_b(c) \geqslant U_b(\tilde{c}) \quad \text{and} \quad \tilde{U}_b(\tilde{c}) \geqslant \tilde{U}_b(c).\]

{[}\ldots{]}

\emph{Definition.} An incentive-compatible allocation \((c, \tilde{c})\) is \emph{profitable
type-by-type} if

\[U_l(c) \geqslant 0 \quad \text{and} \quad \tilde{U}_l(\tilde{c}) \geqslant 0.\]

\emph{Definition.} An incentive-compatible allocation \((c, \tilde{c})\) is \emph{profitable in
expectation} (relative to the prior beliefs) if

\[\alpha U_l(c) + (1 - \alpha)\tilde{U}_l(\tilde{c}) \geqslant 0.\]

An incentive-compatible allocation \((c, \tilde{c})\) that is profitable in expectation is
\emph{interim efficient} if it is Pareto-optimal for the two types of borrower in the set of
incentive-compatible, profitable-in-expectation allocations.

{[}\ldots{]} \emph{Definition.} Utility \(U_b(c_0)\) for borrower type \(b\) is the
\emph{low-information-intensity optimum} for that type if \(c_0\) maximizes type \(b\)'s utility in
the set of incentive-compatible, profitable-type-by-type allocations. That is, it is (part of) the
solution to the following program.

\emph{Program I (type \(b\)):}

\[\max_{\{c, \tilde{c}\}} U_b(c) \quad \text{s.t.} \quad U_b(c) \geqslant U_b(\tilde{c}), \quad \tilde{U}_b(\tilde{c}) \geqslant \tilde{U}_b(c), \quad U_l(c) \geqslant 0, \quad \tilde{U}_l(\tilde{c}) \geqslant 0.\]

The low-information-intensity optimum \(\tilde{c}_0\) for type \(\tilde{b}\) is defined similarly
(Program I (type \(\tilde{b}\))).

The payoff pair \((U_b(c_0), U_b(\tilde{c}_0))\) is called the \emph{low-information-intensity
optimum}. (By abuse of terminology, we will sometimes call the allocation \((c_0, \tilde{c}_0)\)
itself the low-information-intensity optimum.) {[}\ldots{]}

\textbf{Lemma 6.1.} \emph{The borrower can guarantee herself her low-information-intensity optimum
(\(U_b(c_0)\) if she has type \(b\), and \(\tilde{U}_b(\tilde{c}_0)\) if she has type
\(\tilde{b}\)).}

{[}\ldots{]}

\textbf{Proposition 6.1.}

\begin{enumerate}
\def\labelenumi{(\alph{enumi})}
\item
  \emph{The issuance game has a unique perfect Bayesian equilibrium if the low-information-intensity
  optimum is interim efficient (relative to prior beliefs \((\alpha, 1 - \alpha)\)). The borrower
  then obtains her low-information-intensity optimum (\(U_b(c_0)\) for type \(b\),
  \(\tilde{U}_b(\tilde{c}_0)\) for type \(\tilde{b}\)).}
\item
  \emph{If the low-information-intensity optimum is not interim efficient, then the set of
  equilibrium payoffs for the two types of borrowers is the set of payoffs that result from an
  incentive-compatible, profitable-in-expectation allocation and (weakly) Pareto-dominate the
  low-information-intensity optimum.}
\end{enumerate}

{[}\ldots{]} Let \(\tilde{c}^{SI}\) denote the \emph{symmetric information} contractual terms for
the bad borrower. It solves

\[\max_{\{\tilde{c}\}} \tilde{U}_b(\tilde{c}) \quad \text{s.t.} \quad \tilde{U}_l(\tilde{c}) \geqslant 0.\]

\textbf{Assumption 6.4 (weak monotonic profit).} \emph{Investors make a nonnegative profit if the
contractual terms are those of the bad borrower under symmetric information and the borrower is a
good one:}

\[U_l(\tilde{c}^{SI}) \geqslant 0.\]

This assumption is in general satisfied when both types are creditworthy under symmetric information
(and it is satisfied in all of our illustrations). It is always satisfied when the bad borrower is
not creditworthy under symmetric information: in that case, \(\tilde{c}^{SI}\) is the absence of
funding and thus \(U_l(\tilde{c}^{SI}) = 0\).

\emph{Definition.} The \emph{separating allocation} is the allocation \(c^*\) for the good borrower
and the symmetric-information contractual terms \(\tilde{c}^{SI}\) for the bad borrower, where
\(c^*\) maximizes the good borrower's payoff subject to the investors breaking even for the good
borrower and to the bad borrower not preferring \(c^*\) to \(\tilde{c}^{SI}\):

\[\max_{\{c\}} U_b(c) \quad \text{s.t.} \quad U_l(c) \geqslant 0, \quad \tilde{U}_b(c) \leqslant \tilde{U}_b(\tilde{c}^{SI}).\]

Note that the separating allocation is profitable type-by-type.

\textbf{Lemma 6.2.} \emph{Under the weak monotonic-profit assumption, the separating allocation is
the low-information-intensity optimum.}

\emph{Proof.} First, note that the bad borrower's symmetric-information program has the same
objective function and fewer constraints than her low-information-intensity optimum program, Program
I (type \(\tilde{b}\)). Hence,

\[\tilde{U}_b^{SI} \equiv \tilde{U}_b(\tilde{c}^{SI}) \geqslant \tilde{U}_b(\tilde{c}_0).\]

{[}\ldots{]}

\textbf{Lemma 6.3.} \emph{Under the weak monotonic-profit assumption, there exists a threshold level
\(\alpha^*\) for prior beliefs such that the low-information-intensity optimum (that is, the
separating allocation under the weak monotonic-profit assumption) is interim efficient if and only
if \(\alpha \leqslant \alpha^*\).}

\begin{quote}
\begin{enumerate}
\def\labelenumi{\arabic{enumi}.}
\setcounter{enumi}{67}
\item
  So, for example, if the borrower has initial cash \(A\), and has quasi-linear preferences, the net
  utility is equal to the gross utility minus \(A\).
\item
  Very roughly speaking, a perfect Bayesian equilibrium of a game is a set of strategies and beliefs
  such that at any stage of the game players act optimally given their beliefs at that stage (the
  equilibrium is ``perfect'') and beliefs are updated by the players according to Bayes' rule using
  equilibrium strategies and observed actions (the updating is Bayesian). {[}\ldots{]} Here
  investors are assumed to update their beliefs about the borrower's type using the borrower's
  equilibrium type-contingent contract offer and the actual contract offer. The previous condition,
  that they subscribe if and only if they expect a nonnegative profit given their updated beliefs,
  is an optimization requirement.
\end{enumerate}
\end{quote}
\stmtend
%% END-STATEMENT

\asked (i) Write the separating program: maximise the good borrower's utility over her contract, subject to investors breaking even on her, to the bad borrower not preferring the good borrower's terms to her own symmetric-information terms, and to the good borrower working. (ii) Show that the good borrower takes no reward in failure, $R_b^F=0$. (iii) Using the supplementary section, show that the separating outcome is the unique perfect Bayesian equilibrium of the issuance game iff the prior $\alpha$ is at most some $\alpha^*$.

\approach This is the pecking-order application of the privately-known-prospects model (Tirole \S6.2.2, Application 3; lecture block 3.5) with the continuous-investment technology of \S3.4. The separating program is a constrained maximisation; part (ii) is a perturbation argument: shift the reward from failure to success keeping the good type's expected reward fixed, and check that every constraint is relaxed, which, by the information-sensitivity of claims, makes investors hold the safe (debt-like) part of the cash flow. Part (iii) applies the general results on contract design by an informed party (\S6.4: Lemmas 6.2--6.3 and Proposition 6.1) after checking their key assumption, weak monotonic profit.

\notation
\begin{itemize}
  \item $I\ge0$: investment; $R^SI$ / $R^FI$: income in success / failure, $R^S>R^F\ge0$.
  \item $p_H,p_L$ (good) and $q_H,q_L$ (bad): success probabilities working / shirking, $p_H-p_L=q_H-q_L=\Delta p>0$, $q_H<p_H$.
  \item $BI$: private benefit from shirking; $A$: entrepreneur's cash; investors demand a 0 return.
  \item $R_b^S,R_b^F\ge0$: good borrower's rewards in success / failure; $\tilde U_b^{SI}$: bad borrower's symmetric-information gross utility.
  \item $\alpha$: prior probability of the good type; $\alpha^*$: threshold of part (iii).
\end{itemize}

\assume (a) Investment has positive NPV for the good type when she works,
$N\equiv p_HR^S+(1-p_H)R^F-1>0$, and finite scale, $p_H(R^S-B/\Delta p)+(1-p_H)R^F<1$ (analogues of
(3.7) and (3.9)). (b) At the optimum, investors hold a positive claim in success, $R_b^S<R^SI$. (c) Footnote 80's
formula is for $R^F=0$ ($R=R^S$); with $R^F>0$, $\tilde U_b^{SI}$ is computed the same way with the
bad type's expected income per unit.

\step{1}{Payoffs, and why the bad type works when she mimics} The good type's gross utility is
$p_HR_b^S+(1-p_H)R_b^F$. A bad type taking the same terms works, because her gain from working is
\begin{align}
  [q_HR_b^S+(1-q_H)R_b^F]-[q_LR_b^S+(1-q_L)R_b^F+BI]
  = (\Delta p)(R_b^S-R_b^F)-BI, \nonumber
\end{align}
which is the good type's incentive constraint: the two types share the same (IC$_b$). Her utility from
mimicking is therefore $q_HR_b^S+(1-q_H)R_b^F$.

\step{2}{The separating program (part (i))}
\begin{align}
  \max_{\{I,R_b^S,R_b^F\}}\;& p_HR_b^S+(1-p_H)R_b^F \label{t3e12-obj}\\
  \text{s.t.}\;& p_H(R^SI-R_b^S)+(1-p_H)(R^FI-R_b^F)\ge I-A, \tag{IR$_l$}\\
  & q_HR_b^S+(1-q_H)R_b^F\le\tilde U_b^{SI}, \tag{M}\\
  & (\Delta p)(R_b^S-R_b^F)\ge BI, \tag{IC$_b$}\\
  & 0\le R_b^S\le R^SI,\quad 0\le R_b^F\le R^FI. \nonumber
\end{align}
The book's answer lists only the rewards as choice variables, but $I$ is chosen too. With $R^F=0$,
footnote 80's value comes from the bad type's symmetric-information problem in \S3.4.2. Her net utility is the NPV $(q_HR-1)I$, increasing in $I$ when $q_HR\ge1$, so she borrows as much as the constraints allow: (IC$_b$) binds, $R_b=BI/\Delta p$, and substituting it into the break-even constraint with equality,
\begin{align}
  q_H\Bigl(RI-\frac{BI}{\Delta p}\Bigr) &= I-A \nonumber\\
  A &= I\Bigl[1-q_H\Bigl(R-\frac{B}{\Delta p}\Bigr)\Bigr] \nonumber\\
  I &= \frac{A}{1-q_H(R-B/\Delta p)}. \nonumber
\end{align}
Her gross utility is the cash plus the NPV,
\begin{align}
  \tilde U_b^{SI} &= A+(q_HR-1)\frac{A}{1-q_H(R-B/\Delta p)} \nonumber\\
  &= \Bigl[1+\frac{q_HR-1}{1-q_H(R-B/\Delta p)}\Bigr]A \qquad (q_HR\ge1). \label{t3e12-USI}
\end{align}

\answer{(i)} Program \eqref{t3e12-obj}: maximise $p_HR_b^S+(1-p_H)R_b^F$ over $\{I,R_b^S,R_b^F\}$ subject to (IR$_l$), (M) with $\tilde U_b^{SI}$ from footnote 80 (or $A$ if $q_HR<1$), (IC$_b$) and limited liability.

\step{3}{A reward in failure can be moved to success at no cost (part (ii))} Suppose an optimum has
$R_b^F>0$. Change the rewards by $\delta R^S>0$ and $\delta R^F<0$ with
\begin{align}
  p_H\,\delta R^S+(1-p_H)\,\delta R^F = 0
  \iff \delta R^F = -\frac{p_H}{1-p_H}\,\delta R^S. \label{t3e12-pert}
\end{align}
The objective \eqref{t3e12-obj} is unchanged by \eqref{t3e12-pert}. So is (IR$_l$): investors receive
expected income minus the borrower's expected reward, and both are unchanged. (IC$_b$) is relaxed:
its left side changes by $(\Delta p)(\delta R^S-\delta R^F)>0$. (M) is relaxed too; substituting
\eqref{t3e12-pert},
\begin{align}
  q_H\,\delta R^S+(1-q_H)\,\delta R^F
  &= q_H\,\delta R^S-(1-q_H)\frac{p_H}{1-p_H}\,\delta R^S \nonumber\\
  &= \frac{q_H(1-p_H)-p_H(1-q_H)}{1-p_H}\,\delta R^S \nonumber\\
  &= -\frac{p_H-q_H}{1-p_H}\,\delta R^S < 0. \label{t3e12-dM}
\end{align}
By assumption (b), a small $\delta R^S$ keeps $R_b^S\le R^SI$.

\step{4}{Slack constraints allow a strict improvement, so $R_b^F=0$} After the change, (M) and (IC$_b$)
are slack. Now raise $I$ by $\varepsilon>0$ and $R_b^S$ by $\eta=N\varepsilon/p_H>0$. Each constraint changes as follows:
(IR$_l$), left side minus right side, changes by
\begin{align}
  p_H(R^S\varepsilon-\eta)+(1-p_H)R^F\varepsilon-\varepsilon
     &= [p_HR^S+(1-p_H)R^F-1]\varepsilon-p_H\eta \nonumber\
     &= N\varepsilon-N\varepsilon = 0; \nonumber
\end{align}
(IC$_b$), left minus right, changes by $(\Delta p)\eta-B\varepsilon$; the left side of (M) by $+q_H\eta$; and the limited-liability margin $R^SI-R_b^S$ by $R^S\varepsilon-\eta=(R^S-N/p_H)\varepsilon$.
For $\varepsilon$ small enough the slack in (IC$_b$), (M) and $R_b^S<R^SI$ absorbs the last three changes, whatever their sign,
and the objective rises by $p_H\eta=N\varepsilon>0$. That contradicts optimality, so
\begin{align}
  R_b^F = 0. \qquad\blacksquare \label{t3e12-RF0}
\end{align}
The same point in one line: per unit of the bad type's mimicking utility, a reward in success buys
$p_H/q_H>1$ units of the good type's utility and a reward in failure buys $(1-p_H)/(1-q_H)<1$. Investors
receive the whole failure income $R^FI$: they hold the safe, information-insensitive part of the cash
flow, i.e.\ debt. That is the pecking order.

\answer{(ii)} $R_b^F=0$: moving reward from failure to success at constant expected cost leaves (IR$_l$) and the objective unchanged, relaxes (IC$_b$), and relaxes (M) by $(p_H-q_H)\delta R^S/(1-p_H)$; the slack then allows a larger, more profitable investment. Investors hold all of $R^FI$, a debt-like claim.

\step{5}{Uniqueness of the separating outcome iff $\alpha\le\alpha^*$ (part (iii))} First check
Assumption 6.4 (weak monotonic profit). If the bad type is not creditworthy, $\tilde c^{SI}$ is no funding
and $U_l(\tilde c^{SI})=0$. Otherwise her symmetric-information terms are $\tilde I$,
$\tilde R_b^S=B\tilde I/\Delta p$, $\tilde R_b^F=0$, with
$q_H(R^S\tilde I-\tilde R_b^S)+(1-q_H)R^F\tilde I=\tilde I-A$. A good type on those terms works (Step~1),
and the investors' profit is
\begin{align}
  U_l(\tilde c^{SI}) &= p_H(R^S\tilde I-\tilde R_b^S)+(1-p_H)R^F\tilde I-(\tilde I-A) \nonumber\\
  &= p_H(R^S\tilde I-\tilde R_b^S)+(1-p_H)R^F\tilde I-\bigl[q_H(R^S\tilde I-\tilde R_b^S)+(1-q_H)R^F\tilde I\bigr] \nonumber\\
  &= (p_H-q_H)\Bigl[\Bigl(R^S-\frac{B}{\Delta p}\Bigr)-R^F\Bigr]\tilde I \;\ge\; 0
     \quad\text{iff } R^S-\frac{B}{\Delta p}\ge R^F, \label{t3e12-wmp}
\end{align}
which always holds when $R^F=0$ (footnote 80's case). Under Assumption 6.4:
\begin{itemize}
  \item Lemma 6.2: the separating allocation (Step~2 for the good type, $\tilde c^{SI}$ for the bad) is the low-information-intensity optimum;
  \item Lemma 6.3: there is $\alpha^*$ such that this optimum is interim efficient iff $\alpha\le\alpha^*$;
  \item if $\alpha\le\alpha^*$, Proposition 6.1(a) says the issuance game has a unique PBE, and the borrower gets the low-information-intensity optimum, i.e.\ the separating outcome;
  \item if $\alpha>\alpha^*$, the optimum is not interim efficient: some incentive-compatible allocation that is profitable in expectation Pareto-dominates it, strictly for at least one type. By Proposition 6.1(b) that allocation's payoffs are equilibrium payoffs, so the separating outcome is not the only PBE.
\end{itemize}
Hence the separating outcome is the unique PBE iff $\alpha\le\alpha^*$. $\blacksquare$ When good types are
common, a pooling contract with a small cross-subsidy to the rare bad types costs the good type less than
the distortion she needs to separate, so both types can be made better off.

\answer{(iii)} Weak monotonic profit holds (always if $R^F=0$; in general iff $R^S-B/\Delta p\ge R^F$), so the separating allocation is the low-information-intensity optimum (Lemma 6.2). It is interim efficient iff $\alpha\le\alpha^*$ (Lemma 6.3), and then it is the unique PBE outcome (Proposition 6.1(a)); for $\alpha>\alpha^*$ other, Pareto-better equilibria exist (Proposition 6.1(b)).

\step{6}{Numerical check of \eqref{t3e12-USI} and \eqref{t3e12-dM}} With $q_H=0.6$, $R=3$,
$B/\Delta p=2$, $A=1$: $q_H(R-B/\Delta p)=0.6(3-2)=0.6$, so $\tilde I=1/(1-0.6)=1/0.4=2.5$ and, with $q_HR=1.8$,
$\tilde U_b^{SI}=1+(1.8-1)(2.5)=1+2=3=[1+0.8/0.4](1)$. With $p_H=0.8$ and $\delta R^S=0.01$, \eqref{t3e12-pert} gives
$\delta R^F=-(0.8/0.2)(0.01)=-0.04$, and the change in (M) is $(0.6)(0.01)+(0.4)(-0.04)=0.006-0.016=-0.01=-(0.2/0.2)(0.01)$, as \eqref{t3e12-dM} predicts.

\intuition Under asymmetric information the good type's problem is to give investors a claim whose value
depends as little as possible on her type. Income in failure is worth relatively more to the bad type (she
fails more often), so promising herself anything there invites mimicry. The good type takes her reward
only in success and gives investors all the safe salvage value $R^FI$. Investors thus hold debt-like
claims first, and that is the pecking order.

\pitfalls
\begin{itemize}
  \item Writing the mimicking constraint with the bad type shirking: because $p_H-p_L=q_H-q_L$, she faces the same (IC$_b$) and works under the good type's terms, so her mimicking utility uses $q_H$.
  \item Treating $I$ as fixed: with variable investment $I$ is a choice variable, and it is the slack created by the perturbation that lets $I$ (and the good type's utility) rise.
  \item Using the good type's symmetric-information utility in (M): the bad type compares the good type's terms with \emph{her own} symmetric-information utility $\tilde U_b^{SI}$.
\end{itemize}

\bookans Matches the book's answer (pp.~599--600): the same program (with $I$ made explicit here), the
same perturbation argument for $R_b^F=0$, and uniqueness iff $\alpha\le\alpha^*$. The book invokes
``Proposition 6.2''; with the excerpt copied in the list, the argument is Lemmas 6.2--6.3 and Proposition
6.1. The book asserts that weak monotonic profit holds; \eqref{t3e12-wmp} shows it needs
$R^S-B/\Delta p\ge R^F$, automatic in footnote 80's case $R^F=0$.

\refbib{Tirole Ch.~6, \S6.2.2 Application 3 ``Pecking-Order Hypothesis'', pp.~246--249; \S3.4.1--3.4.2,
pp.~127--128; supplementary \S6.4, pp.~264--268; Exercise 6.8, p.~278; answer pp.~599--600. Lecture
Topic~3, pecking-order slides.}
%% END-FROM

% =============================================================================
\exsec{15. Reputational capital --- Tirole 6.3, pp.~274--275}
%% BEGIN-STATEMENT 15
\stmtbegin
\textbf{Source:} Tirole, Ch. 6, Exercise 6.3, pp.~274--275 (PDF 287) · Answer in book: yes, p.~597

Consider the fixed-investment model. All parameters are common knowledge between the borrower and
the investors, except the private benefit which is known only to the borrower. The private benefit
is equal to \(B\) with probability \(1 - \alpha\) and to \(b\) with probability \(\alpha\), where
\(B > b > 0\).

\begin{enumerate}
\def\labelenumi{(\roman{enumi})}
\tightlist
\item
  Consider first the one-period adverse-selection problem. Suppose that the borrower has assets
  \(A > 0\) such that
\end{enumerate}

\[p_H\left(R - \frac{b}{\Delta p}\right) > I - A > \max\left[p_H\left(R - \frac{B}{\Delta p}\right), p_L R\right].\]

Show that the project receives funding if and only if

\[(p_H - (1 - \alpha)\Delta p)\left(R - \frac{b}{\Delta p}\right) \geqslant I - A.\]

\begin{enumerate}
\def\labelenumi{(\roman{enumi})}
\setcounter{enumi}{1}
\tightlist
\item
  Suppose now that there are two periods (\(t = 1, 2\)). The second period is described as in
  question (i), except that the belief \(\tilde{\alpha}\) at date 2 is the posterior belief updated
  from the prior belief \(\alpha\), and that the borrower has cash \(A\) only if she has been
  successful at date 1 (and has 0 and is not funded if she has been unsuccessful). So, suppose that
  the first-period project is funded and that the borrower receives at the end of date 1 a reward
  \(A\) when successful and 0 when unsuccessful. The first-period funding is project finance and
  does not specify any funding for the second project. Suppose for notational simplicity that the
  private benefit is the same (\(B\) or \(b\)) in period 1 and in period 2. Let \(\Delta p_1\)
  denote the increase in the probability of success when diligent in period 1. Assume that
\end{enumerate}

\[b < (\Delta p_1)A < B < (\Delta p_1)\left[p_L\left(R - \frac{I - A}{p_H - (1 - \alpha_S)\Delta p}\right) + B\right]\]

and

\[(p_H - (1 - \alpha)\Delta p)\left(R - \frac{b}{\Delta p}\right) < I - A < (p_H - (1 - \alpha_S)\Delta p)\left(R - \frac{b}{\Delta p}\right),\]

where \(1 - \alpha_S \equiv (1 - \alpha)p_L/((1 - \alpha)p_L + \alpha p_H)\).

A ``pooling equilibrium'' is an equilibrium in which the borrower's first-period effort is
independent of her private benefit. A ``separating equilibrium'' is (here) an equilibrium in which
the \(b\)-type works and the \(B\)-type shirks in period 1. A ``semiseparating'' equilibrium is
(here) an equilibrium in which in period 1 the \(b\)-type works and the \(B\)-type randomizes
between working and shirking.

\begin{itemize}
\tightlist
\item
  Show that there exists no pooling and no separating equilibrium.
\item
  Compute the semiseparating equilibrium. Does this model formalize the notion of reputational
  capital?
\end{itemize}

\paragraph{Context needed to solve it}

\emph{From Tirole §3.2.1--3.2.2 (the fixed-investment model), pp.~115--117 (PDF 128--130):}

\textbf{3.2.1 The Fixed-Investment Model.} Variants of the following entrepreneurial model will be
used in the following: an entrepreneur (also called the ``insider'' or the ``borrower,'' ``she'')
has a project. This project requires a fixed investment \(I\). The entrepreneur initially has
``assets'' (``cash on hand'' or ``net worth'') \(A < I\). For the moment we interpret these assets
as being cash or liquid securities that can be used toward covering the cost of investment.
{[}\ldots{]} The entrepreneur's cash can either be invested in the project or used for consumption.
To implement the project the entrepreneur must borrow \(I - A\) from lenders. {[}\ldots{]}

\emph{Project.} If undertaken, the project either succeeds, that is, yields verifiable income
\(R > 0\), or fails, that is, yields no income. The probability of success is denoted by \(p\). The
project is subject to moral hazard. The entrepreneur can ``behave'' (``work,'' ``exert effort,''
``take no private benefit'') or ``misbehave'' (``shirk,'' ``take a private benefit''); or,
equivalently, the entrepreneur chooses between a project with a high probability of success and
another project which ceteris paribus she prefers (is easier to implement, is more fun, has greater
spinoffs in the future for the entrepreneur, benefits a friend, delivers perks, is more
``glamorous,'' etc.) but has a lower probability of success.

Behaving yields probability \(p = p_H\) of success and no private benefit to the entrepreneur, and
misbehaving results in probability \(p = p_L < p_H\) of success and private benefit \(B > 0\)
(measured in units of account) to the entrepreneur. In the ``effort interpretation,'' \(B\) can also
be interpreted as a disutility of effort saved by the entrepreneur when shirking. Let
\(\Delta p \equiv p_H - p_L\).

\emph{Preferences and the loan agreement.} Both the borrower and the potential lenders (or
``investors'') are risk neutral. For notational simplicity, there is no time preference; the rate of
return expected by investors (which is also the riskless rate, due to risk neutrality) is taken to
be 0. The borrower is protected by limited liability, and so her income cannot take negative values.

Lenders behave competitively in the sense that the loan, if any, makes zero profit. That is, we have
in mind that several prospective lenders compete for issuing a loan to the borrower, and that, if
the most attractive loan offer made a positive profit, the borrower could turn to an alternative
lender and offer to switch for a slightly lower interest rate. {[}\ldots{]}

Let us turn to the loan contract. A contract first stipulates whether the project is financed. If
so, it further specifies how the profit is shared between the lenders and the borrower. The
borrower's limited liability will imply that both sides receive 0 in the case of failure (the gross
payoffs are the \emph{ex post} monetary payoffs and take no account of past investments and private
benefit). Intuitively, there is no point in specifying a positive transfer from the lenders to the
borrower, as such a transfer can only weaken incentives, while it has no insurance benefit under
risk neutrality. This property will be proved more rigorously and is here taken for granted. In the
case of success, the two parties share the profit \(R\); \(R_b\) goes to the borrower and \(R_l\) to
the lenders. To sum up, we posit an incentive scheme for the entrepreneur of the following form:
\(R_b\) in the case of success, 0 in the case of failure.

The zero-profit constraint for the lenders can be written as

\[p_H R_l = I - A,\]

assuming that the loan agreement induces the borrower to behave (which under our assumptions will be
the case). {[}\ldots{]} We summarize the timing in Figure 3.1.

\emph{Figure 3.1 (timeline, left to right):} Loan agreement (sharing rule) → Investment → Moral
hazard → Outcome.

We assume that the project is viable only in the absence of moral hazard. That is, the project has
positive NPV if the entrepreneur behaves,

\[p_H R - I > 0, \qquad (3.1)\]

but negative NPV, even if one includes the borrower's private benefit, if she does not,

\[p_L R - I + B < 0. \qquad (3.2)\]

{[}\ldots{]}

\textbf{3.2.2 The Lenders' Credit Analysis.} Because the project has negative NPV in the case of
misbehavior, the loan agreement must be careful to preserve enough of a stake for the borrower in
the enterprise. The borrower faces the following tradeoff once the financing has been secured: by
misbehaving, she obtains private benefit \(B\), but she reduces the probability of success from
\(p_H\) to \(p_L\). Because she has stake \(R_b\) in the firm's income (she receives \(R_b\) in the
case of success and 0 in the case of failure), the borrower will therefore behave if the following
``incentive compatibility constraint'' is satisfied:

\[p_H R_b \geqslant p_L R_b + B \quad \text{or} \quad (\Delta p) R_b \geqslant B. \qquad (\text{IC}_b)\]

From this incentive compatibility constraint we infer that the highest income in the case of success
that can be pledged to the lenders without jeopardizing the borrower's incentives is
\(R - \dfrac{B}{\Delta p}\). The \emph{(expected) pledgeable income} is then
\(\mathcal{P} = p_H\left(R - \dfrac{B}{\Delta p}\right)\).

Because the lenders must break even in order to be willing to finance the project, a necessary
condition for the borrower to receive a loan is that the expected pledgeable income exceed the
lenders' initial outlay:

\[\mathcal{P} \equiv p_H\left(R - \frac{B}{\Delta p}\right) \geqslant I - A, \qquad (\text{IR}_l)\]

where ``IR\(_l\)'' stands for the lenders' individual rationality constraint (which we will also
often call the ``breakeven constraint'' or the ``participation constraint''). Thus a necessary
condition for financing to be arranged is

\[A \geqslant \bar{A} = p_H \frac{B}{\Delta p} - (p_H R - I). \qquad (3.3)\]

{[}\ldots{]} Conversely, if \(A \geqslant \bar{A}\), the entrepreneur is able to secure financing,
and so condition (3.3) is both a necessary and a sufficient condition for financing.
\stmtend
%% END-STATEMENT

\asked (i) In one period with privately known private benefit ($b$ with prob.\ $\alpha$, $B$ with prob.\ $1-\alpha$), show that the project is funded iff $(p_H-(1-\alpha)\Delta p)(R-b/\Delta p)\ge I-A$. (ii) In two periods, where date-2 cash $A$ comes only from date-1 success and the date-2 belief is the posterior: show that there is no pooling and no separating equilibrium in date-1 effort; compute the semiseparating equilibrium (the $B$-type randomizes); say whether this formalizes reputational capital.

\approach (i) is the pooling analysis of Exercise 6.1 with assets $A$: under asymmetric information all types accept the same contract, so the lender's revenue uses the pooled success probability of those who behave. (ii) is a dynamic signalling (career-concerns) problem: date-1 effort affects the probability of success, success moves the market's belief (Bayes' rule), and the belief decides date-2 funding terms. The value of a good reputation is the date-2 rent; the $B$-type's date-1 incentive constraint plus Bayes' rule pin down the equilibrium.

\notation
\begin{itemize}
  \item Date 2 (as in (i)): success probability $p_H$ (work) or $p_L$ (shirk), $\Delta p=p_H-p_L$; income $R$; investment $I$; borrower cash $A$; private benefit $b$ or $B$, $B>b>0$.
  \item Date 1: success probability $p^1_H$ (work) or $p^1_L$ (shirk), $\Delta p_1=p^1_H-p^1_L$; reward $A$ on success, 0 on failure.
  \item $\alpha$: prior that the type is $b$; $\tilde\alpha$: posterior after a date-1 success; $m(\tilde\alpha)\equiv p_H-(1-\tilde\alpha)\Delta p=\tilde\alpha p_H+(1-\tilde\alpha)p_L$.
  \item $y\in[0,1]$: probability that the $B$-type works at date 1; $V^B_S(\tilde\alpha)$, $V^b_S(\tilde\alpha)$: date-2 continuation values after a date-1 success.
\end{itemize}

\assume The statement defines $1-\alpha_S=(1-\alpha)p_L/((1-\alpha)p_L+\alpha p_H)$, i.e.\ it uses the date-2 probabilities in Bayes' rule. If date-1 probabilities equalled the date-2 ones ($\Delta p_1=\Delta p$), the assumptions would be inconsistent (Step~4 shows they need $\Delta p_1>\Delta p/p_H$, impossible when $\Delta p_1=\Delta p$ because $p_H<1$). We therefore let date-1 probabilities differ and use them in Bayes' rule: $1-\alpha_S=(1-\alpha)p^1_L/((1-\alpha)p^1_L+\alpha p^1_H)$. We also assume, as the book implicitly does, that the $b$-type prefers to reinvest $A$ at date 2 ($p_HR_b\ge A$).

\step{1}{(i) The three kinds of contract} Both types want funding and both prefer a larger $R_b$, so any contract offered is accepted by both (pooling). By (IC$_b$) a type works iff $(\Delta p)R_b\ge$ her private benefit. Case $R_b\ge B/\Delta p$: both work, lender revenue
\begin{align}
  p_H(R-R_b) &\le p_H\Bigl(R-\frac{B}{\Delta p}\Bigr) < I-A. \nonumber
\end{align}
Case $R_b<b/\Delta p$: both shirk, revenue $p_L(R-R_b)\le p_LR<I-A$. Case $b/\Delta p\le R_b<B/\Delta p$: the $b$-type works, the $B$-type shirks, and revenue is largest at $R_b=b/\Delta p$:
\begin{align}
  [\alpha p_H+(1-\alpha)p_L]\Bigl(R-\frac{b}{\Delta p}\Bigr) &= (p_H-(1-\alpha)\Delta p)\Bigl(R-\frac{b}{\Delta p}\Bigr), \label{r5e15-i}
\end{align}
using $\alpha p_H+(1-\alpha)p_L = p_H-(1-\alpha)(p_H-p_L)$. Only this case can cover $I-A$.
\answer{(i)} Funded iff $(p_H-(1-\alpha)\Delta p)(R-b/\Delta p)\ge I-A$; the contract makes the good type work and lets the bad type shirk.

\step{2}{(ii) Date-2 terms as a function of the belief} With belief $\tilde\alpha$, funding requires \eqref{r5e15-i} at $\tilde\alpha$. Let $\hat\alpha$ solve $m(\hat\alpha)(R-b/\Delta p)=I-A$; funding occurs iff $\tilde\alpha\ge\hat\alpha$. The borrower (competitive lenders) then gets
\begin{align}
  m(\tilde\alpha)(R-R_b) &= I-A \quad\Longrightarrow\quad R_b(\tilde\alpha) = R - \frac{I-A}{m(\tilde\alpha)}, \label{r5e15-Rb}
\end{align}
increasing in $\tilde\alpha$. Because $I-A>p_H(R-B/\Delta p)$ and $m<p_H$, $R_b(\tilde\alpha)<R-(I-A)/p_H<B/\Delta p$: the $B$-type shirks at date 2. Continuation values after a date-1 success:
\begin{align}
  V^B_S(\tilde\alpha) &= \begin{cases} p_LR_b(\tilde\alpha)+B & \tilde\alpha\ge\hat\alpha \\ A & \tilde\alpha<\hat\alpha\end{cases},\qquad
  V^b_S(\tilde\alpha) = \begin{cases} p_HR_b(\tilde\alpha) & \tilde\alpha\ge\hat\alpha \\ A & \tilde\alpha<\hat\alpha\end{cases}. \nonumber
\end{align}
After a failure the borrower has nothing and is not funded (value 0). Date-1 incentive constraint of a type with private benefit $\beta\in\{b,B\}$: work iff $p^1_HV_S \ge p^1_LV_S+\beta$, i.e.
\begin{align}
  (\Delta p_1)\,V_S &\ge \beta. \label{r5e15-IC1}
\end{align}
The $b$-type always works: $V^b_S\ge A$ (she can always keep the cash), so $(\Delta p_1)V^b_S\ge(\Delta p_1)A>b$.

\step{3}{No pooling equilibrium} Pooling means the $B$-type also works ($y=1$). Both types then succeed with $p^1_H$, so Bayes' rule leaves the belief unchanged:
\begin{align}
  \tilde\alpha &= \frac{\alpha p^1_H}{\alpha p^1_H+(1-\alpha)p^1_H} = \alpha. \nonumber
\end{align}
By the first assumed inequality, $m(\alpha)(R-b/\Delta p)<I-A$, so $\alpha<\hat\alpha$: no date-2 funding and $V^B_S=A$. Then $(\Delta p_1)A<B$ violates \eqref{r5e15-IC1} for the $B$-type, who would rather shirk. Contradiction.

\step{4}{No separating equilibrium} Separation means $y=0$. A success then raises the belief to
\begin{align}
  \alpha_S &= \frac{\alpha p^1_H}{\alpha p^1_H+(1-\alpha)p^1_L}, \qquad 1-\alpha_S = \frac{(1-\alpha)p^1_L}{(1-\alpha)p^1_L+\alpha p^1_H}. \nonumber
\end{align}
By the second assumed inequality, $\alpha_S>\hat\alpha$: funding at date 2, and $V^B_S=p_LR_b(\alpha_S)+B$. The assumption $B<(\Delta p_1)[p_LR_b(\alpha_S)+B]$ says \eqref{r5e15-IC1} holds strictly: the $B$-type would work. Contradiction. (Consistency: since $R_b<B/\Delta p$, this assumption implies $B<(\Delta p_1)(p_LB/\Delta p+B)=(\Delta p_1)Bp_H/\Delta p$, i.e.\ $\Delta p_1>\Delta p/p_H$, the condition used in the Assumption above.)

\step{5}{Semiseparating equilibrium: indifference} With $y\in(0,1)$ the $B$-type must be indifferent, \eqref{r5e15-IC1} with equality. Since $(\Delta p_1)A<B$, this needs funding, and
\begin{align}
  (\Delta p_1)\Bigl[p_L\Bigl(R-\frac{I-A}{m(\alpha'_S)}\Bigr)+B\Bigr] &= B. \label{r5e15-ind}
\end{align}
Solve for the reward first:
\begin{align}
  p_LR_b^* + B &= \frac{B}{\Delta p_1} \nonumber\\
  R_b^* &= \frac{B(1-\Delta p_1)}{\Delta p_1\,p_L}, \label{r5e15-Rbs}
\end{align}
then for the belief, from \eqref{r5e15-Rb}:
\begin{align}
  m(\alpha'_S) &= \frac{I-A}{R-R_b^*}, \qquad 1-\alpha'_S = \frac{1}{\Delta p}\Bigl[p_H - \frac{I-A}{R-R_b^*}\Bigr]. \label{r5e15-alpha}
\end{align}
An interior solution needs $R_b^*\ge b/\Delta p$ (so $\alpha'_S\ge\hat\alpha$), i.e.\ $(\Delta p_1)(p_Lb/\Delta p+B)\le B$; $\alpha'_S<\alpha_S$ follows from Step~4, because $V^B_S$ increases with the belief.

\step{6}{Semiseparating equilibrium: the mixing probability} If the $B$-type works with probability $y$, her date-1 success probability is $\pi_B=p^1_L+y\,\Delta p_1$, and Bayes' rule gives
\begin{align}
  \frac{1-\alpha'_S}{\alpha'_S} &= \frac{(1-\alpha)\,\pi_B}{\alpha\,p^1_H} \nonumber\\
  \pi_B &= \frac{\alpha\,p^1_H}{1-\alpha}\cdot\frac{1-\alpha'_S}{\alpha'_S}, \qquad y = \frac{\pi_B-p^1_L}{\Delta p_1}. \label{r5e15-y}
\end{align}
The more often the $B$-type works, the less a success says, which is why a unique $y$ makes her exactly indifferent.

\step{7}{Numbers} Date 2: $p_H=0.9$, $p_L=0.5$ ($\Delta p=0.4$), $R=10$, $I=8$, $A=1$, $b=0.2$, $B=2.4$; date 1: $p^1_H=0.9$, $p^1_L=0.1$ ($\Delta p_1=0.8$); $\alpha=0.5$. Assumptions: $p_H(R-b/\Delta p)=(0.9)(9.5)=8.55>7=I-A>\max[(0.9)(4),(0.5)(10)]=5$;
\begin{align}
  1-\alpha_S &= \frac{(0.5)(0.1)}{(0.5)(0.1)+(0.5)(0.9)} = 0.1, \nonumber\\
  m(\alpha_S) &= 0.9-(0.1)(0.4) = 0.86, \qquad m(\alpha) = 0.9-(0.5)(0.4) = 0.7, \nonumber\\
  m(\alpha)(9.5) &= 6.65 < 7 < m(\alpha_S)(9.5) = 8.17, \nonumber\\
  R_b(\alpha_S) &= 10-\frac{(7)}{(0.86)} = 1.8605, \qquad (\Delta p_1)A = 0.8\in(0.2,\,2.4), \nonumber\\
  (\Delta p_1)\bigl[p_LR_b(\alpha_S)+B\bigr] &= (0.8)[(0.5)(1.8605)+2.4] = 2.6642 > 2.4. \nonumber
\end{align}
Equilibrium, from \eqref{r5e15-Rbs}--\eqref{r5e15-y}:
\begin{align}
  R_b^* &= \frac{(2.4)(1-0.8)}{(0.8)(0.5)} = 1.2 \;\ge\; b/\Delta p = 0.5, \nonumber\\
  m(\alpha'_S) &= \frac{(7)}{(10-1.2)} = 0.79545, \qquad 1-\alpha'_S = \frac{0.9-0.79545}{0.4} = 0.26136, \qquad \alpha'_S = 0.73864, \nonumber\\
  \pi_B &= \frac{(0.5)(0.9)}{(0.5)}\cdot\frac{0.26136}{0.73864} = 0.31846, \qquad y = \frac{0.31846-0.1}{0.8} = 0.27308. \nonumber
\end{align}
Checks: $\hat\alpha=1-(0.9-7/9.5)/0.4=0.59211<\alpha'_S=0.73864<\alpha_S=0.9$; the $b$-type reinvests, $p_HR_b^*=(0.9)(1.2)=1.08\ge A=1$.
\answer{(ii)} No pooling, no separating equilibrium; the semiseparating one has the $B$-type working with probability $y$ from \eqref{r5e15-y} and posterior $\alpha'_S$ from \eqref{r5e15-alpha} (here $y=0.273$, $\alpha'_S=0.739$). Yes: the date-2 rent from a good name is reputational capital that disciplines date-1 behaviour.

\intuition A date-1 success is valuable not for its own cash but for what it does to the market's opinion: it raises the posterior and unlocks funding with a larger reward at date 2. That expected rent is reputational capital, and it gives even the bad type a reason to behave. But the reward to reputation cannot be too high or too low: if all bad types behaved, success would say nothing and the rent would vanish; if none did, the rent would be large and all would behave. Equilibrium settles in between, with just enough bad types working to make success informative to the degree that leaves them indifferent.

\pitfalls
\begin{itemize}
  \item Updating on effort rather than on outcomes: the market sees success or failure, not effort, so the posterior uses the success probabilities $p^1_H$, $p^1_L$ and the mixing probability $y$.
  \item Forgetting that without date-2 funding the continuation value is $A$ (keeping the cash), which is what kills the pooling equilibrium.
  \item Setting up the date-2 contract so that both types work: by (i), the only fundable contract lets the $B$-type shirk at date 2, which is why her continuation value is $p_LR_b+B$.
\end{itemize}

\bookans Tirole p.~597: (i) the only fundable contract makes the good type work and the bad type shirk, pledgeable income $(p_H-(1-\alpha)\Delta p)(R-b/\Delta p)$; (ii) no pooling (belief stays $\alpha$, no date-2 funding, so the bad type shirks), no separating (belief $\alpha_S$, funding, so the bad type strictly prefers to work); the semiseparating equilibrium solves $B=(\Delta p_1)[p_L(R-(I-A)/(p_H-(1-\alpha'_S)\Delta p))+B]$, which determines $\alpha'_S$ and hence the bad type's probability of working. Matches the book's answer (p.~597); we add the closed forms \eqref{r5e15-Rbs}--\eqref{r5e15-y}, and note that Bayes' rule must use date-1 probabilities for the assumptions to be consistent.
\refbib{Tirole \S3.2 ``The Fixed-Investment Model'', pp.~115--117; \S6.2 privately-known-prospects model, pp.~241--243; Exercise 6.1 (pooling with private benefits), p.~273; lecture Topic~3, ``Reputation and dynamic adverse selection''.}

\end{document}
